% Final_Report_PES.tex - Main LaTeX file
\documentclass[12pt,a4paper]{report}
\usepackage[utf8]{inputenc}
\usepackage[T1]{fontenc}
\usepackage{mathptmx} % Times New Roman font
\usepackage{geometry}
\geometry{
  a4paper,
  left=25mm,
  right=25mm,
  top=25mm,
  bottom=25mm,
}
\usepackage{microtype}
\usepackage{graphicx}
\usepackage{amsmath,amssymb,amsfonts}
\usepackage{amsthm}
\newtheorem{theorem}{Theorem}[chapter]
\newtheorem{proposition}[theorem]{Proposition}
\newtheorem{lemma}[theorem]{Lemma}
\newtheorem{corollary}[theorem]{Corollary}
\theoremstyle{definition}
\newtheorem{definition}[theorem]{Definition}
\newtheorem{remark}[theorem]{Remark}
\newtheorem{observation}[theorem]{Observation}
\newtheorem{heuristic}[theorem]{Heuristic}
\newtheorem{ruleofthumb}[theorem]{Rule of Thumb}
\newtheorem{toymodel}[theorem]{Simplified Model}
\usepackage{booktabs}
\usepackage{array}
\usepackage{tabularx}
\usepackage{pifont} % For checkmarks
\usepackage{multirow}
\usepackage{longtable}
\usepackage{enumitem}
\usepackage{multicol}
\usepackage{colortbl}
\usepackage{setspace}
\usepackage[numbers,sort&compress]{natbib}
\usepackage{titlesec}
\usepackage{fancyhdr}
\usepackage{float}
\usepackage{algorithm}
\usepackage{algpseudocode}
\usepackage{chngcntr}
\counterwithin{algorithm}{chapter}
\usepackage{caption}
\usepackage{subcaption}
\usepackage{tikz}
\usetikzlibrary{shapes.geometric, arrows.meta, positioning, fit, backgrounds}
\usepackage{pgfplots}
\pgfplotsset{compat=1.18}
\usepackage{tcolorbox}

% hyperref should be loaded last (after other packages)
\usepackage{hyperref}
\usepackage{bookmark}

% Custom commands
\newcommand{\source}[1]{\vspace{-0.5\baselineskip}\caption*{\small Source: #1}}

% Report metadata
\newcommand{\reportdate}{January 18, 2026}
\newcommand{\methodN}{80}

\hypersetup{
    colorlinks=true,
    linkcolor=black,
    citecolor=black,
    urlcolor=blue,
    pdftitle={Fine-Tuning Agentic AI Workflows Using Reinforcement Learning},
    pdfauthor={Group 6}
}

% Chapter formatting
\titleformat{\chapter}[display]
{\normalfont\bfseries\Huge}
{\chaptertitlename\ \thechapter}{20pt}{\Huge}

% Header/Footer
\pagestyle{fancy}
\fancyhf{}
\setlength{\headheight}{15pt}
\fancyhead[L]{\nouppercase{\leftmark}}
\fancyhead[R]{}
\fancyfoot[C]{\thepage}
\renewcommand{\headrulewidth}{0.4pt}
\renewcommand{\footrulewidth}{0pt}

% Make chapter-opening pages match the main style
\fancypagestyle{plain}{%
  \fancyhf{}%
  \setlength{\headheight}{15pt}%
  \fancyhead[L]{\nouppercase{\leftmark}}%
  \fancyhead[R]{}%
  \fancyfoot[C]{\thepage}%
  \renewcommand{\headrulewidth}{0.4pt}%
  \renewcommand{\footrulewidth}{0pt}%
}

\onehalfspacing
\raggedbottom

\begin{document}

\pagenumbering{gobble}

% Title Page
\hypersetup{pageanchor=false}
\begin{titlepage}
    \thispagestyle{empty}
    \centering
    \vspace*{0.8cm}

    {\Huge\bfseries Fine-Tuning Agentic AI Workflows Using Reinforcement Learning\par}
    \vspace{0.2cm}
    {\Large\bfseries A Survey and Experimental Framework\par}

    \vspace{0.8cm}
    {\Large\bfseries Final Report -- Literature Survey and Experimental Framework\par}
    \vspace{0.6cm}
    {\large\bfseries Group 6\par}

    \vspace{0.6cm}
    {\large\textbf{Members}\par}
    \vspace{0.2cm}
    \begin{tabular}{@{}ll@{}}
        Madhu Kumara L & PES2PGE24DS176 \\
        Sandhya Jeyaraj & PES2PGE24DS144 \\
        Mohammad Rafi & PES2PGE24DS157 \\
        Arumugam Chetty K & PES2PGE24DS137 \\
        Arvind C R & PES2PGE24DS140 \\
        Dhruva & PES2PGE23DS169 \\
    \end{tabular}

    \vspace{0.6cm}
    {\large\textbf{Mentors}\par}
    \vspace{0.2cm}
    Narayana Darapaneni (Director -- AIML)\\
    Anwesh Reddy Paduri (Senior Data Scientist)

    \vfill

    % TODO: Add PES University logo before final submission
    % Download from official PES resources and save as `pes_logo.png` in this directory.
    % Alternatively, use pes_logo.pdf or pes_logo.eps (LaTeX auto-detects format).
    \IfFileExists{pes_logo.png}{%
        \includegraphics[width=0.22\textwidth]{pes_logo.png}\par
    }{%
        \IfFileExists{pes_logo.pdf}{%
            \includegraphics[width=0.22\textwidth]{pes_logo.pdf}\par
        }{%
            % No logo file present: print nothing rather than a placeholder box
            \relax
        }%
    }

	    \vspace{0.4cm}
	    {\large Department of Artificial Intelligence and Machine Learning\par}
	    {\large PES University, Bengaluru, India\par}
	    \vspace{0.4cm}
	    {\large \reportdate\par}
\end{titlepage}
\hypersetup{pageanchor=true}
\pagenumbering{roman}

% Abstract
\chapter*{Abstract}
\addcontentsline{toc}{chapter}{Abstract}

Modern generative models acquire transformative capabilities when equipped with the ability to act autonomously---planning multi-step workflows, invoking external tools, analyzing feedback from execution, and iteratively refining their approach.
 This survey synthesizes \textbf{60+ papers} published between January 2023 and December 2025, systematically surveying the evolving reinforcement learning methodologies used to build such autonomous agents. We organize the discussion into \textbf{six topical domains} (core policy optimization, tool orchestration, memory architectures, multi-agent coordination, evaluation benchmarks, and safety); separately, papers are classified into \textbf{19 methodological categories} grouped into five clusters (Figure~\ref{fig:rl_taxonomy}).

\textbf{Key Findings.} We identify three primary branches of algorithms in current research. One approach builds on \textbf{Proximal Policy Optimization (PPO)-style alignment}, which integrates trained reward networks with value function baselines.
 This method is effective but requires significant memory. Another branch, characterized by \textbf{Direct Preference Optimization (DPO)}, skips explicit reward modeling
 by optimizing using preference rankings directly, making it more computationally efficient for offline data.
 A third, increasingly popular family includes \textbf{Group Relative Policy Optimization (GRPO)}, which uses group-level comparisons to optimize the policy without needing a critic network. GRPO-style methods are frequent among the 2024--2025 papers we reviewed, reflecting their role in memory-efficient training (we do not report a numeric share; see Section~\ref{sec:grpo}). We also examine the trade-offs between these algorithms across various deployment contexts.

\textbf{Analytical Insights.} We complement the empirical survey with theoretical discussion. We examine why PPO-style clipping improves optimization stability in practice and how success rates on agentic benchmarks may scale with model capacity. We also propose a simplified model clarifying when retrieval improves factuality versus when it fails, and provide a communication-cost analysis for common multi-agent coordination topologies.

\textbf{Empirical Results.} Reported results have improved substantially, although the commonly quoted pairs mix models, benchmark splits, and scaffolds: GPT-4 solved 1.7\% of the original SWE-bench~\cite{jimenez2024swebench}, while an agent scaffold around Claude~3.5 Sonnet is reported at 49\% on the 500-task SWE-bench Verified subset (a different split, so no ratio is computed); end-to-end RL raised Llama-3.1-8B on WebArena-Lite from 8.5\% to 44.8\% (WebAgent-R1); on OSWorld the best model in the original paper reached 12.2\% versus 72.4\% for humans~\cite{xie2024osworld}, and Mobile-Agent-v3 reports 37.7\% in 2025; and on AIME 2024, GPT-4o scores 9.3\% versus 79.8\% for DeepSeek-R1~\cite{deepseek2025r1}.\footnote{Figures are as reported by the cited papers or vendors as of December 2025.}

\textbf{Open Problems.} Several fundamental challenges persist: credit assignment in long-horizon trajectories, robust tool invocation with graceful failure handling, benchmark validity for real-world deployment prediction, and safety considerations for agents with environmental access. We propose a three-stage training pipeline (Supervised Fine-Tuning [SFT]$\rightarrow$RL$\rightarrow$Retrieval-Augmented Generation [RAG]) with testable hypotheses designed for academic computing resources; it was not executed here. Our executed experiments are limited to low-budget LoRA GRPO runs on GSM8K, where five Qwen3-8B seeds score 83.3\% on average on 200 held-out problems versus 82.0\% for the base model ($p=0.26$; no demonstrated gain).

\vspace{0.6cm}
\noindent \textbf{Keywords:} Autonomous Agents, RL-based Alignment, PPO, DPO, GRPO, RAG, Multi-Agent Coordination, Tool Orchestration, Code Generation.

\vspace{0.4cm}
\noindent \textbf{Highlights:}
\begin{itemize}
    \item Comprehensive synthesis of 60+ papers on RL-trained agents published 2023--2025
    \item Systematic comparison of PPO, DPO, and GRPO algorithm families with theoretical foundations
    \item Qualitative analysis of GRPO adoption and memory-efficiency implications
    \item Documentation of reported benchmark progress, separating like-for-like RL gains from cross-model comparisons
    \item Proposed (unexecuted) experimental framework, plus low-budget GRPO runs on GSM8K with released records
\end{itemize}

\begin{figure}[htbp]
\centering
\noindent \textbf{Graphical Abstract}\par\medskip
\begin{tikzpicture}[
    scale=0.82,  % Slightly reduced for more canvas width
    % Removed transform shape to prevent text scaling issues in PDF viewers
    topicdomain/.style={rectangle, draw, thick, rounded corners, minimum width=2.8cm, minimum height=1.4cm, align=center, font=\footnotesize\bfseries, inner sep=4pt},
    algo/.style={rectangle, draw, fill=blue!15, rounded corners, minimum width=2cm, minimum height=0.8cm, align=center, font=\scriptsize},
    arrow/.style={->, >=stealth, thick, gray!70},
    biarrow/.style={<->, >=stealth, thick, gray!50, dashed}
]
    % === STEP 1: Define all coordinates first (invisible) ===
    \coordinate (agent_pos) at (0,0);
    \coordinate (rl_pos) at (0,3.2);
    \coordinate (ppo_pos) at (-3.0,5.2);   % Spread wider
    \coordinate (dpo_pos) at (0,5.2);
    \coordinate (grpo_pos) at (3.0,5.2);   % Spread wider
    \coordinate (other_pos) at (5.8,6.2);  % Moved further right for clear separation
    \coordinate (tools_pos) at (-5.2,0.5);
    \coordinate (rag_pos) at (5.2,0.5);
    \coordinate (multi_pos) at (-4,-2.8);
    \coordinate (bench_pos) at (4,-2.8);
    \coordinate (safety_pos) at (0,-4.2);

    % === STEP 2: Draw ALL arrows FIRST (background layer) ===
    % Algorithm connections to RL box
    \draw[gray!60, thin] (ppo_pos) -- ([yshift=0.7cm]rl_pos);
    \draw[gray!60, thin] (dpo_pos) -- ([yshift=0.7cm]rl_pos);
    \draw[gray!60, thin] (grpo_pos) -- ([yshift=0.7cm]rl_pos);

    % Cross-domain connections (curved, drawn BEHIND nodes)
    \draw[biarrow] ([yshift=-0.9cm]tools_pos) to[bend right=25] ([yshift=-0.9cm]rag_pos);
    \draw[biarrow] ([xshift=-1.4cm]rl_pos) to[bend left=20] ([yshift=0.9cm]tools_pos);
    \draw[biarrow] ([xshift=1.4cm]rl_pos) to[bend right=20] ([yshift=0.9cm]rag_pos);

    % Main flow arrows
    \draw[arrow, line width=1.2pt] ([yshift=-0.7cm]rl_pos) -- ([yshift=0.8cm]agent_pos);
    \draw[arrow, line width=1.2pt] ([xshift=1.4cm]tools_pos) -- ([xshift=-2cm]agent_pos);
    \draw[arrow, line width=1.2pt] ([xshift=-1.4cm]rag_pos) -- ([xshift=2cm]agent_pos);
    \draw[arrow, line width=1.2pt] ([xshift=-1cm,yshift=-0.8cm]agent_pos) -- ([yshift=0.9cm]multi_pos);
    \draw[arrow, line width=1.2pt] ([xshift=1cm,yshift=-0.8cm]agent_pos) -- ([yshift=0.9cm]bench_pos);
    \draw[arrow, line width=1.2pt] ([yshift=-0.9cm]bench_pos) -- ([xshift=1.4cm]safety_pos);
    \draw[arrow, line width=1.2pt] ([yshift=-0.9cm]multi_pos) -- ([xshift=-1.4cm]safety_pos);

    % === STEP 3: Draw ALL nodes ON TOP of arrows ===
    % Central agent
    \node[topicdomain, fill=yellow!30, minimum width=4cm, minimum height=1.6cm] (agent) at (agent_pos) {Autonomous\\AI Agent};

    % Core RL domain (top)
    \node[topicdomain, fill=blue!20] (rl) at (rl_pos) {RL Algorithms\\(Ch. 2)};
    \node[algo] (ppo) at (ppo_pos) {PPO\\22\%};
    \node[algo] (dpo) at (dpo_pos) {DPO\\10\%};
    \node[algo, fill=green!25] (grpo) at (grpo_pos) {GRPO};
    \node[algo, fill=gray!15] (other) at (other_pos) {Other\\29\%};

    % Tool Learning (left)
    \node[topicdomain, fill=orange!20, minimum height=1.8cm] (tools) at (tools_pos) {Tool Learning\\(Ch. 3)\\{\tiny APIs, code, browsing}};

    % RAG domain (right)
    \node[topicdomain, fill=purple!20, minimum height=1.8cm] (rag) at (rag_pos) {RAG\\(Ch. 4)\\{\tiny Memory, retrieval}};

    % Multi-agent (bottom left) - Multi-Agent is covered in Ch. 4 with RAG
    \node[topicdomain, fill=red!20, minimum height=1.8cm] (multi) at (multi_pos) {Multi-Agent\\(Ch. 4)\\{\tiny Coordination}};

    % Benchmarks (bottom right) - Evaluation covered in Ch. 3
    \node[topicdomain, fill=cyan!20, minimum height=1.8cm] (bench) at (bench_pos) {Evaluation\\(Ch. 3)\\{\tiny SWE-bench, etc.}};

    % Safety (bottom) - Safety discussed in Ch. 6 (Limitations)
    \node[topicdomain, fill=gray!25] (safety) at (safety_pos) {Safety (Ch. 3 \& 6)};

    % === STEP 4: Info boxes (far from main diagram) ===
    \node[draw, thick, fill=green!10, rounded corners, font=\scriptsize, align=center, text width=4.2cm, inner sep=5pt] at (8,3.5) {
        \textbf{Key Insight:}\\
        GRPO eliminates critic network,\\
        cutting memory needs\\
        from 4$\times$ to $\sim$2$\times$ base model size
    };

    \node[draw, thick, fill=blue!10, rounded corners, font=\scriptsize, align=center, text width=3.8cm, inner sep=5pt] at (-8,3.5) {
        \textbf{Progress 2023--2025:}\\
        WebArena-Lite (8B, RL): 8.5\%$\rightarrow$44.8\%\\
        OSWorld: 12.2\% (2024) / 37.7\% (2025)\\
        AIME 2024: 9.3\%$\rightarrow$79.8\%
    };
\end{tikzpicture}
\caption*{\textbf{Survey Overview:} This survey maps 60+ papers across six interconnected research areas. Key finding: methods like GRPO have made practical training possible for autonomous agents that execute complex multi-step tasks without human intervention.
}
\end{figure}

\addtocontents{toc}{\protect\thispagestyle{empty}}
\tableofcontents

\cleardoublepage\phantomsection
\addcontentsline{toc}{chapter}{List of Figures}
\listoffigures

\cleardoublepage\phantomsection
\addcontentsline{toc}{chapter}{List of Tables}
\listoftables

\chapter*{Notation}
\addcontentsline{toc}{chapter}{Notation}
\label{ch:notation}

\noindent Table~\ref{tab:notation} summarizes the mathematical notation used throughout this survey. Standard reinforcement learning conventions follow \cite{sutton2018reinforcement}.

\begin{table}[h]
\centering
\begin{tabular}{@{}cl@{}}
\toprule
\textbf{Symbol} & \textbf{Description} \\
\midrule
$\mathcal{S}, \mathcal{A}, \Omega$ & State, action, and observation spaces \\
$\pi_\theta(a|s)$ & Policy parameterized by $\theta$; $s$ may denote belief-state/history embedding \\
$\pi_{\text{ref}}$ & Reference policy (frozen) \\
$r(s,a)$ & Reward function (lowercase $r$ preferred; some papers use $R$) \\
$H$ & Episode horizon (number of timesteps) \\
$\gamma$ & Discount factor $\in [0,1]$ \\
$\beta_{\text{KL}}$ & KL penalty coefficient (PPO/GRPO) \\
$\beta_{\text{DPO}}$ & DPO preference strength/temperature parameter \\
$\epsilon$ & PPO clipping threshold (typically 0.1--0.2) \\
$\hat{A}_t$ & Advantage estimate at time $t$ \\
$G$ & Group size for GRPO \\
$y_w, y_l$ & Preferred and non-preferred outputs (DPO) \\
$T(s'|s,a)$ & Transition dynamics \\
$h_t$ & Observation history up to time $t$ \\
$O(o \mid s, a)$ & Observation emission kernel of the POMDP \\
$\alpha$ & Learning rate (also denoted $\eta$ in some literature) \\
$\text{KL}[\cdot\|\cdot]$ & Kullback-Leibler divergence \\
$\sigma(z)$ & Logistic sigmoid: $\sigma(z) = 1/(1+e^{-z})$ \\
$s_G$ & Group standard deviation (within-batch) \\
\bottomrule
\end{tabular}
\caption{Mathematical notation used in this survey. See individual chapters for extended definitions. Note: $\beta$ is overloaded in RL literature; we distinguish $\beta_{\text{KL}}$ (KL penalty) from $\beta_{\text{DPO}}$ (preference strength) where context requires clarity.}
\label{tab:notation}
\end{table}

\clearpage
\pagenumbering{arabic}

\chapter{Introduction}
\label{ch:introduction}

\textit{This chapter motivates the need for reinforcement learning in training autonomous AI agents, introduces the mathematical framework (POMDPs), and outlines our survey's scope and contributions.}

\vspace{0.3cm}

Picture a developer staring at a stack trace. They scan the error, form a hypothesis, make a change, run the code again, and see what happens. Rinse and repeat until the bug is squashed. A year ago, getting an AI to do this meant holding its hand through every step. Today, systems built on ReAct-style reasoning \cite{yao2023react} can poke around websites, spin up test suites, and grind through multi-step problems without a human in the loop. That shift didn't happen by showing models more examples---it happened because researchers figured out how to train for outcomes, not imitation.

\subsection*{How to Read This Survey}

\textbf{Reading paths by audience:}
\begin{itemize}
    \item \textbf{Practitioners} seeking implementation guidance: Ch 2 (algorithms) $\rightarrow$ Ch 3 (tools) $\rightarrow$ Ch 5 (experiments)
    \item \textbf{Researchers} exploring open problems: Ch 1 (foundations) $\rightarrow$ Ch 4 (architectures) $\rightarrow$ Ch 6 (limitations)
    \item \textbf{Managers/Executives} wanting the big picture: Abstract $\rightarrow$ Key Takeaways boxes (end of each chapter) $\rightarrow$ Conclusion
\end{itemize}

\noindent\begin{minipage}{\linewidth}
\textbf{Chapter dependencies:}
\begin{center}
\begin{tikzpicture}[
    node distance=1.5cm,
    box/.style={rectangle, draw, fill=blue!10, text width=2.2cm, text centered, rounded corners, minimum height=0.7cm, font=\small},
    arrow/.style={->, >=stealth, thick}
]
    \node[box] (ch1) {Ch 1: Intro};
    \node[box, right=1cm of ch1] (ch2) {Ch 2: Algorithms};
    \node[box, right=1cm of ch2] (ch3) {Ch 3: Tools};
    \node[box, below=0.8cm of ch2] (ch4) {Ch 4: RAG/Multi-Agent};
    \node[box, right=1cm of ch4] (ch5) {Ch 5: Experiments};
    \node[box, right=1cm of ch5] (ch6) {Ch 6: Limitations};

    \draw[arrow] (ch1) -- (ch2);
    \draw[arrow] (ch2) -- (ch3);
    \draw[arrow] (ch2) -- (ch4);
    \draw[arrow] (ch3) -- (ch5);
    \draw[arrow] (ch4) -- (ch5);
    \draw[arrow] (ch5) -- (ch6);
\end{tikzpicture}
\end{center}
\end{minipage}

\textbf{Notation:} All mathematical symbols are defined in Table~\ref{tab:notation} (p.~\pageref{tab:notation}). Key terms are \textbf{bolded} and expanded on first use.

\section{Why Supervised Learning Isn't Enough}
Show a model a thousand correct solutions, and it learns to generate text that mimics those solutions. This represents pattern matching, not problem-solving. When unexpected errors occur, surface-level similarity to training data fails; the model needs to figure out what actually works through trial and error.

Chess makes the distinction clear. You can memorize grandmaster games and still lose to anyone who can calculate three moves ahead. Opening theory is useful; it's not the whole game. Same with agents: watching demos teaches the format, but navigating back from a failed state requires a different kind of knowledge---the kind you only get from trial and error.

Here's another angle: sequential reasoning trades time for depth. A single forward pass gives you one shot at the answer. Chain-of-thought prompting gives you multiple passes where each step can build on the last, catch mistakes, and adjust. That's not just ``more compute''---it's structured compute where the model's own outputs become working memory. GRPO (Group Relative Policy Optimization---a memory-efficient RL algorithm we detail in Section~\ref{sec:grpo}) takes a similar idea to training: instead of learning from one attempt per prompt, you learn from the relative quality of multiple attempts. More signal, same data budget.

\begin{figure}[htbp]
\centering
\begin{tikzpicture}[
    node distance=1.2cm,
    thought/.style={rectangle, draw, fill=yellow!20, text width=3.5cm, text centered, rounded corners, minimum height=0.8cm},
    action/.style={rectangle, draw, fill=orange!20, text width=3.5cm, text centered, rounded corners, minimum height=0.8cm},
    arrow/.style={->, >=stealth, thick}
]
    % Chain of Thought
    \node[thought] (cot1) {\textbf{CoT:} Thought 1};
    \node[thought, below=0.8cm of cot1] (cot2) {Thought 2};
    \node[thought, below=0.8cm of cot2] (cot3) {Thought 3};
    \node[action, below=0.8cm of cot3] (cot_ans) {Final Answer};

    \draw[arrow] (cot1) -- (cot2);
    \draw[arrow] (cot2) -- (cot3);
    \draw[arrow] (cot3) -- (cot_ans);

    % ReAct
    \node[thought, right=4cm of cot1] (react1) {\textbf{ReAct:} Thought};
    \node[action, below=0.6cm of react1] (action1) {Action $\to$ Obs};
    \node[thought, below=0.6cm of action1] (react2) {Thought};
    \node[action, below=0.6cm of react2] (action2) {Action $\to$ Obs};
    \node[action, below=0.6cm of action2] (react_ans) {Final Answer};

    \draw[arrow] (react1) -- (action1);
    \draw[arrow] (action1) -- (react2);
    \draw[arrow] (react2) -- (action2);
    \draw[arrow] (action2) -- (react_ans);
\end{tikzpicture}
\caption{Comparing Chain-of-Thought (CoT) and ReAct reasoning. \textbf{Left:} CoT lays out a chain of reasoning, then answers. \textbf{Right:} ReAct interleaves reasoning with actions, letting models pull in information and adapt on the fly.}
\label{fig:reasoning_comparison}
\source{Adapted from Wei et al.~\cite{wei2022chain} and Yao et al.~\cite{yao2023react}}
\end{figure}

\section{What We're Actually Trying to Solve}
The gap we care about isn't between ``describing'' and ``explaining''---it's between \textit{talking about doing something} and \textit{actually doing it}. LLMs read mountains of text where people describe how to fix bugs, configure servers, or navigate codebases. But that text almost never shows the messy reality: the wrong turns, the retries, the moment when the stack trace doesn't match anything in Stack Overflow.

Describing how to fix a bug and actually fixing a bug are different skills. One involves recalling patterns from training data. The other involves navigating files, running commands, reading output, and adjusting when things don't work. Standard supervised learning can't teach the second skill because the training data doesn't include the feedback loop---the ``I tried X, it failed, so I tried Y'' part.

Our angle of attack: training algorithms---PPO (Proximal Policy Optimization), DPO (Direct Preference Optimization), and GRPO---that optimize for whether the task got done, not whether the output looks like the training examples. Plus architectures (ReAct, tool integration) that give models access to the full execution cycle---act, observe, adjust.

\subsection{The Math: Agents as POMDPs}
We model agentic AI as a Partially Observable MDP (POMDP) because it captures the key constraint: agents never see everything. They get text snippets, screenshots, error messages---not the full state of the filesystem or the browser's internal DOM. Using a standard MDP would pretend observability exists when it doesn't. The POMDP tuple $\mathcal{M} = \langle \mathcal{S}, \mathcal{A}, \Omega, T, O, R, \gamma \rangle$ breaks down as:
\begin{itemize}
    \item $\mathcal{S}$: The true state of the world (file contents, browser state, server responses). Hidden from the agent.
    \item $\mathcal{A}$: What the agent can do---token sequences that might be tool calls, code snippets, or plain text. Unlike Atari where actions are discrete button presses, here the action space is combinatorially huge.
    \item $\Omega$: What the agent actually sees (text, images, whatever the interface shows it).
    \item $T(s'|s,a)$: How the environment responds to actions. Modeling this explicitly reminds us that environments are dynamic and unpredictable.
    \item $O(o|s)$: The lossy mapping from true state to observations. Critical---observations hide information.
    \item $r(s,a)$: The reward signal. Could be task success, human ratings, or intermediate process scores.
    \item $\gamma \in [0,1]$: How much future rewards are discounted. Tuning this trades off myopia against long-horizon planning.
\end{itemize}

\begin{tcolorbox}[colback=gray!5!white, colframe=gray!60!black, title=\textbf{Concrete Example: Bug-Fixing Agent as POMDP}]
\small
Consider an agent tasked with fixing a failing test in a Python project:
\begin{itemize}
    \item $\mathcal{S}$: Full codebase state (all files, git history, test results)---agent cannot see everything
    \item $\Omega$: Error message shown: \texttt{AssertionError: expected 42, got None}
    \item $\mathcal{A}$: Agent chooses action \texttt{read\_file("utils.py")} to investigate
    \item $T$: Environment returns file contents (state transition)
    \item $O$: Agent sees only the 50 lines around the relevant function (partial observation)
    \item $R$: +1 if test passes after edit, 0 otherwise
\end{itemize}
The agent must infer the hidden bug location from partial observations and choose actions that maximize expected reward.
\end{tcolorbox}

The agent policy $\pi_\theta(a|h_t)$ maps observation history $h_t = (o_1, a_1, \ldots, a_{t-1}, o_t)$ to action distributions. The objective is:
\begin{equation}
\max_\theta \; \mathbb{E}_{\tau \sim \pi_\theta} \left[ \sum_{t=0}^{H-1} \gamma^t \Big( r(s_t, a_t) - \beta_{\text{KL}}\, \text{KL}\big[\pi_\theta(\cdot \mid h_t) \,\|\, \pi_{\text{ref}}(\cdot \mid h_t)\big] \Big) \right]
\label{eq:objective}
\end{equation}
where the expectation is over trajectories $\tau$ generated by $\pi_\theta$, the per-step KL is evaluated at the realized history $h_t$, and the KL term keeps the learned policy $\pi_\theta$ tethered to a reference policy $\pi_{\text{ref}}$ (typically the SFT model), discouraging the exploitation of reward shortcuts.

\section{Survey Scope and Taxonomy}
We organize the scattered research on autonomous systems into six connected areas. \textit{Note: To capture the rapidly evolving frontier of agentic technologies, this survey includes analysis of emerging research and preprints through early 2026.}

\begin{enumerate}
    \item \textbf{Tool Invocation:} How agents turn ideas into actions---API calls, code generation, function composition (e.g., Toolformer's self-supervised learning).
    \item \textbf{Policy Optimization:} The math behind agent improvement---from PPO's stability tricks to DPO's simplified approach and GRPO's group-based learning.
    \item \textbf{Knowledge Augmentation:} Ways to extend agent memory beyond context limits using retrieval, external databases, and memory systems like HippoRAG.
    \item \textbf{Alignment Methods:} How to keep agents helpful and safe---through preference learning, self-critique, and frameworks for scalable oversight.
    \item \textbf{Collaborative Architectures:} How multiple agents work together---role-based teams, emergent coordination, and decentralized decision-making.
    \item \textbf{Evaluation Methods:} Real-world testbeds that measure actual competence---code repair (SWE-bench), web tasks (WebArena), and desktop automation (OSWorld).
\end{enumerate}

\subsection{Systematic Review Methodology}
We took a structured approach to cast a wide net. We queried arXiv, Semantic Scholar, ACL Anthology, and Google Scholar (January 2023--December 2025) using terms: ``agentic AI,'' ``LLM agents,'' ``RLHF,'' ``DPO,'' ``GRPO,'' ``tool use LLM,'' and ``multi-agent coordination.'' Papers were included if they presented novel agent training methods, introduced benchmarks, or addressed alignment/safety. Two reviewers independently categorized papers across our 19-direction taxonomy; disagreements were resolved through discussion. The final corpus comprises \textbf{60+ papers}.

\section{Related Surveys and Our Positioning}
Table~\ref{tab:survey_comparison} positions our work relative to existing surveys in the agentic AI and LLM alignment literature.

\begin{table}[h]
\centering
\caption{Comparison with Existing Surveys on Agentic AI and LLM Alignment}
\label{tab:survey_comparison}
\small
\begin{tabular}{@{}p{2.5cm}ccccccc@{}}
\toprule
\textbf{Survey} & \textbf{Year} & \textbf{RL Algos} & \textbf{Tools} & \textbf{RAG} & \textbf{Multi-Agent} & \textbf{Benchmarks} & \textbf{Taxonomy} \\
\midrule
\citet{xi2023rise} & 2023 & \ding{55} & \ding{51} & \ding{55} & \ding{51} & \ding{51} & Partial \\
\citet{sumers2023cognitive} & 2023 & \ding{55} & \ding{51} & \ding{55} & \ding{55} & \ding{55} & \ding{51} \\
\citet{wang2024planning} & 2024 & \ding{55} & \ding{51} & \ding{55} & \ding{55} & \ding{51} & \ding{51} \\
\citet{wang2025landscape} & 2025 & \ding{51} & \ding{55} & \ding{55} & \ding{55} & \ding{51} & \ding{51} \\
\citet{gao2023retrieval} & 2023 & \ding{55} & \ding{55} & \ding{51} & \ding{55} & \ding{51} & \ding{51} \\
\midrule
\textbf{Ours} & 2026 & \ding{51} & \ding{51} & \ding{51} & \ding{51} & \ding{51} & \ding{51} \\
\bottomrule
\end{tabular}
\vspace{0.2cm}

\scriptsize
\ding{51} = Covered in depth, \ding{55} = Not covered or mentioned briefly. Our survey uniquely integrates RL algorithms (PPO/DPO/GRPO), tool use, RAG, multi-agent systems, and benchmarks into a unified framework with a 19-category taxonomy.
\end{table}

\section{Our Contributions}
This report makes the following contributions to autonomous AI agent research:

\begin{enumerate}
    \item \textbf{Comprehensive Systematic Survey (C1):} We comb through \textbf{60+ papers} (2023--2025), analyzing the literature to identify key patterns. The field has fragmented across dozens of algorithmic variants without clear principles for selection. We classify papers into \textbf{19 research categories} (grouped into five methodological clusters, Figure~\ref{fig:rl_taxonomy}) and discuss them across six topical domains because this structure makes explicit the fundamental trade-offs: core algorithms differ by optimization target, tool learning differs by supervision type, multi-agent systems differ by coordination mechanism. This isn't just organization—it's pattern extraction from noise.

    \item \textbf{Algorithmic Landscape Analysis (C2):} We describe algorithm adoption in our corpus qualitatively (no retained variant catalog supports a numeric share). We identify \textbf{15 algorithm variants} (Flow-GRPO, M-GRPO, AT-GRPO, RAPO, ToolPO, etc.) and lay out when to use which algorithm.

    \item \textbf{Formal Analysis (C3):} We work through \textbf{four results with appropriately scoped labels} that adapt existing frameworks to the agentic setting: trust-region motivation for PPO vs.\ TRPO (Proposition~\ref{prop:convergence}), agentic scaling considerations (Heuristic~\ref{lem:scaling}), a simplified model for when retrieval helps factuality (Simplified Model~\ref{prop:hallucination}), and communication overhead by topology (Heuristic~\ref{thm:multiagent}). We use ``heuristic'' and ``simplified model'' labels where formal guarantees do not exist.

    \item \textbf{Benchmark Progress Documentation (C4):} We document reported progress on SWE-bench, WebArena-Lite, OSWorld, AIME 2024, and GAIA (Table~\ref{tab:sota_benchmarks}), distinguishing like-for-like RL gains (e.g., WebArena-Lite 8.5\%$\rightarrow$44.8\% for Llama-3.1-8B) from cross-model, cross-split comparisons that do not support a single improvement multiplier.

    \item \textbf{Training Instability Taxonomy (C5):} We catalog six failure modes because without understanding failure, optimization is blind. The Echo Trap happens because agents overfit to patterns that work in training but fail in deployment—this reveals a fundamental issue with reward optimization without distributional robustness. Entropy collapse occurs because exploration decays too fast, an inevitable consequence of exponential policy updates without explicit entropy maintenance. Reward instability reflects the fundamental problem of optimizing against a learned proxy that's itself changing. Gradient spikes emerge from long-horizon credit assignment where small estimation errors compound. Tool hallucination shows that models don't truly understand tool semantics, just statistical correlations. Credit assignment failure is the classic RL problem: sparse rewards provide insufficient signal. Each failure mode reveals a structural problem, not just a bug—mapping these failures lets practitioners apply focused fixes instead of blindly tweaking hyperparameters.

    \item \textbf{Proposed Experimental Framework (C6):} We design (but do not execute) a three-stage training pipeline (SFT$\rightarrow$RL$\rightarrow$RAG) with \textbf{five testable hypotheses} and ablation grids. The executed experiments are limited to low-budget LoRA GRPO runs on GSM8K (Chapter~\ref{ch:experiments}).

    \item \textbf{Research Roadmap (C7):} We identify seven critical gaps because the field cannot progress without acknowledging fundamental barriers. The benchmark-reality gap exists because benchmarks are sanitized while reality is messy—this won't fix itself without new evaluation paradigms. Long-horizon planning fails because credit assignment breaks beyond 20 steps—this requires algorithmic innovation, not just scale. Safety-capability tradeoffs emerge because capability improvements expand action spaces exponentially—the blowup is combinatorial and cannot be dodged without rethinking how agents are structured. Computational barriers reflect fundamental memory-time tradeoffs—algorithms must become more efficient, not just run on bigger hardware. These gaps are not implementation issues; they are the next frontier.
\end{enumerate}

\noindent Table~\ref{tab:contribution_summary} shows where each contribution fills holes left by prior surveys.

\begin{table}[h]
\centering
\caption{Comparison with Existing Surveys on LLM Agents}
\label{tab:contribution_summary}
\small
\begin{tabular}{@{}p{3.5cm}ccccc@{}}
\toprule
\textbf{Feature} & \textbf{Xi+23} & \textbf{Wang+24} & \textbf{Sumers+23} & \textbf{Gao+23} & \textbf{Ours} \\
\midrule
RL Algorithm Depth & \ding{55} & Partial & \ding{55} & \ding{55} & \ding{51} \\
GRPO Coverage & \ding{55} & \ding{55} & \ding{55} & \ding{55} & \ding{51} \\
Quantitative Analysis & \ding{55} & Partial & \ding{55} & \ding{55} & \ding{51} \\
Training Instabilities & \ding{55} & \ding{55} & \ding{55} & \ding{55} & \ding{51} \\
2024--2025 Coverage & \ding{55} & Partial & \ding{55} & \ding{55} & \ding{51} \\
Experimental Framework & \ding{55} & \ding{55} & \ding{55} & \ding{55} & \ding{51} \\
Theoretical Results & \ding{55} & \ding{55} & Partial & \ding{55} & \ding{51} \\
\bottomrule
\end{tabular}
\end{table}

\begin{figure}[htbp]
\centering
\begin{tikzpicture}[scale=0.85, transform shape]
    % Timeline axis
    \draw[->, thick, gray] (0,0) -- (15,0);
    \node[below, font=\small] at (15,-0.2) {Time};

    % Year markers
    \foreach \x/\year in {1/2022, 4.5/2023, 8/2024, 11.5/2025} {
        \draw[thick] (\x,-0.15) -- (\x,0.15);
        \node[below, font=\bfseries] at (\x,-0.3) {\year};
    }

    % Era backgrounds (extended height for labels)
    \fill[blue!10, rounded corners] (0.3,0.5) rectangle (3.5,5.2);
    \fill[green!10, rounded corners] (3.7,0.5) rectangle (7.3,5.2);
    \fill[orange!10, rounded corners] (7.5,0.5) rectangle (11,5.2);
    \fill[red!10, rounded corners] (11.2,0.5) rectangle (14.5,5.2);

    % Era labels (at top, with white background for clarity)
    \node[font=\bfseries\small, fill=blue!20, rounded corners, inner sep=3pt] at (1.9,4.8) {Foundation Era};
    \node[font=\bfseries\small, fill=green!20, rounded corners, inner sep=3pt] at (5.5,4.8) {Emergence Era};
    \node[font=\bfseries\small, fill=orange!20, rounded corners, inner sep=3pt] at (9.25,4.8) {Scaling Era};
    \node[font=\bfseries\small, fill=red!20, rounded corners, inner sep=3pt] at (12.85,4.8) {Autonomy Era};

    % 2022 milestones (adjusted positions)
    \node[draw, fill=white, rounded corners, text width=2.5cm, align=center, font=\scriptsize, inner sep=4pt] at (1.9,3.6) {InstructGPT\\(PPO + RLHF)};
    \node[draw, fill=white, rounded corners, text width=2.5cm, align=center, font=\scriptsize, inner sep=4pt] at (1.9,2.3) {Constitutional AI\\(Self-supervision)};
    \node[draw, fill=white, rounded corners, text width=2.5cm, align=center, font=\scriptsize, inner sep=4pt] at (1.9,1.0) {Chain-of-Thought\\(Reasoning)};

    % 2023 milestones
    \node[draw, fill=white, rounded corners, text width=2.6cm, align=center, font=\scriptsize, inner sep=4pt] at (5.5,3.6) {DPO\\(No learned RM)};
    \node[draw, fill=white, rounded corners, text width=2.6cm, align=center, font=\scriptsize, inner sep=4pt] at (5.5,2.3) {ReAct Framework\\(Reasoning + Action)};
    \node[draw, fill=white, rounded corners, text width=2.6cm, align=center, font=\scriptsize, inner sep=4pt] at (5.5,1.0) {Toolformer, GPT-4\\(Tool use emerges)};

    % 2024 milestones
    \node[draw, fill=white, rounded corners, text width=2.6cm, align=center, font=\scriptsize, inner sep=4pt] at (9.25,3.6) {OpenHands, SWE-bench\\(Code agents)};
    \node[draw, fill=white, rounded corners, text width=2.6cm, align=center, font=\scriptsize, inner sep=4pt] at (9.25,2.3) {WebArena, OSWorld\\(GUI agents)};
    \node[draw, fill=white, rounded corners, text width=2.6cm, align=center, font=\scriptsize, inner sep=4pt] at (9.25,1.0) {Multi-agent systems\\(Coordination)};

    % 2025 milestones
    \node[draw, fill=white, rounded corners, text width=2.6cm, align=center, font=\scriptsize, inner sep=4pt] at (12.85,3.6) {DeepSeek-R1\\(GRPO dominance)};
    \node[draw, fill=white, rounded corners, text width=2.6cm, align=center, font=\scriptsize, inner sep=4pt] at (12.85,2.3) {79.8\% AIME 2024\\(Expert-competitive)};
    \node[draw, fill=white, rounded corners, text width=2.6cm, align=center, font=\scriptsize, inner sep=4pt] at (12.85,1.0) {37.7\% OSWorld\\(Mobile-Agent-v3)};

    % Benchmark progress arrow
    \draw[->, thick, purple!70] (2,-1.2) -- (13,-1.2);
    \node[below, font=\scriptsize, purple!70] at (7.5,-1.4) {AIME 2024: 9.3\% (GPT-4o) $\rightarrow$ 79.8\% (DeepSeek-R1)};
\end{tikzpicture}
\caption{Evolution of autonomous agents from 2022 to 2025. Research unfolded across four eras: \textbf{Foundation} (RLHF methods established), \textbf{Emergence} (agents acquire tool-calling and multi-step reasoning), \textbf{Scaling} (benchmarks proliferate, code/GUI agents mature), and \textbf{Autonomy} (GRPO unlocks lean self-improvement loops, with some tasks now rivaling human scores).}
\label{fig:timeline}
\end{figure}

\section{Report Organization}

The report breaks down as follows:
\begin{itemize}
    \item \textbf{Chapter~\ref{ch:foundations}} covers the optimization methods (PPO, DPO, GRPO) and the RLHF basics that power agent training.
    \item \textbf{Chapter~\ref{ch:tooluse}} examines tool interaction, benchmark evaluation, and safety considerations for grounded agent behavior.
    \item \textbf{Chapter~\ref{ch:advancedarch}} discusses memory systems (RAG) and multi-agent coordination architectures.
    \item \textbf{Chapter~\ref{ch:experiments}} proposes our experimental framework with testable hypotheses and ablation designs.
    \item \textbf{Chapter~\ref{ch:limitations}} lays out where current approaches fall short and pinpoints research gaps.
    \item \textbf{Chapter~\ref{ch:conclusion}} pulls together takeaways and offers advice for practitioners.
\end{itemize}

\begin{tcolorbox}[colback=cyan!5!white, colframe=cyan!60!black, title=\textbf{Chapter 1 Key Takeaways}]
\begin{enumerate}
    \item \textbf{RL $>$ imitation:} Supervised learning teaches format; RL teaches problem-solving through trial and error.
    \item \textbf{POMDP framing is essential:} Agents operate under partial observability---they never see the full state.
    \item \textbf{Fast reported progress:} e.g., AIME 2024 9.3\%$\rightarrow$79.8\% (GPT-4o vs.\ DeepSeek-R1); quoted SWE-bench pairs mix splits and scaffolds and are not attributable to RL alone.
    \item \textbf{Three algorithm families:} PPO (exploration), DPO (preferences), GRPO (memory-efficient reasoning).
\end{enumerate}
\end{tcolorbox}

\chapter{Foundations and Alignment Algorithms}
\label{ch:foundations}

\textit{This chapter covers the three dominant algorithm families for training agentic AI: PPO-style RLHF (\S\ref{sec:rlhf_mech}, \S\ref{sec:ppo}), DPO-style direct optimization (\S\ref{sec:dpo}), and GRPO-style group comparison (\S\ref{sec:grpo}). We present a quantitative analysis of algorithm adoption and provide selection guidelines.}

\vspace{0.3cm}

\noindent\textit{Mathematical notation used in this chapter follows Table~\ref{tab:notation}.}

\section{Mechanisms of RLHF}
\label{sec:rlhf_mech}
The primary challenge is that gradient descent needs something to optimize, but ``make this response more helpful'' isn't a loss function you can differentiate. You can't write down an equation for what makes one answer better than another---ask ten people and you'll get ten different rankings. \citet{christiano2017deep} figured out a workaround: don't ask humans to score responses on some scale (they're terrible at it), just ask them to pick which response is better. Collect enough of these A-vs-B comparisons and you can train a neural network to approximate the pattern in those choices. It's a noisy proxy for what people actually want, but it's enough to give gradient descent something concrete to chase.

\begin{figure}[htbp]
\centering
\begin{tikzpicture}[
    stage/.style={rectangle, draw, fill=green!20, text width=2.2cm, text centered, minimum height=1.4cm, rounded corners},
    arrow/.style={->, >=stealth, very thick}
]
    \node[stage] (base) {Base Model\\{\scriptsize(Llama, Qwen)}};
    \node[stage, right=2cm of base] (sft) {SFT\\{\scriptsize(Supervised\\Fine-Tuning)}};
    \node[stage, right=2cm of sft] (rl) {RL Tuning\\{\scriptsize(PPO/DPO/\\GRPO)}};
    \node[stage, right=2cm of rl] (agent) {Fine-tuned\\Agent};

    \draw[arrow] (base) -- node[above, font=\scriptsize] {Demos} (sft);
    \draw[arrow] (sft) -- node[above, font=\scriptsize] {Rewards} (rl);
    \draw[arrow] (rl) -- node[above, font=\scriptsize] {Deploy} (agent);

    % Annotations
    \node[below=0.3cm of sft, font=\tiny] {5K--10K demos};
    \node[below=0.3cm of rl, font=\tiny] {Environment interaction};
\end{tikzpicture}
\caption{Pretraining and two-stage post-training (Base $\rightarrow$ SFT $\rightarrow$ RL); distinct from the proposed SFT$\rightarrow$RL$\rightarrow$RAG pipeline of Chapter~\ref{ch:experiments}. \textbf{Stage 1:} A base model supplies language understanding. \textbf{Stage 2:} Imitation learning on 5K--10K examples drills task formats. \textbf{Stage 3:} RL tunes the policy through environment interaction, letting agents cope with novel situations outside their training distribution.}
\label{fig:rl_pipeline}
\source{Based on DeepSeek-R1 methodology~\cite{deepseek2025r1}}
\end{figure}

\begin{figure}[htbp]
\centering
\begin{tikzpicture}[scale=0.9, transform shape]
    % === AGENT BOX ===
    \node[draw, thick, fill=blue!20, rounded corners, minimum width=4cm, minimum height=2.5cm] (agent) at (0, 2.5) {};
    \node[font=\bfseries] at (0, 3.3) {Agent};

    % Agent internals
    \node[draw, fill=blue!10, rounded corners, font=\scriptsize, minimum width=1.5cm] (policy) at (-0.8, 2.3) {Policy $\pi_\theta$};
    \node[draw, fill=blue!10, rounded corners, font=\scriptsize, minimum width=1.5cm] (value) at (0.8, 2.3) {Value $V_\phi$};

    % === ENVIRONMENT BOX ===
    \node[draw, thick, fill=green!20, rounded corners, minimum width=4cm, minimum height=2cm] (env) at (0, -1.5) {};
    \node[font=\bfseries] at (0, -0.8) {Environment};
    \node[font=\scriptsize, gray] at (0, -1.8) {(Web, Code, OS, APIs)};

    % === ARROWS WITH LABELS ===
    % Action arrow (Agent → Environment)
    \draw[->, thick, red!70!black, line width=1.5pt] (1.5, 1.1) -- (1.5, -0.3);
    \node[right, font=\small, red!70!black] at (1.6, 0.4) {Action $a_t$};

    % State arrow (Environment → Agent)
    \draw[->, thick, blue!70!black, line width=1.5pt] (-1.5, -0.3) -- (-1.5, 1.1);
    \node[left, font=\small, blue!70!black] at (-1.6, 0.4) {State $s_{t+1}$};

    % Reward arrow (curved, from environment to agent)
    \draw[->, thick, orange!80!black, line width=1.5pt] (2.3, -1.5) to[out=0, in=0] (2.3, 2.5);
    \node[right, font=\small, orange!80!black] at (2.4, 0.5) {Reward $r_t$};

    % === TIME STEP INDICATOR ===
    \node[draw, dashed, gray, rounded corners, font=\scriptsize, inner sep=4pt] at (0, -3.2) {
        $t \rightarrow t+1$: Agent observes $s_t$, takes $a_t$, receives $r_t$, transitions to $s_{t+1}$
    };

    % === RL OBJECTIVE ===
    \node[draw, thick, fill=yellow!20, rounded corners, font=\small, align=center, inner sep=6pt] at (6, 0.5) {
        \textbf{Objective:}\\[2pt]
        $\max_\theta \mathbb{E}\left[\sum_{t=0}^{H} \gamma^t r_t\right]$
    };

\end{tikzpicture}
\caption{The reinforcement learning loop. At each timestep, the \textbf{agent} observes state $s_t$, selects action $a_t$ according to policy $\pi_\theta$, and receives reward $r_t$ from the \textbf{environment}. We want a policy that racks up as much long-run reward as possible. In the agentic setting, the ``environment'' spans web browsers, code interpreters, operating systems, and API back-ends.}
\label{fig:rl_loop}
\source{Standard RL formulation~\cite{sutton2018reinforcement}}
\end{figure}

The recipe that emerged from summarization work \cite{stiennon2020learning} goes like this:
\begin{enumerate}
    \item \textbf{Teach the basics:} Fine-tune your base model on demonstration data. This gets it speaking the right language and following the task format---think of it as making sure the student knows how to hold a pencil before you grade their essays.
    \item \textbf{Build a scorer:} Train a separate network $r_\phi(x, y)$ that predicts which response humans would prefer. Feed it thousands of A-vs-B comparisons; it learns to spot the patterns in what annotators picked.
    \item \textbf{Optimize against the scorer:} Run PPO to push the model toward higher scores while keeping it from wandering too far from its initialization. That KL penalty matters---without it, models find bizarre loopholes that score well but produce garbage.
\end{enumerate}

\section{Limitations of RLHF and Mitigation Strategies}

The vanilla RLHF setup has some ugly failure modes that won't go away with better hyperparameters:

\textit{Humans are slow.} Every preference judgment needs a person's eyeballs, which means your annotation budget puts a hard ceiling on training scale. Want 10$\times$ more data? Hire 10$\times$ more annotators.

\textit{Reward models can be gamed.} Push the optimization hard enough and your model will discover responses that score high but read like nonsense---the reward network learned to recognize surface patterns, not actual quality.

\textit{Errors compound.} You're stacking three imperfect steps: SFT introduces biases, reward modeling adds noise, and RL amplifies whatever mistakes survived. Small issues become big ones.

Researchers have tackled each of these:

\textbf{Let the model supervise itself:} \citet{bai2022constitutional} taught models a set of principles, then had them critique and revise their own outputs in a loop. Harmful responses dropped substantially across their test suite, and the human annotators barely had to lift a finger.

\textbf{Let AI do the ranking:} \citet{lee2023rlaif} showed you can swap human annotators for a capable LLM and still get good alignment---at least for summarization. Annotation costs fell by more than 10$\times$ with no obvious quality hit.

\textbf{Borrow from behavioral economics:} \citet{ethayarajh2024kto} noticed that people feel losses more sharply than equivalent gains (classic prospect theory stuff). Their KTO loss function bakes this asymmetry in directly, and they got comparable results using 2--4$\times$ less training data.

\section{Taxonomy of Agentic RL}
\label{sec:taxonomy}

Building on recent surveys \cite{wang2025landscape}, whose capability-by-environment framing we adopt rather than claim as new, we organize agentic RL research along two axes: the cognitive capability being optimized and the environment where it's applied. This grid exposes recurring themes as well as domain-specific headaches.

\subsection{Core Agentic Capabilities}
Each capability represents a different optimization target:
\begin{itemize}
    \item \textbf{Hierarchical Goal Decomposition:} Training agents to break complex problems into manageable sub-problems, whether through lookahead planning or recursive planning with reflection.
    \item \textbf{Tool Mastery:} Developing competence in using APIs---selecting the right ones, constructing valid calls, and interpreting results---through trajectory-level feedback.
    \item \textbf{Memory Management:} Breaking free from rigid context limits through learned memory operations---deciding what to stash, when to fetch, and how to compress history without losing the important bits.
    \item \textbf{Deliberate Reasoning:} Reinforcing explicit multi-step thinking (chains, trees, graphs) with process rewards that value valid intermediate deduction steps, not just final answers.
    \item \textbf{Self-Improvement:} Enabling agents to learn new skills without external curriculum---through self-play, curiosity signals, or synthetic task generation.
    \item \textbf{Active Perception:} For vision-equipped agents, learning where to look, when to capture screenshots, and how to balance observation cost against information gain.
\end{itemize}

\subsection{Task Domains}
Agents specialize for different environments:
\begin{itemize}
    \item \textbf{Software Engineering:} Agents working in codebases---navigating repositories, diagnosing bugs from stack traces, writing patches, and validating fixes through tests. SWE-bench exemplifies these challenges.
    \item \textbf{Mathematics:} From word problems and competition math to formal proofs in theorem provers like Lean or Coq.
    \item \textbf{Digital Interfaces:} Controlling graphical interfaces---web browsers, mobile apps, desktop OS---where observation spaces explode and actions require pixel-perfect precision.
    \item \textbf{Physical Interaction:} Connecting language understanding with motor control for robotics and spatial navigation, where simulation-to-reality transfer remains challenging.
    \item \textbf{Agent Networks:} Optimizing how multiple language-model agents coordinate---whether cooperating on shared goals, competing for resources, or negotiating with incomplete information.
\end{itemize}

\subsection{Comprehensive RL Taxonomy}
We organize papers into 19 categories not by convenience but by what they optimize. The five-cluster organization (Figure~\ref{fig:rl_taxonomy} and Appendix~\ref{app:paper_mapping}) reflects fundamental divisions in how agents learn. Core algorithms differ by their optimization target: policy-based methods optimize action selection directly, value-based methods optimize state evaluations, model-based methods learn environment dynamics. Learning paradigms differ by their data source: RLHF uses human preferences, imitation uses demonstrations, offline RL uses logged data. These categories aren't just labels—they reflect different assumptions about what's learnable and how to learn it.

\begin{figure}[htbp]
\centering
\begin{tikzpicture}[
    cluster/.style={rectangle, draw, fill=gray!10, rounded corners, minimum width=4cm, minimum height=2.6cm},
    cat/.style={rectangle, draw, fill=white, text width=3.2cm, text centered, font=\tiny, rounded corners, minimum height=0.35cm},
    title/.style={font=\scriptsize\bfseries},
    scale=0.85, transform shape
]
    % Cluster 1: Core Algorithms (top left)
    \node[cluster] (c1) at (0,4.5) {};
    \node[title] at (0,5.6) {Core Algorithms};
    \node[cat, fill=blue!15] at (0,5) {1. Policy-Based};
    \node[cat, fill=blue!15] at (0,4.5) {2. Value-Based};
    \node[cat, fill=blue!15] at (0,4) {3. Model-Based};

    % Cluster 2: Learning Paradigms (top right)
    \node[cluster] (c2) at (5.5,4.5) {};
    \node[title] at (5.5,5.6) {Learning Paradigms};
    \node[cat, fill=green!15] at (5.5,5) {4. RLHF \& Alignment};
    \node[cat, fill=green!15] at (5.5,4.5) {5. Imitation Learning};
    \node[cat, fill=green!15] at (5.5,4) {6. Offline RL};

    % Cluster 3: Advanced Methods (middle left)
    \node[cluster, minimum height=3.2cm] (c3) at (0,1.2) {};
    \node[title] at (0,2.6) {Advanced Methods};
    \node[cat, fill=orange!15] at (0,2.1) {7. Multi-Agent RL};
    \node[cat, fill=orange!15] at (0,1.6) {8. Hierarchical RL};
    \node[cat, fill=orange!15] at (0,1.1) {9. Meta-Learning};
    \node[cat, fill=orange!15] at (0,0.6) {10. Safe RL};
    \node[cat, fill=orange!15] at (0,0.1) {11. Exploration};

    % Cluster 4: Applications (middle right)
    \node[cluster, minimum height=3.2cm] (c4) at (5.5,1.2) {};
    \node[title] at (5.5,2.6) {Applications};
    \node[cat, fill=purple!15] at (5.5,2.1) {12. Agentic RL};
    \node[cat, fill=purple!15] at (5.5,1.6) {13. Robotics};
    \node[cat, fill=purple!15] at (5.5,1.1) {14. Reasoning};
    \node[cat, fill=purple!15] at (5.5,0.6) {15. Security};
    \node[cat, fill=purple!15] at (5.5,0.1) {16. Benchmarks};

    % Infrastructure cluster (bottom)
    \node[cluster, minimum width=9.5cm, minimum height=1.1cm] (c5) at (2.75,-1.5) {};
    \node[title] at (2.75,-1) {Infrastructure \& Meta};
    \node[cat, fill=cyan!15] at (0,-1.5) {17. Frameworks};
    \node[cat, fill=cyan!15] at (2.75,-1.5) {18. Surveys};
    \node[cat, fill=cyan!15] at (5.5,-1.5) {19. Open-Ended};

    % Arrows - direct vertical connections (no overlap)
    \draw[->, thick, gray!70] (c1.south) -- (c3.north);
    \draw[->, thick, gray!70] (c2.south) -- (c4.north);
    \draw[->, thick, gray!70] (c3.south) -- ++(0,-0.15) -| ([xshift=-1.5cm]c5.north);
    \draw[->, thick, gray!70] (c4.south) -- ++(0,-0.15) -| ([xshift=1.5cm]c5.north);
\end{tikzpicture}
\caption{Taxonomy of 19 RL research categories organized into five clusters: \textbf{Core Algorithms} (policy/value/model-based methods), \textbf{Learning Paradigms} (how agents acquire behavior), \textbf{Advanced Methods} (multi-agent, hierarchical, safe exploration), \textbf{Applications} (domain-specific implementations), and \textbf{Infrastructure} (frameworks, benchmarks, surveys). Arrows indicate foundational dependencies.}
\label{fig:rl_taxonomy}
\source{Author's synthesis based on surveyed literature}
\end{figure}

\subsubsection*{1. Core Methods}
At its heart, RL is about credit assignment: figuring out which past action caused the current outcome. Core methods solve the basic loop of trial, error, and feedback---whether estimating state values (Value-Based) or directly optimizing action probabilities (Policy-Based). Figure~\ref{fig:core_methods} sketches the two main families.

\begin{figure}[htbp]
\centering
\begin{tikzpicture}[
    box/.style={rectangle, draw, text width=2.8cm, text centered, rounded corners, minimum height=0.8cm, font=\scriptsize},
    arrow/.style={->, >=stealth, thick},
    scale=0.95, transform shape
]
    % Value-Based side
    \node[box, fill=blue!20] (vb_title) at (0,3) {\textbf{Value-Based}};
    \node[box, fill=blue!10] (state) at (0,2) {State $s$};
    \node[box, fill=blue!10] (qfunc) at (0,1) {$Q(s,a)$ Network};
    \node[box, fill=blue!10] (argmax) at (0,0) {$\arg\max_a Q(s,a)$};
    \node[box, fill=blue!10] (action_v) at (0,-1) {Action $a$};

    \draw[arrow] (state) -- (qfunc);
    \draw[arrow] (qfunc) -- (argmax);
    \draw[arrow] (argmax) -- (action_v);

    % Policy-Based side
    \node[box, fill=green!20] (pb_title) at (5,3) {\textbf{Policy-Based}};
    \node[box, fill=green!10] (state_p) at (5,2) {State $s$};
    \node[box, fill=green!10] (policy) at (5,1) {$\pi_\theta(a|s)$};
    \node[box, fill=green!10] (sample) at (5,0) {Sample $a \sim \pi$};
    \node[box, fill=green!10] (action_p) at (5,-1) {Action $a$};

    \draw[arrow] (state_p) -- (policy);
    \draw[arrow] (policy) -- (sample);
    \draw[arrow] (sample) -- (action_p);

    % Examples
    \node[font=\tiny, text width=2.5cm, align=center] at (0,-1.8) {DQN, Rainbow,\\Distributional RL};
    \node[font=\tiny, text width=2.5cm, align=center] at (5,-1.8) {PPO, SAC,\\REINFORCE};
\end{tikzpicture}
\caption{Value-based methods learn $Q(s,a)$ and derive policy via $\arg\max$; policy-based methods directly parameterize $\pi_\theta(a|s)$ and sample actions stochastically.}
\label{fig:core_methods}
\end{figure}

\subsubsection*{2. RLHF \& Alignment}
Optimization needs a clear objective, but human intent is too complex to write as math equations. Alignment techniques treat human preference as ground truth. Instead of hard-coding rewards, we approximate what humans want by training models to predict human preferences, allowing optimization for fuzzy concepts like "helpfulness."

\subsubsection*{3. Multi-Agent RL}
No single entity can handle everything. Multi-agent systems tackle complex problems by distributing work across specialized roles. The challenge shifts from individual optimization to finding equilibrium strategies (like Nash Equilibrium) where multiple optimizers co-adapt---mirroring how human organizations trade coordination costs for total capability.

\subsubsection*{4. Hierarchical RL}
Complex tasks suffer from the ``curse of dimensionality'' over time. Planning a trip needs high-level decisions (destination) and low-level execution (stepping on the gas). Hierarchical RL abstracts time: a ``manager'' sets abstract goals over long horizons while ``workers'' execute atomic actions to achieve them, reducing the search space for long tasks. Figure~\ref{fig:hierarchical} sketches this two-tier setup.

\begin{figure}[htbp]
\centering
\begin{tikzpicture}[
    box/.style={rectangle, draw, text width=2.2cm, text centered, rounded corners, minimum height=0.6cm, font=\scriptsize},
    arrow/.style={->, >=stealth, thick},
    scale=0.9, transform shape
]
    % High level (top row)
    \node[box, fill=purple!30] (manager) at (0,2) {Manager\\$\pi_{high}(g|s)$};
    \node[box, fill=purple!10] (goal) at (3.5,2) {Subgoal $g$};
    \node[box, fill=gray!20] (env) at (7,2) {Environment};

    % Low level (bottom row)
    \node[box, fill=blue!30] (worker) at (3.5,0.5) {Worker\\$\pi_{low}(a|s,g)$};
    \node[box, fill=blue!10] (action) at (7,0.5) {Primitive\\Action $a$};

    % Forward flow arrows
    \draw[arrow] (manager) -- node[above, font=\tiny] {every $k$ steps} (goal);
    \draw[arrow] (goal) -- (worker);
    \draw[arrow] (worker) -- node[above, font=\tiny] {every step} (action);
    \draw[arrow] (action) -- (env);
    
    % Feedback arrow - routed below to avoid crossing boxes
    \draw[arrow, dashed] (env.north) -- ++(0,0.5) -| node[pos=0.25, above, font=\tiny] {reward $r$} (manager.north);

    % Examples
    \node[font=\tiny, text width=5cm, align=center] at (3.5,-0.6) {Options Framework, HAM, HIRO, Feudal Networks};
\end{tikzpicture}
\caption{Hierarchical RL: manager policy selects subgoals every $k$ steps; worker policy executes primitive actions to achieve them. Temporal abstraction reduces effective horizon.}
\label{fig:hierarchical}
\end{figure}

\subsubsection*{5. Safe \& Constrained RL}
Standard RL assumes maximizing reward is the only goal, but in reality, the path matters. Safe RL introduces constrained optimization---agents must maximize utility while staying within safety boundaries. This stops the agent from gaming the objective via risky or unacceptable shortcuts.

\subsubsection*{6. Curiosity \& Exploration}
In environments with sparse feedback, reward-driven agents learn nothing because they never stumble upon success. Exploration methods inject an intrinsic drive---optimizing not just for external reward but for shrinking uncertainty or chasing novelty. This lets agents piece together world models even absent explicit tasks. Figure~\ref{fig:exploration} shows how intrinsic rewards patch the sparse-feedback problem.

\begin{figure}[htbp]
\centering
\begin{tikzpicture}[
    box/.style={rectangle, draw, text width=2cm, text centered, rounded corners, minimum height=0.6cm, font=\scriptsize},
    arrow/.style={->, >=stealth, thick},
    scale=0.9, transform shape
]
    % Agent
    \node[box, fill=blue!20] (agent) at (0,0) {Agent};
    \node[box, fill=gray!20] (env) at (3,0) {Environment};

    % Rewards
    \node[box, fill=red!20] (ext) at (6,0.5) {Extrinsic $r_e$\\{\tiny(sparse)}};
    \node[box, fill=green!20] (int) at (6,-0.5) {Intrinsic $r_i$\\{\tiny(curiosity)}};
    \node[box, fill=orange!20] (total) at (9,0) {$r = r_e + \beta r_i$};

    \draw[arrow] (agent) -- node[above, font=\tiny] {action} (env);
    \draw[arrow] (env) -- (ext);
    \draw[arrow] (env) -- (int);
    \draw[arrow] (ext) -- (total);
    \draw[arrow] (int) -- (total);
    \draw[arrow] (total.south) -- ++(0,-0.8) -| node[near end, below, font=\tiny] {learn} (agent.south);

    % Methods
    \node[font=\tiny, text width=5cm, align=center] at (4.5,-1.5) {RND: prediction error | ICM: forward model error | Go-Explore: state archive};
\end{tikzpicture}
\caption{Exploration via intrinsic motivation: curiosity rewards ($r_i$) supplement sparse external rewards ($r_e$), enabling learning in hard-exploration domains.}
\label{fig:exploration}
\end{figure}

\subsubsection*{7. Model-Based RL}
Trial-and-error in the real world is slow and expensive. Model-based methods learn environment dynamics, allowing agents to ``dream'' or simulate futures without acting. By shifting computation from physical interaction to internal simulation, agents squeeze far more learning out of each interaction---swapping real-world trials for simulated ones. Figure~\ref{fig:model_based} shows the model-based learning loop.

\begin{figure}[htbp]
\centering
\begin{tikzpicture}[
    box/.style={rectangle, draw, text width=2.2cm, text centered, rounded corners, minimum height=0.7cm, font=\scriptsize},
    arrow/.style={->, >=stealth, thick},
    scale=0.9, transform shape
]
    % Real world
    \node[box, fill=red!20] (real) at (0,0) {Real World};
    \node[box, fill=orange!20] (data) at (3,0) {Experience\\$(s,a,r,s')$};
    \node[box, fill=yellow!20] (model) at (6,0) {Learned Model\\$\hat{T}(s'|s,a)$};

    % Imagination
    \node[box, fill=cyan!20] (imagine) at (6,-1.5) {Imagined\\Rollouts};
    \node[box, fill=green!20] (plan) at (3,-1.5) {Planning /\\Policy Update};
    \node[box, fill=blue!20] (policy) at (0,-1.5) {Improved\\Policy $\pi$};

    \draw[arrow] (real) -- node[above, font=\tiny] {interact} (data);
    \draw[arrow] (data) -- node[above, font=\tiny] {train} (model);
    \draw[arrow] (model) -- (imagine);
    \draw[arrow] (imagine) -- node[below, font=\tiny] {simulate} (plan);
    \draw[arrow] (plan) -- (policy);
    \draw[arrow] (policy) -- (real);

    % Label
    \node[font=\tiny, text width=4cm, align=center] at (3,-2.5) {World Models, Dreamer, MuZero};
\end{tikzpicture}
\caption{Model-based RL loop: learn dynamics model from real experience, then use imagined rollouts for planning, cutting data requirements 10--100$\times$.}
\label{fig:model_based}
\end{figure}

\subsubsection*{8. Imitation Learning}
Searching from scratch is inefficient when solutions already exist. Imitation learning treats expert behavior as a strong prior, narrowing the search to "proven effective strategies." This bootstraps learning, avoiding rediscovery of basic competencies.

\subsubsection*{9. Agentic RL}
This applies control theory to language and digital tools. Unlike physical control, the "action space" involves generating tokens or API calls, and the "state space" is the context window. The twist: reasoning is no longer static output---it is a latent policy that pokes at the digital environment until the answer clicks into place.

\subsubsection*{10. Offline RL}
Online interaction is often impossible due to safety or cost (medical treatment, power grids). Offline RL tackles the counterfactual problem: learning solid policies from frozen historical data. The hard part is staying conservative---avoiding actions that look good but lead to unseen states.

\subsubsection*{11. Meta \& Continual RL}
Standard models are brittle---they master one task but fail if rules change slightly. Meta-learning shifts the target from ``performance on Task A'' to ``speed of adaptation to Task B.'' It optimizes the learning algorithm itself, creating agents that are not just skilled but adaptable to changing environments. Figure~\ref{fig:meta_rl} shows the meta-learning loop.

\begin{figure}[htbp]
\centering
\begin{tikzpicture}[
    box/.style={rectangle, draw, text width=2.4cm, text centered, rounded corners, minimum height=0.6cm, font=\scriptsize},
    arrow/.style={->, >=stealth, thick},
    scale=0.85, transform shape
]
    % Meta level (left)
    \node[box, fill=red!20, minimum height=1.8cm] (meta) at (0,1) {Meta-Learner\\$\theta$};

    % Task distribution (middle-left) - stacked vertically
    \node[box, fill=orange!20] (task1) at (3.5,2) {Task 1};
    \node[box, fill=orange!20] (task2) at (3.5,1) {Task 2};
    \node[box, fill=orange!20] (taskn) at (3.5,0) {Task $n$};

    % Adapted parameters (middle-right)
    \node[box, fill=green!20, minimum height=1.8cm] (adapt) at (7,1) {Adapted\\$\theta'_i$};

    % Evaluation (right)
    \node[box, fill=blue!20, minimum height=1.8cm] (eval) at (10.2,1) {Fast\\Adaptation};

    % Forward arrows - fan out from meta to tasks
    \draw[arrow] (meta.east) -- (task1.west);
    \draw[arrow] (meta.east) -- (task2.west);
    \draw[arrow] (meta.east) -- (taskn.west);
    
    % Fan in from tasks to adapted
    \draw[arrow] (task1.east) -- (adapt.west);
    \draw[arrow] (task2.east) -- (adapt.west);
    \draw[arrow] (taskn.east) -- (adapt.west);
    
    % Forward to evaluation
    \draw[arrow] (adapt) -- node[above, font=\tiny] {few shots} (eval);

    % Feedback loop - routed below with clear path
    \draw[arrow, dashed, gray] (eval.south) -- ++(0,-0.8) -| node[pos=0.75, below, font=\tiny] {meta-gradient} (meta.south);

    % Method examples
    \node[font=\tiny, text width=5cm, align=center] at (5.1,-1.2) {MAML, RL$^2$, PEARL};
\end{tikzpicture}
\caption{Meta-RL: learn initialization $\theta$ that enables fast adaptation to new tasks with few gradient steps. Outer loop optimizes ``learning to learn.''}
\label{fig:meta_rl}
\end{figure}

\subsubsection*{12. Open-Ended Learning \& Generalization}
Fixed benchmarks lead to overfitting. Open-ended learning mimics evolution, where environment and agent co-evolve. The system generates its own curriculum of increasingly difficult tasks, ensuring capabilities must constantly expand. The goal isn't solving a task but indefinitely increasing solvable complexity.

\subsubsection*{13. Systems \& Modular RL}
The bottleneck in modern AI is often data throughput, not algorithms. Systems RL tackles the engineering physics of training: moving massive experience data from environments to learners efficiently. It treats the training loop as a high-performance distributed computing problem.

\subsubsection*{14. Multi-Objective RL}
Real problems rarely have a single success metric. Agents must balance speed against accuracy, or profit against risk. Multi-objective RL rejects scalar rewards, instead solving for the Pareto frontier---all optimal trade-offs. This allows dynamic selection between conflicting goals at runtime.

\subsubsection*{15. Embodied \& Robotics RL}
Intelligence needs grounding in physical reality. Embodied RL deals with noisy, continuous, unforgiving physics. The hard part is the ``sim-to-real'' gap: learning policies in perfect simulation that survive real-world friction, latency, and sensor noise.

\subsubsection*{16. Surveys \& Overviews}
Science needs consolidation. As research explodes, these works serve as map-makers, synthesizing scattered discoveries into coherent frameworks. They chart what we know, spotlight blind spots, and stake out the questions that will drive the next wave of work.

\subsubsection*{17. Reasoning \& Search}
Standard LLMs use "System 1" thinking---fast, heuristic pattern matching. Reasoning RL adds "System 2" capabilities by treating thought generation as a search problem. Models learn to navigate reasoning trees, backtracking and verifying, turning compute time into logical accuracy.

\subsubsection*{18. LLM Security \& Red Teaming}
Complex systems have hidden failure modes. Red teaming applies adversarial optimization to find these weaknesses. By treating attackers as RL agents trying to maximize target model failure, we automatically discover and patch vulnerabilities humans might never think to test.

\subsubsection*{19. Benchmarks \& Evaluations}
Progress demands measurement. Without agreed-upon yardsticks, we cannot distinguish genuine advances from cherry-picked results. Good benchmarks act as the community's shared exam---forcing models to prove their worth on problems they have not memorized.

\section{Policy Gradient Methods: PPO and TRPO}
\label{sec:ppo}
Policy optimization has a cliff-edge problem: a single bad update can destroy the policy, and because the data distribution depends on the policy, it can never recover. PPO solves this by enforcing a "trust region"—a mathematical guardrail ensuring we never step too far into the unknown. Trust-region methods prevent catastrophic performance degradation by limiting policy updates. PPO \cite{schulman2017ppo} makes this practical with a clipping mechanism:
\begin{equation}
L^{CLIP}(\theta) = \mathbb{E}_t \left[ \min(r_t(\theta) \hat{A}_t, \text{clip}(r_t(\theta), 1-\epsilon, 1+\epsilon) \hat{A}_t) \right]
\end{equation}
where $r_t(\theta)$ is the probability ratio between new and old policies, and $\hat{A}_t$ estimates advantage. The clipping threshold $\epsilon$ (typically 0.1-0.2) bounds how much the policy can change per iteration. Despite proven stability, PPO needs lots of memory---maintaining actor, critic, reference, and reward models---and is sensitive to learning rate, clipping range, and advantage estimation.

\begin{tcolorbox}[colback=gray!5!white, colframe=gray!60!black, title=\textbf{Worked Example: PPO Clipping in Action}]
\small
Suppose $\epsilon = 0.2$ and the advantage $\hat{A}_t = +3$ (a good action). The probability ratio $r_t(\theta)$ measures how much more likely the new policy makes this action:

\begin{center}
\begin{tabular}{lccl}
\toprule
\textbf{Scenario} & $r_t(\theta)$ & \textbf{Clipped?} & \textbf{Objective term / gradient w.r.t.\ $r_t$} \\
\midrule
Small update & 1.1 & No (in $[0.8, 1.2]$) & $1.1 \times 3 = 3.3$; gradient $\hat{A}_t = 3$ \\
Large update & 1.5 & Yes $\rightarrow$ 1.2 & $1.2 \times 3 = 3.6$ (constant); gradient $0$ \\
\bottomrule
\end{tabular}
\end{center}

\textbf{Key insight:} The numbers $3.3$, $3.6$, and $4.5$ are values of the surrogate objective term $r_t\hat{A}_t$, not gradients. Without clipping, the objective at $r_t = 1.5$ would be $1.5 \times 3 = 4.5$ and its gradient with respect to $r_t$ would remain $\hat{A}_t = 3$, so the optimizer would keep pushing the ratio upward. With clipping, the $\min$ selects the constant clipped term $3.6$ once $r_t > 1+\epsilon$ (for $\hat{A}_t>0$), so the gradient from this sample is $0$ and the policy stops moving further in that direction.
\end{tcolorbox}

\section{DPO: Skip the Reward Model Entirely}
\label{sec:dpo}
\textbf{Transition:} PPO's multi-model architecture works but creates engineering headaches. What if we could eliminate the reward model entirely?

Building an independent reward model is not always necessary if the objective can be solved algebraically. The authors of the DPO framework \cite{rafailov2023direct} demonstrated that the RLHF objective can be restructured until the reward model cancels out. What's left is a loss you can calculate based on pairs of preferences---no intermediate scoring step:
\begin{equation}
L_{\text{DPO}} = -\mathbb{E}_{(x,y_w,y_l)} \left[ \log \sigma \left( \beta_{\text{DPO}} \log \frac{\pi_\theta(y_w|x)}{\pi_{\text{ref}}(y_w|x)} - \beta_{\text{DPO}} \log \frac{\pi_\theta(y_l|x)}{\pi_{\text{ref}}(y_l|x)} \right) \right]
\end{equation}
Here $y_w$ is the response annotators liked better and $y_l$ is the one they didn't; $\beta$ controls how aggressively you push toward preferred responses (crank it too high and you'll overfit to idiosyncrasies in your preference data).

The appeal is obvious: simpler pipeline, fewer moving parts, less stuff that can break. The downside is equally clear---you're locked into whatever preferences you collected at the start. DPO can't explore; it just learns to mimic the patterns in your dataset. For tasks where the optimal strategy isn't well-represented in training, that's a real limitation.

\begin{tcolorbox}[colback=gray!5!white, colframe=gray!60!black, title=\textbf{Worked Example: DPO $\beta$ Parameter Tuning}]
\small
The $\beta$ parameter controls preference strength. Consider a prompt where annotators preferred response A over B:

\begin{center}
\begin{tabular}{lccc}
\toprule
$\beta$ & \textbf{Effect} & \textbf{Risk} & \textbf{Use Case} \\
\midrule
0.1 (low) & Weak preference signal & Slow convergence & Noisy/inconsistent labels \\
0.5 (typical) & Balanced learning & --- & Most applications \\
1.0+ (high) & Strong preference push & Overfitting to quirks & Clean, consistent data \\
\bottomrule
\end{tabular}
\end{center}

\textbf{For practical implementation:} Start with $\beta = 0.1$--$0.5$.
 If the model ignores preferences, increase $\beta$. If it becomes repetitive or loses fluency, decrease $\beta$. Monitor KL divergence from the reference model---if it spikes, $\beta$ is too high.
\end{tcolorbox}

\section{GRPO: The Memory-Efficient Alternative That Took Over}
\label{sec:grpo}
\textbf{Transition:} DPO simplifies training but sacrifices exploration. GRPO takes a different path---keeping online learning while dramatically reducing memory requirements.

Standard RL uses a ``critic'' network to estimate how good a state is, which helps reduce variance in gradient estimates. But here's the thing: you don't actually need to know how good something is in absolute terms---you just need to know whether it's better or worse than the alternatives. The DeepSeek-R1 team \cite{deepseek2025r1} ran with this idea and stripped out the critic entirely. Their approach, GRPO, works like this: create multiple outputs for a given query, score them all
 (using whatever reward signal you have---test cases, a learned model, even just rule-based checks), then rank them within the batch. Responses that beat the batch average get reinforced; responses below average get penalized. That's it.

The implications for memory usage are significant. PPO needs to keep four models in VRAM: the policy, a reference copy, a critic, and often a reward model. GRPO eliminates the critic entirely, requiring only the policy and a reference model for the KL penalty (2$\times$ vs 4$\times$ model size), plus a scoring function. For a 70B parameter model, that's a 50\% memory reduction that can make the difference between needing multiple nodes and fitting on a single machine.

Did it work? DeepSeek-R1 hit 79.8\% on AIME 2024 math competition problems \cite{deepseek2025r1}, putting it in the same ballpark as human medalists. And they did it without the infrastructure burden that made PPO-style training prohibitively expensive for most labs.

\subsection{The GRPO Family Tree}
Once GRPO proved itself, people started tweaking it for specific use cases:

\begin{itemize}
    \item \textbf{Flow-GRPO} \cite{li2025flowgrpo}: Trains the planner of a multi-module agent (AgentFlow) on-policy by broadcasting a single verifiable trajectory-level outcome to every turn, turning multi-turn optimization into a sequence of single-turn updates with group-normalized advantages.

    \item \textbf{M-GRPO} \cite{hong2025mgrpo}: Built for multi-agent setups where a ``manager'' delegates subtasks to ``workers.'' Each level in the hierarchy keeps its own baseline, so agents can specialize without fighting over a shared reward signal.

    \item \textbf{GRPO-TCR:} Adds a penalty for malformed tool calls. Generate an invalid API invocation and your advantage score takes a hit---a simple way to teach syntactic correctness without separate supervised training.

\end{itemize}

\noindent\textbf{Critic-free refinements of GRPO:}
\begin{itemize}
    \item \textbf{DAPO (Decoupled Clip and Dynamic Sampling Policy Optimization)} \cite{bytedance2025dapo}: A critic-free GRPO-style method that keeps clipping but decouples the lower and upper clip ranges (``Clip-Higher'') to counter entropy collapse, adds dynamic sampling that filters prompts whose groups are all-correct or all-wrong, and uses token-level loss normalization and overlong-response shaping. It reports 50 points on AIME 2024 with Qwen2.5-32B.
\end{itemize}

\noindent\textbf{Additional Novel Algorithms (2024--2025):} Beyond GRPO and PPO variants, new algorithms have emerged:
\begin{itemize}
    \item \textbf{Tree-GRPO} \cite{ji2025treegrpo}: Applies tree search to LLM agent reinforcement learning, outperforming chain-based approaches over a selection of 11 datasets and 3 QA task types.
 Multi-child branching enables multiple policy updates per forward pass with fine-grained credit assignment via reward backpropagation through the tree structure.

    \item \textbf{ReMax} \cite{li2024remax}: A REINFORCE-style method for RLHF that drops PPO's value model and uses the reward of a greedily decoded response as the baseline; it predates GRPO and is not a GRPO simplification.
\end{itemize}

\subsection{Algorithm Distribution in Current Research}
We do not report a numeric breakdown of algorithm families. An earlier draft gave counts over 80 ``cataloged method variants'' (31 GRPO, 18 PPO, 8 DPO, 23 custom) with an inter-rater $\kappa$, but no catalog listing those variants or the agreement data was retained, so the counts cannot be checked and are withdrawn. Qualitatively, GRPO-style critic-free methods are frequent among the 2024--2025 papers in Appendix~\ref{app:paper_mapping}.

\noindent\textbf{Architecture Paradigms:} The studies use three recurring architectural patterns (no retained tally supports percentages):
\begin{itemize}
    \item \textbf{Modular architectures:} separate components for perception, planning, and execution.
    \item \textbf{End-to-end architectures:} unified models mapping observations directly to actions.
    \item \textbf{Hierarchical architectures:} multi-level policies with manager--worker dynamics.
\end{itemize}

\subsection{Comparative Analysis of Alignment Algorithms}
Table~\ref{tab:algorithm_comparison} lines up the three main alignment algorithms across the dimensions that matter for agentic AI work.

\begin{table}[h]
\centering
\caption{Head-to-Head Rundown of PPO, DPO, and GRPO for Agentic Fine-Tuning}
\label{tab:algorithm_comparison}
\small
\begin{tabular}{@{}p{2.8cm}p{3.8cm}p{3.8cm}p{3.8cm}@{}}
\toprule
\textbf{Dimension} & \textbf{PPO} & \textbf{DPO} & \textbf{GRPO} \\
\midrule
\textbf{Core Mechanism} & Trust-region policy gradient with clipped surrogate & Direct preference mapping via log-ratio &
 Group-normalized advantage estimation \\
\addlinespace
\textbf{Models Required} & Actor, Critic, Reward, Reference (4) & Policy, Reference (2) & Policy, Reference (2)\textsuperscript{$\dagger$} \\
\addlinespace
\textbf{Training Mode} & Online (env interaction) & Offline (preference data) & Online (self-generated) \\
\addlinespace
\textbf{Sample Efficiency} & Low ($\sim$10K interactions) & High (preference pairs) & Medium (group sampling) \\
\addlinespace
\textbf{Memory Cost} & High (4$\times$ model size) & Medium (2$\times$ model size) & Low (2$\times$ model size)\textsuperscript{*} \\
\addlinespace
\textbf{Exploration} & Good (stochastic policy) & Limited (no exploration) & Moderate (group diversity) \\
\addlinespace
\textbf{Stability} & High with proper $\epsilon$ & Very High & High \\
\addlinespace
\textbf{Key Hyperparams} & $\epsilon=0.2$, $\beta_{\text{KL}}=0.01$ & $\beta=0.1$ & Group size $G=8$ \\
\addlinespace
\textbf{Best For} & Complex multi-step tasks & Static preference alignment & Reasoning-heavy tasks \\
\addlinespace
\textbf{SOTA Results} & InstructGPT, Llama-2-Chat & Zephyr & DeepSeek-R1 (79.8\% AIME) \\
\bottomrule
\end{tabular}

\vspace{0.2cm}
\small\textsuperscript{*}GRPO requires policy + reference model (2$\times$), eliminating only the critic; still 50\% reduction vs. PPO.\\
\textsuperscript{$\dagger$}GRPO uses a programmatic/rule-based reward function rather than a learned reward model.
\end{table}

\begin{figure}[htbp]
\centering
\begin{tikzpicture}[scale=0.9]
    % Radar chart for algorithm comparison
    \def\n{6} % number of dimensions
    \def\radius{3}

    % Draw grid FIRST (background)
    \foreach \r in {0.6, 1.2, 1.8, 2.4, 3.0} {
        \draw[gray!40, thin] (90:\r) \foreach \i in {1,...,\n} { -- ({90-\i*360/\n}:\r)} -- cycle;
    }

    % Axes
    \foreach \i in {1,...,\n} {
        \pgfmathsetmacro{\angle}{90 - (\i-1)*360/\n}
        \draw[gray!50, thin] (0,0) -- (\angle:\radius);
    }

    % PPO data - separate fill and stroke for clarity
    \fill[blue!25, opacity=0.6]
        (90:0.6) -- (30:3.0) -- (-30:2.4) -- (-90:1.2) -- (-150:1.2) -- (150:2.4) -- cycle;
    \draw[blue!80, thick]
        (90:0.6) -- (30:3.0) -- (-30:2.4) -- (-90:1.2) -- (-150:1.2) -- (150:2.4) -- cycle;

    % DPO data
    \fill[red!25, opacity=0.5]
        (90:1.8) -- (30:0.6) -- (-30:3.0) -- (-90:3.0) -- (-150:2.4) -- (150:1.8) -- cycle;
    \draw[red!80, thick]
        (90:1.8) -- (30:0.6) -- (-30:3.0) -- (-90:3.0) -- (-150:2.4) -- (150:1.8) -- cycle;

    % GRPO data
    \fill[green!30, opacity=0.5]
        (90:3.0) -- (30:1.8) -- (-30:2.4) -- (-90:2.4) -- (-150:3.0) -- (150:3.0) -- cycle;
    \draw[green!60!black, thick]
        (90:3.0) -- (30:1.8) -- (-30:2.4) -- (-90:2.4) -- (-150:3.0) -- (150:3.0) -- cycle;

    % Dimension labels with white background (placed LAST, on top of everything)
    \node[font=\scriptsize, fill=white, inner sep=2pt, rounded corners] at (90:\radius+0.9) {Memory Efficiency};
    \node[font=\scriptsize, fill=white, inner sep=2pt, rounded corners] at (30:\radius+1.0) {Exploration};
    \node[font=\scriptsize, fill=white, inner sep=2pt, rounded corners] at (-30:\radius+0.9) {Stability};
    \node[font=\scriptsize, fill=white, inner sep=2pt, rounded corners] at (-90:\radius+1.0) {Sample Efficiency};
    \node[font=\scriptsize, fill=white, inner sep=2pt, rounded corners] at (-150:\radius+0.9) {Simplicity};
    \node[font=\scriptsize, fill=white, inner sep=2pt, rounded corners] at (150:\radius+1.1) {SOTA Performance};

    % Scale labels with white background
    \node[font=\tiny, fill=white, inner sep=1pt] at (75:0.5) {1};
    \node[font=\tiny, fill=white, inner sep=1pt] at (75:1.5) {3};
    \node[font=\tiny, fill=white, inner sep=1pt] at (75:2.5) {5};

    % Legend with proper box
    \draw[thick, fill=white, rounded corners] (4.3, 2.5) rectangle (6.5, 0.3);
    \node[draw=blue!80, fill=blue!25, minimum width=0.35cm, minimum height=0.25cm] at (4.7, 2.1) {};
    \node[right, font=\scriptsize] at (5.0, 2.1) {PPO};
    \node[draw=red!80, fill=red!25, minimum width=0.35cm, minimum height=0.25cm] at (4.7, 1.4) {};
    \node[right, font=\scriptsize] at (5.0, 1.4) {DPO};
    \node[draw=green!60!black, fill=green!30, minimum width=0.35cm, minimum height=0.25cm] at (4.7, 0.7) {};
    \node[right, font=\scriptsize] at (5.0, 0.7) {GRPO};
\end{tikzpicture}
\caption{Radar comparison of PPO, DPO, and GRPO across six dimensions (scale: 1--5, higher is better). GRPO lands in the sweet spot---low RAM footprint yet strong results---while PPO shines at exploration and DPO at stability and data efficiency. This visualization explains why GRPO now dominates in GPU-limited labs.}
\label{fig:algorithm_radar}
\end{figure}

\subsection{Theoretical Properties}

We formalize key properties that distinguish the three alignment algorithms:

\begin{ruleofthumb}[Memory-Exploration Trade-off]
\label{prop:tradeoff}
Let $M_{\text{PPO}}$, $M_{\text{DPO}}$, $M_{\text{GRPO}}$ denote memory requirements and $E_{\text{PPO}}$, $E_{\text{DPO}}$, $E_{\text{GRPO}}$ denote exploration capacity. As an order-of-magnitude comparison:
\begin{enumerate}
    \item $M_{\text{PPO}} \approx 4M_{\text{base}} > M_{\text{DPO}} \approx M_{\text{GRPO}} \approx 2M_{\text{base}}$
    \item $E_{\text{PPO}} > E_{\text{GRPO}} > E_{\text{DPO}}$
\end{enumerate}
where $M_{\text{base}}$ is the base model memory footprint. These multipliers (4$\times$, 2$\times$, 2$\times$) are approximate and assume distinct resident networks (policy, reference, reward, value) without aggressive sharing or offloading; real memory footprints vary by implementation. PPO explores the most via online environment interaction but needs actor, critic, learned reward, and reference models. DPO sidesteps exploration altogether---it learns only from a fixed preference dataset. GRPO lands in the middle on exploration via group sampling diversity; it eliminates the critic and learned reward model but retains the reference model for KL regularization. \textit{Note: GRPO can be implemented in online or offline mode; here we assume online for comparison.}
\end{ruleofthumb}

\begin{remark}
The trade-off in Rule of Thumb~\ref{prop:tradeoff} explains why GRPO hit top reasoning marks (79.8\% on AIME 2024) without a huge compute budget: reasoning tasks benefit more from self-improvement through group comparison than from exhaustive environment exploration, while the reduced memory allows training larger models.
\end{remark}

\begin{proposition}[GRPO Memory Efficiency]
\label{prop:grpo_memory}
Let $G$ denote a group of $n$ generated responses $\{o_1, \ldots, o_n\}$ for query $q$, with rewards $\{r_1, \ldots, r_n\}$ from a verifiable reward function $R(\cdot)$. In its simplest on-policy form (one gradient step per batch, so the PPO ratio equals 1 and clipping is inactive), the reward-driven part of the GRPO gradient is
\begin{equation}
\nabla_\theta \mathcal{L}_{\text{GRPO}} = -\frac{1}{n} \sum_{i=1}^{n} \left( \frac{r_i - \mu_G}{s_G + \epsilon} \right) \nabla_\theta \log \pi_\theta(o_i | q) \;+\; \beta_{\text{KL}}\, \nabla_\theta \mathbb{D}_{\text{KL}}\!\left[\pi_\theta \,\|\, \pi_{\text{ref}}\right],
\end{equation}
where $\mu_G = \frac{1}{n}\sum_j r_j$ and $s_G = \sqrt{\frac{1}{n}\sum_j (r_j - \mu_G)^2}$ are within-group statistics. As in DeepSeekMath~\cite{shao2024deepseekmath}, the KL penalty to the reference policy is a separate loss term; it is not folded into the rewards before normalization. With multiple updates per batch, DeepSeekMath additionally applies the PPO-style ratio and clipping to the advantage term. This formulation eliminates the need for:
\begin{enumerate}
    \item A \textbf{learned reward model} (replaced by verifiable $R(\cdot)$)
    \item A \textbf{critic/value network} (replaced by group normalization baseline)
\end{enumerate}
Consequently, GRPO requires storing only the policy network $\pi_\theta$ and frozen reference $\pi_{\text{ref}}$ during training, versus four networks for PPO (policy, reference, reward model, critic). Note that the verifiable reward function $R(\cdot)$ is either programmatic (e.g., code execution, math verification) or a lightweight external call---it does not add GPU memory overhead as it need not be stored as model weights during gradient computation.
\end{proposition}

\begin{proof}[Derivation Sketch]
Standard policy gradient requires a baseline $b(s)$ to reduce variance: $\nabla_\theta J = \mathbb{E}[(R - b(s)) \nabla_\theta \log \pi_\theta]$. PPO learns $b(s) = V_\phi(s)$ via a separate critic. GRPO instead uses the empirical group mean $\mu_G$ as a zero-cost baseline: for i.i.d.\ samples from the current policy, $\mathbb{E}[\mu_G] = V(q)$. Because $\mu_G$ includes the sample's own reward $r_i$, the baseline is not independent of the sample it scores, and division by $s_G$ rescales each prompt's contribution; the resulting estimator is therefore biased relative to the plain policy gradient (leave-one-out baselines such as RLOO~\cite{ahmadian2024back} avoid the first effect). The substitution is a practical variance-reduction device, not an exact replacement for a learned baseline.
\end{proof}

\begin{proposition}[Trust-Region Motivation: TRPO Guarantee vs.\ PPO Heuristic]
\label{prop:convergence}
TRPO \cite{schulman2015trust} optimizes a surrogate objective under an explicit KL-divergence trust-region constraint:
\begin{equation}
\max_\theta \mathbb{E}_t\left[\frac{\pi_\theta(a_t|s_t)}{\pi_{\theta_{\text{old}}}(a_t|s_t)} \hat{A}_t\right] \quad \text{s.t.} \quad \mathbb{E}_t\big[\text{KL}[\pi_{\theta_{\text{old}}}(\cdot|s_t) \,\|\, \pi_\theta(\cdot|s_t)]\big] \leq \delta
\end{equation}
Under bounded-advantage assumptions, TRPO is accompanied by a conservative improvement guarantee. PPO \cite{schulman2017ppo} replaces the hard KL constraint with a clipped surrogate objective:
\begin{equation}
L^{\text{CLIP}}(\theta) = \mathbb{E}_t\left[\min\left(r_t(\theta)\hat{A}_t, \text{clip}(r_t(\theta), 1-\epsilon, 1+\epsilon)\hat{A}_t\right)\right]
\end{equation}
where $r_t(\theta) = \pi_\theta(a_t|s_t)/\pi_{\theta_{\text{old}}}(a_t|s_t)$. This clipping empirically constrains update magnitude but \textbf{does not provide the same formal monotonic-improvement guarantee as TRPO}. The practical success of PPO comes from its simplicity and empirical stability, not from theoretical guarantees matching TRPO's.
\end{proposition}

\begin{remark}
The distinction matters for practitioners: TRPO offers provable (conservative) improvement per step under its assumptions; PPO offers no such guarantee but trains faster and scales better. Most agentic RL work uses PPO or GRPO for pragmatic reasons, accepting the lack of formal guarantees in exchange for reduced engineering complexity.
\end{remark}

\begin{heuristic}[Scaling Considerations for Agentic Benchmarks]
\label{lem:scaling}
Established pretraining scaling work (e.g., compute-optimal training \cite{hoffmann2022training}) demonstrates predictable tradeoffs between model size and the volume of training data given a static allocation of computational resources for \emph{pretraining loss}.
 However, the scaling behavior of \emph{agentic benchmark success rates} (as opposed to pretraining loss) is less settled and may saturate due to environment/tool bottlenecks. Early observations suggest:
\begin{itemize}
    \item Larger models (e.g., GPT-4 vs.\ GPT-3.5) show disproportionate gains on agentic tasks
    \item Model capacity ($N$) appears more important than data ($D$) or compute ($C$) for tool-use tasks
    \item Success rates may plateau as tasks hit environment-side limits (API latency, simulator fidelity)
\end{itemize}
We treat agentic scaling trends as an empirical question evaluated via H5 (Section~\ref{sec:design_philosophy}) rather than as a formal scaling law. Unlike pretraining loss (which follows well-characterized power-law relationships~\cite{hoffmann2022training}), agentic success rates lack established theoretical foundations---observed trends should be viewed as empirical regularities subject to revision as more data accumulates.
\end{heuristic}

\begin{remark}
This heuristic motivates memory-efficient approaches like GRPO: if model capacity is the primary bottleneck for agentic tasks, then freeing up GPU memory to train larger models may matter more than sophisticated exploration or reward modeling. Further rigorous study is needed to establish whether formal scaling laws exist for agentic capabilities.
\end{remark}

\subsection{Algorithm Selection Guide: When to Use Which}
Drawing on our 60+ paper survey, here is a hands-on guide for picking among the three main alignment algorithms:

\begin{figure}[htbp]
\centering
\begin{tikzpicture}[
    decision/.style={diamond, draw, fill=yellow!20, text width=2.8cm, text centered, inner sep=1pt, font=\small},
    result/.style={rectangle, draw, fill=green!20, text width=2cm, text centered, rounded corners, minimum height=0.8cm, font=\small\bfseries},
    arrow/.style={->, >=stealth, thick},
    level distance=1.8cm,
    sibling distance=4cm
]
    % Root decision
    \node[decision] (q1) {Memory constrained?};

    % Yes branch - GRPO path
    \node[decision, below left=1.5cm and 2cm of q1] (q2) {Verifiable rewards?};
    \node[result, below left=1.2cm and 0.5cm of q2] (grpo) {GRPO};
    \node[result, below right=1.2cm and 0.5cm of q2] (dpo1) {DPO};

    % No branch - PPO/DPO path
    \node[decision, below right=1.5cm and 2cm of q1] (q3) {Need exploration?};
    \node[result, below left=1.2cm and 0.5cm of q3] (dpo2) {DPO};
    \node[result, below right=1.2cm and 0.5cm of q3] (ppo) {PPO};

    % Arrows
    \draw[arrow] (q1) -- node[above left, font=\scriptsize] {Yes} (q2);
    \draw[arrow] (q1) -- node[above right, font=\scriptsize] {No} (q3);
    \draw[arrow] (q2) -- node[left, font=\scriptsize] {Yes} (grpo);
    \draw[arrow] (q2) -- node[right, font=\scriptsize] {No} (dpo1);
    \draw[arrow] (q3) -- node[left, font=\scriptsize] {No} (dpo2);
    \draw[arrow] (q3) -- node[right, font=\scriptsize] {Yes} (ppo);
\end{tikzpicture}
\caption{Algorithm selection decision tree. Start with infrastructure constraints (memory), then consider reward availability and exploration needs.}
\label{fig:algo_decision_tree}
\end{figure}

\noindent\textbf{Choose GRPO When:}
\begin{itemize}
    \item Memory constraints limit your infrastructure (policy plus frozen reference; no critic or learned reward model)
    \item Tasks emphasize reasoning, mathematics, or code generation
    \item You have access to verifiable reward signals (test cases, correctness checkers)
    \item Self-improvement without human feedback is acceptable
    \item Training stability is a priority (group normalization provides natural regularization)
\end{itemize}

\noindent\textbf{Choose PPO When:}
\begin{itemize}
    \item Maximum exploration is critical (novel environments, sparse rewards)
    \item You need fine-grained credit assignment (learned critic)
    \item Multi-step sequential decision making dominates (game playing, robotics)
    \item Infrastructure supports 4$\times$ model memory overhead
    \item You have significant computational resources available for environment interaction
\end{itemize}

\noindent\textbf{Choose DPO When:}
\begin{itemize}
    \item Human preferences are abundant and high-quality
    \item Training must be offline (no environment access during training)
    \item Stability matters more than exploration capability
    \item The task is mainly about style, tone, or preference alignment
    \item Simplicity of implementation is a priority
\end{itemize}

\noindent\textbf{Hybrid Approaches (Recommended for Production):} The best-performing systems in 2024--2025 combine algorithms:
\begin{enumerate}
    \item \textbf{SFT $\rightarrow$ DPO $\rightarrow$ GRPO:} Begin with imitation learning on demonstrations, align preferences via DPO, then sharpen reasoning with GRPO.
    \item \textbf{PPO + GRPO:} Use PPO for environment exploration, then switch to GRPO for computational efficiency during scaling.
    \item \textbf{HRPO (Hybrid RPO):} Combine GRPO's group normalization with DPO's preference pairs in a single objective.
\end{enumerate}

\subsection{Algorithm Pseudocode}
Algorithm~\ref{alg:grpo} presents the GRPO training procedure, which eliminates the need for a separate \emph{value/critic network} by using group-based normalization. A reward signal is still required---this can be verifiable/programmatic (as in DeepSeek-R1's math tasks) or from a learned reward model.

\begin{algorithm}[htbp]
\caption{Group Relative Policy Optimization (GRPO) algorithm. Unlike PPO, GRPO normalizes advantages within each group of completions, \textbf{eliminating the need for a separate value/critic network}. A reward signal $r(\cdot)$ is still required---this can be verifiable/programmatic or from a learned reward model.}
\label{alg:grpo}
\small
\begin{algorithmic}[1]
\Require Policy $\pi_\theta$, reference policy $\pi_{\text{ref}}$, prompts $\mathcal{D}$, reward function $r(\cdot)$, group size $G$, KL coefficient $\beta_{\text{KL}}$, learning rate $\eta$
\Ensure Optimized policy $\pi_{\theta^*}$
\For{each iteration}
    \State Sample prompt $x \sim \mathcal{D}$
    \State Generate $G$ completions $\{y_1, \ldots, y_G\} \sim \pi_\theta(\cdot \mid x)$
    \State Compute rewards $r_i = r(x, y_i)$ for $i \in [G]$
    \State Compute group mean $\mu_G = \frac{1}{G}\sum_{j=1}^{G} r_j$ and std $s_G = \sqrt{\frac{1}{G}\sum_{j=1}^{G}(r_j - \mu_G)^2}$
    \State Normalize advantages $\hat{A}_i = \frac{r_i - \mu_G}{s_G + \epsilon}$ \Comment{$\epsilon = 10^{-8}$ for numerical stability}
    \State Update $\theta \leftarrow \theta + \eta \nabla_\theta \Big[\frac{1}{G}\sum_{i=1}^{G} \hat{A}_i \log \pi_\theta(y_i \mid x) - \beta_{\text{KL}}\, \mathbb{D}_{\text{KL}}[\pi_\theta \| \pi_{\text{ref}}]\Big]$ \Comment{KL is a separate term}
\EndFor
\State \Return $\pi_{\theta^*} \gets \pi_\theta$
\end{algorithmic}
\source{Simplified single-update form adapted from DeepSeek-Math~\cite{shao2024deepseekmath} and DeepSeek-R1~\cite{deepseek2025r1}; the PPO-style ratio and clipping used with multiple updates per batch are omitted.}
\end{algorithm}

\begin{tcolorbox}[colback=red!5!white, colframe=red!70!black, title=\textbf{Common Pitfalls in RL-Based Agent Training}]
\small
\textbf{PPO Pitfalls:}
\begin{itemize}[nosep,leftmargin=*]
    \item \textit{Surrogate reward optimization:} Agent maximizes proxy instead of true objective. \textbf{Fix:} Use process rewards, not just outcome rewards.
    \item \textit{KL explosion:} Policy diverges too far from reference. \textbf{Fix:} Reduce learning rate; increase KL penalty coefficient.
    \item \textit{Value function lag:} Critic estimates become stale. \textbf{Fix:} Train critic more frequently than actor.
\end{itemize}

\textbf{DPO Pitfalls:}
\begin{itemize}[nosep,leftmargin=*]
    \item \textit{Preference inconsistency:} Contradictory labels confuse training. \textbf{Fix:} Filter low-agreement pairs; use lower $\beta$.
    \item \textit{Distribution shift:} Model output drifts from the distribution of collected preferences.
 \textbf{Fix:} Iterative DPO with fresh preferences.
    \item \textit{Mode collapse:} Model produces repetitive outputs. \textbf{Fix:} Add entropy bonus; reduce $\beta$.
\end{itemize}

\textbf{GRPO Pitfalls:}
\begin{itemize}[nosep,leftmargin=*]
    \item \textit{Group size too small:} Noisy advantage estimates. \textbf{Fix:} Use $G \geq 8$; increase for high-variance tasks.
    \item \textit{Reward function bugs:} Verifier has edge cases. \textbf{Fix:} Test reward function extensively before training.
    \item \textit{Collapse to single mode:} All responses become identical. \textbf{Fix:} Increase temperature; add diversity bonus.
\end{itemize}
\end{tcolorbox}

\section{Process Supervision}

Rewarding only final answers creates bad incentives: models might find shortcuts that produce correct outputs through flawed reasoning. Process supervision analyzes individual steps within the reasoning chain:

\textbf{Step-Level Verification:} \citet{lightman2023lets} found that process reward models (PRMs) that score each deduction step substantially outperform outcome reward models (ORMs) on MATH. In their best-of-$n$ sampling experiments, PRMs achieved higher solve rates---the magnitude depends on model size and $n$, but consistently favors step-level supervision. Rewarding sound logic rather than fortunate guesses helps the model build habits that travel to new problems.

\textbf{Trajectory-Stable Optimization:} \citet{li2025stppo} identified instabilities when PPO encounters multi-turn agent trajectories with delayed rewards. Their ST-PPO fixes push training completion from 34\% (vanilla PPO) to 89\%, making long-horizon agentic fine-tuning actually work.

\begin{figure}[htbp]
\centering
\begin{tikzpicture}[scale=0.9, transform shape]
    % === LEFT SIDE: Outcome Supervision ===
    \node[font=\bfseries\small, fill=red!20, rounded corners, inner sep=6pt] at (-4.5, 4.5) {Outcome Supervision};

    % Reasoning chain (left)
    \foreach \i/\status in {1/gray!30, 2/gray!30, 3/gray!30, 4/gray!30, 5/green!40} {
        \pgfmathsetmacro{\ypos}{3.5 - (\i-1)*0.9}
        \node[draw, fill=\status, rounded corners, minimum width=2.2cm, minimum height=0.6cm, font=\scriptsize] (L\i) at (-4.5, \ypos) {Step \i};
    }

    % Arrows between steps (left)
    \foreach \i in {1,...,4} {
        \pgfmathtruncatemacro{\j}{\i+1}
        \draw[->, thick, gray!60] (L\i) -- (L\j);
    }

    % Final reward indicator
    \node[draw, thick, fill=green!60, circle, minimum size=0.8cm, font=\bfseries\small] at (-2.3, -0.1) {\checkmark};
    \draw[->, thick, green!60!black] (L5.east) -- (-2.7, -0.1);
    \node[font=\tiny, align=center] at (-2.3, -0.7) {Final\\Reward};

    % Problem annotation
    \node[draw, dashed, red!60, rounded corners, text width=2.5cm, align=center, font=\tiny, inner sep=3pt] at (-4.5, -1.5) {Only final answer evaluated\\Flawed reasoning can pass};

    % === RIGHT SIDE: Process Supervision ===
    \node[font=\bfseries\small, fill=green!20, rounded corners, inner sep=6pt] at (4.5, 4.5) {Process Supervision};

    % Reasoning chain with step rewards (right)
    \foreach \i/\reward in {1/+0.2, 2/+0.3, 3/-0.1, 4/+0.4, 5/+0.5} {
        \pgfmathsetmacro{\ypos}{3.5 - (\i-1)*0.9}
        \node[draw, fill=blue!25, rounded corners, minimum width=2.2cm, minimum height=0.6cm, font=\scriptsize] (R\i) at (4.5, \ypos) {Step \i};
        % Reward indicator
        \node[fill=white, draw, rounded corners, font=\tiny\bfseries, inner sep=2pt] at (6.3, \ypos) {\reward};
        \draw[->, thick, blue!50] (R\i.east) -- (5.9, \ypos);
    }

    % Arrows between steps (right)
    \foreach \i in {1,...,4} {
        \pgfmathtruncatemacro{\j}{\i+1}
        \draw[->, thick, blue!60] (R\i) -- (R\j);
    }

    % Benefit annotation
    \node[draw, dashed, green!60!black, rounded corners, text width=2.8cm, align=center, font=\tiny, inner sep=3pt] at (4.5, -1.5) {Each step verified\\+12\% accuracy (MATH)\\Transferable reasoning};

    % === CENTER: Comparison ===
    \draw[thick, gray!40] (0, 4.8) -- (0, -2);
    \node[font=\scriptsize, gray, rotate=90] at (0.3, 1.5) {vs};

    % Key insight box at bottom
    \node[draw, thick, fill=yellow!20, rounded corners, font=\scriptsize, align=center, text width=10cm, inner sep=5pt] at (0, -2.8) {
        \textbf{Key Trade-off:} Process supervision requires step-level annotations (expensive) but yields more robust reasoning. Outcome supervision is cheaper but may reward ``lucky'' correct answers through flawed logic.
    };
\end{tikzpicture}
\caption{Outcome supervision vs process supervision for reasoning tasks. \textbf{Left:} Outcome-only rewards evaluate the final answer, potentially rewarding flawed reasoning that happens to land on the right conclusion. \textbf{Right:} Process supervision gives step-level feedback, validating every step of the reasoning process.
 Research has shown that step-level reward models typically exceed the performance of outcome-based models on MATH, with the gap widening as sampling budgets increase---sound reasoning chains generalize better to new problems.
}
\label{fig:process_supervision}
\end{figure}

\begin{tcolorbox}[colback=blue!5!white, colframe=blue!75!black, title=\textbf{Chapter 2 Key Takeaways}]
\begin{enumerate}
    \item \textbf{GRPO-style methods are common:} memory efficiency is a major practical driver of algorithm choice.
    \item \textbf{No single best algorithm:} PPO excels at exploration, DPO at preference alignment, GRPO at reasoning tasks.
    \item \textbf{Hybrid pipelines win:} Production systems typically combine SFT $\rightarrow$ DPO $\rightarrow$ GRPO.
    \item \textbf{Process supervision outperforms outcome supervision:} Step-level rewards improve generalization.
\end{enumerate}
\end{tcolorbox}

\chapter{Tool Use and Environment Interaction}
\label{ch:tooluse}

\textit{This chapter examines how agents engage with outside utilities and various environments.
 We cover API integration and tool learning (\S3.1), vision-language GUI automation (\S3.2), evaluation benchmarks and safety (\S3.3), and scaling for agents (\S3.4). Key insight: tool mastery is now essential for production deployment.}

\vspace{0.3cm}

\noindent\textit{See Table~\ref{tab:notation} for notation conventions.}

\section{Giving Models Hands}
Fundamentally, an LLM is a next-token predictor located in a datacenter. No eyes, no hands, no network connection. If you want it to actually \textit{do} things---book a flight, run a test suite, query a database---you need to give it interfaces: ways to convert sequences of tokens into API calls and button clicks.
 Think of it like wiring up nerves between a brain and muscles.

ReAct \cite{yao2023react} popularized a clean way to do this: let the model alternate between thinking out loud (``I should search for the user's email'') and taking action (``Search[user email]''). The thinking part keeps the reasoning interpretable; the action part keeps it grounded in what actually happens. On ALFWorld household tasks, this interleaved approach beat action-only baselines by 34\%.

\begin{figure}[htbp]
\centering
\begin{tikzpicture}[
    box/.style={rectangle, draw, fill=blue!20, text width=2.5cm, text centered, rounded corners, minimum height=1cm},
    arrow/.style={->, >=stealth, thick}
]
    % Nodes
    \node[box] (llm) {LLM Core};
    \node[box, below=1.5cm of llm] (reasoning) {Reasoning\\{\scriptsize(CoT/ReAct)}};
    \node[box, right=2.5cm of reasoning] (tools) {Tool Use\\{\scriptsize(APIs, Code)}};
    \node[box, below=1.5cm of reasoning] (env) {Environment\\{\scriptsize(Web, OS)}};
    \node[box, left=2.5cm of reasoning] (memory) {Memory\\{\scriptsize(RAG)}};

    % Arrows
    \draw[arrow] (llm) -- (reasoning);
    \draw[arrow] (reasoning) -- (tools);
    \draw[arrow] (tools.south) -- ++(0,-0.5) -| (env.east);
    \draw[arrow] (env.north) -- node[right, font=\scriptsize] {Obs} (reasoning.south);
    \draw[arrow] (memory) -- (reasoning);
    \draw[arrow] (reasoning) -- node[above, font=\scriptsize] {Store} (memory);
    \draw[arrow, dashed] (env.west) -- ++(-1,0) |- node[near start, left, font=\scriptsize] {Feedback} (llm.west);
\end{tikzpicture}
\caption{How agentic workflows connect LLM-based logic, tool integration, environment feedback, and memory.
 The agent sharpens its actions step by step based on what it observes.}
\label{fig:agentic_workflow}
\source{Adapted from Yao et al.~\cite{yao2023react}}
\end{figure}

\textbf{Toolformer} \cite{schick2023toolformer} figured out how to teach tool use without labeled examples. The trick: filter training data to keep only cases where calling a tool (calculator, calendar, search) actually helped predict the next token better. No explicit ``use the calculator here'' annotations---the model just learns that certain situations benefit from external calls.

\subsection{Different Approaches to Teaching Tool Use}
\label{sec:tool_learning}

Researchers have tried several angles:

\textbf{Learn from execution traces:} \citet{qin2023toolllm} grabbed 16,000+ real API specs and recorded what happened when agents used them correctly. Train on those traces and you get models that score 85.1\% on multi-API tasks---proof that tool skills transfer when you show enough examples of successful use.

\textbf{Look stuff up at runtime:} \citet{patil2023gorilla} attacked the hallucinated-API-call problem by retrieving actual documentation during inference. Ground your output in real specs and signature accuracy jumps to 94.2\%. In production, where a wrong parameter silently corrupts data, that matters.

\textbf{Route to specialists:} \citet{shen2023hugginggpt} and \citet{liang2023taskmatrix} utilize the model as a central dispatcher
 that breaks requests into pieces and farms them out to specialized models (image generation, classification, audio).
 The LLM does the reasoning; dedicated models do the heavy lifting. Scores 72.3\% on multimodal benchmarks.

\textbf{Discover skills through play:} \citet{wang2023voyager} let a Minecraft agent explore freely and bank whatever skills it discovered. Result: it finds diamond tools 15.3$\times$ faster than baselines by reusing what it already learned. Early taste of lifelong learning.

\subsection{What Shipped in 2025}
\label{sec:tool_learning_2025}

The last year brought a wave of new approaches:

\textbf{DeepAgent} \cite{li2025deepagent}: Most systems choke beyond 20 steps because they try to process the entire execution history. DeepAgent sidesteps this with autonomous ``memory folding''---compressing past interactions into structured episodic, working, and tool memories. It is trained with ToolPO, an end-to-end RL method that uses LLM-simulated APIs and assigns credit to tool-invocation tokens.

\textbf{Tool-Star}: Two-phase training: first imitate demos that include explicit ``thinking'' tokens, then tune with RL on outcomes. Those thinking tokens give you step-level signal for free, even when you only check the final answer.

\textbf{ReTool} \cite{feng2025retool}: Built-in reflection loop. Agent runs a tool, looks at the result, and decides: accept, retry with different parameters, or switch tools entirely. This self-correction hits \textbf{67\% on AIME 2024}---competitive with much larger models.

\textbf{VerlTool} \cite{jiang2025verltool}: An agentic RL training framework that extends veRL to multi-turn tool use (code execution, search, SQL, visual and web tools) with asynchronous rollouts; it is a training infrastructure, not a grammar-constrained decoder.

\textbf{ToolBrain} \cite{le2025toolbrain}: When you have 1000+ possible APIs, in-context examples don't scale. ToolBrain learns embeddings for tool capabilities and picks the right one via retrieval. 89\% precision, sub-100ms latency.

\textbf{AGENTFLOW} \cite{chen2025agentflow}: Parallelizes multi-tool workflows by running subtasks concurrently and coordinating through a message bus. 3.2$\times$ speedup over sequential execution, no accuracy hit.

\section{When Text Isn't Enough: Vision-Language Agents}
\label{sec:vision_agents}

HTML will tell you there's a button labeled ``Submit,'' but it won't tell you the button is crammed into the bottom-right corner behind a cookie banner. Screen interaction is fundamentally spatial---you need to know \textit{where} things are, not just what they're called. Pure text models are working with one hand tied behind their back.

The core challenge: translate ``click the search bar'' into an (x, y) coordinate that actually lands on the right pixel. That's visual grounding, and it requires architectural tricks beyond slapping a vision encoder onto GPT:

\textbf{Dual-pathway vision processing.} \citet{hong2023cogagent} built CogAgent with two separate pathways---one for fine-grained pixel details, one for high-level semantics. Result: 96.2\% accuracy at localizing UI elements, even on cluttered screens where the ``X'' button to close a popup is 12 pixels wide.

\textbf{Autonomous app exploration.} \citet{yang2023appagent} skipped the demonstration dataset entirely. Their agent explores apps autonomously, builds up a mental map of what buttons do, and figures out workflows on its own. Hit 67.8\% success on real smartphone tasks---not bad for an agent that was never shown the app before.

\textbf{Grounding at scale.} \citet{cheng2024seeclick} released 1.2 million labeled GUI elements. Train on that and you get 89.4\% accuracy mapping natural language to screen coordinates. Without data like this, models for visual grounding are flying blind.

% === VISION-LANGUAGE ATTENTION MECHANISM DIAGRAM ===
\begin{figure}[h]
\centering
\begin{tikzpicture}[scale=0.8, transform shape]
    % === BACKGROUND REGIONS ===
    \fill[blue!8, rounded corners] (-6.2, -0.5) rectangle (-2.8, 5.5);
    \fill[green!8, rounded corners] (2.8, -0.5) rectangle (6.2, 5.5);
    \fill[purple!10, rounded corners] (-1.8, 1.5) rectangle (1.8, 4.5);

    % === LEFT: VISUAL ENCODER ===
    \node[font=\bfseries\small, blue!70!black] at (-4.5, 5.2) {Visual Encoder};

    % Image representation
    \node[draw, thick, fill=blue!20, minimum width=2.5cm, minimum height=1.8cm, rounded corners]
        (image) at (-4.5, 3.8) {};
    \draw[step=0.4, blue!40, thin] (-5.7, 2.95) grid (-3.3, 4.65);
    \node[font=\tiny, blue!70!black] at (-4.5, 3.0) {Image Patches};

    % Patch embeddings
    \foreach \i/\y in {1/2.0, 2/1.4, 3/0.8, 4/0.2} {
        \node[draw, fill=blue!30, minimum width=1.8cm, minimum height=0.4cm, rounded corners=2pt, font=\tiny]
            (vpatch\i) at (-4.5, \y) {$\mathbf{v}_{\i}$};
    }

    % Arrow from image to patches
    \draw[->, thick, blue!60] (-4.5, 2.9) -- (-4.5, 2.3);

    % === RIGHT: TEXT ENCODER ===
    \node[font=\bfseries\small, green!60!black] at (4.5, 5.2) {Text Encoder};

    % Text tokens
    \node[draw, thick, fill=green!20, minimum width=2.5cm, minimum height=1.2cm, rounded corners, font=\scriptsize, align=center]
        (text) at (4.5, 4.2) {"Click the\\search button"};

    % Token embeddings
    \foreach \i/\y/\tok in {1/2.0/click, 2/1.4/the, 3/0.8/search, 4/0.2/button} {
        \node[draw, fill=green!30, minimum width=1.8cm, minimum height=0.4cm, rounded corners=2pt, font=\tiny]
            (tpatch\i) at (4.5, \y) {$\mathbf{t}_{\i}$: \tok};
    }

    % Arrow from text to tokens
    \draw[->, thick, green!60] (4.5, 3.5) -- (4.5, 2.3);

    % === CENTER: CROSS-ATTENTION BLOCK ===
    \node[font=\bfseries\small, purple!70!black] at (0, 4.2) {Cross-Attention};

    % Attention computation
    \node[draw, thick, fill=purple!25, minimum width=3cm, minimum height=1.8cm, rounded corners, align=center, font=\scriptsize]
        (attn) at (0, 2.8) {
            $\text{Attention}(Q, K, V)$\\[2pt]
            $= \text{softmax}\left(\frac{QK^T}{\sqrt{d}}\right)V$
        };

    % Query, Key, Value labels
    \node[draw, fill=green!40, rounded corners=3pt, font=\tiny, minimum width=0.8cm] (Q) at (-0.9, 1.7) {Q};
    \node[draw, fill=blue!40, rounded corners=3pt, font=\tiny, minimum width=0.8cm] (K) at (0, 1.7) {K};
    \node[draw, fill=blue!40, rounded corners=3pt, font=\tiny, minimum width=0.8cm] (V) at (0.9, 1.7) {V};

    % === ARROWS: Visual to Cross-Attention ===
    \draw[->, thick, blue!60, dashed] (-3.5, 1.4) -- (-1.2, 2.5);
    \draw[->, thick, blue!60, dashed] (-3.5, 0.8) -- (-1.2, 2.2);
    \node[font=\tiny, blue!60!black, rotate=20] at (-2.5, 2.4) {K, V};

    % === ARROWS: Text to Cross-Attention ===
    \draw[->, thick, green!60, dashed] (3.5, 1.4) -- (1.2, 2.5);
    \draw[->, thick, green!60, dashed] (3.5, 0.8) -- (1.2, 2.2);
    \node[font=\tiny, green!60!black, rotate=-20] at (2.5, 2.4) {Q};

    % === OUTPUT ===
    \node[draw, thick, fill=orange!30, minimum width=3.5cm, minimum height=0.6cm, rounded corners, font=\small]
        (output) at (0, 0) {\textbf{Grounded Action}};

    \draw[->, thick, purple!70] (0, 1.5) -- (0, 0.4);

    % Output description
    \node[font=\tiny, align=center] at (0, -0.6) {$(x, y)$ coordinates + action type};

    % === ANNOTATION ===
    \node[draw, thick, fill=yellow!20, rounded corners, font=\scriptsize, align=left, text width=5.5cm]
        at (0, -1.8) {
            \textbf{Key Insight:} Cross-attention allows text queries\\
            to ``attend'' to relevant visual regions, enabling\\
            precise grounding of natural language to pixels.
        };

\end{tikzpicture}
\caption{Cross-attention mechanism in vision-language agents. \textbf{Left:} The visual module converts image patches into embeddings $\mathbf{v}_i$. \textbf{Right:} The text module converts tokens into embeddings $\mathbf{t}_i$.
 \textbf{Center:} Cross-attention uses text as queries (Q) and visual features as keys/values (K, V), enabling the model to locate the visual region corresponding to ``search button.'' This architecture underlies GUI automation systems like CogAgent~\cite{hong2023cogagent} and SeeClick~\cite{cheng2024seeclick}, achieving 89--96\% grounding accuracy.}
\label{fig:cross-attention-mechanism}
\end{figure}

\section{Measuring What Matters: Benchmarks and Safety}
\label{sec:eval_safety}

How do you test whether an agent can actually do useful work? You throw it at problems that can't be gamed with pattern matching. Three benchmarks emerged as the standard gauntlet because each one breaks agents in different ways:

\textbf{SWE-bench} \cite{jimenez2024swebench} hands the agent a real GitHub issue from a real repo---cryptic bug report, thousands of files, tangled dependencies---and asks for a patch that passes the test suite. GPT-4 scored 1.7\% out of the gate. Turns out generating a function is easy; navigating a codebase to find \textit{which} function to change is hard.

\textbf{WebArena} \cite{zhou2024webarena} sets up realistic websites---e-commerce stores, content management systems---and assigns tasks like ``book the cheapest flight to Denver'' or ``cancel the pending order.'' Humans hit 78.2\%. Agents? 14.4\%. The gap isn't about language understanding; it's about clicking the right buttons in the right order when the UI doesn't cooperate.

\textbf{OSWorld} \cite{xie2024osworld} goes deeper: full desktop control. Move files, open applications, change system settings. The agent sees screenshots and outputs mouse/keyboard actions. This tests vision, planning, and motor precision all at once---the closest thing we have to measuring ``can this agent do my job.''

\begin{figure}[htbp]
\centering
\begin{tikzpicture}
    \begin{axis}[
        ybar,
        ylabel={Success Rate (\%)},
        symbolic x coords={SWE-bench, WebArena, OSWorld},
        xtick=data,
        width=12cm,
        height=6cm,
        ymin=0,
        ymax=100,
        legend style={at={(0.5,-0.15)}, anchor=north, legend columns=-1},
        bar width=15pt,
    ]
    \addplot coordinates {(SWE-bench,1.7) (WebArena,14.4) (OSWorld,12.2)};
    \addplot coordinates {(WebArena,78.2) (OSWorld,72.4)};
    \legend{GPT-4, Human}
    \end{axis}
\end{tikzpicture}
\caption{Performance gap between GPT-4 and humans on agentic benchmarks. \textbf{SWE-bench} tests code repair on real GitHub issues; \textbf{WebArena} checks web navigation (booking, shopping); \textbf{OSWorld} evaluates desktop automation requiring vision and control. SWE-bench has no published human baseline, so no human bar is shown. The large gap (GPT-4/GPT-4V at 1.7--14.4\% vs.\ humans at 72--78\% where measured) shows how hard autonomous multi-step execution is and motivates RL-based fine-tuning.}
\label{fig:benchmark_comparison}
\source{Data from ~\cite{jimenez2024swebench,zhou2024webarena,xie2024osworld}}
\end{figure}

Table~\ref{tab:sota_benchmarks} rounds up current top scores on major agentic benchmarks, tracing the jump from prompting-only methods to RL-fine-tuned agents.

\begin{table}[h]
\centering
\caption{Reported results on agentic and reasoning benchmarks (as of December 2025). Rows are as reported by the cited sources; values from different models, splits, or scaffolds are not converted into improvement ratios. Unattributed historical rows from earlier drafts are removed.}
\label{tab:sota_benchmarks}
\small
\begin{tabular}{@{}llccc@{}}
\toprule
\textbf{Benchmark} & \textbf{Model/Method} & \textbf{Success (\%)} & \textbf{Training} & \textbf{Year} \\
\midrule
SWE-bench (original) & GPT-4, retrieval baseline~\cite{jimenez2024swebench} & 1.7 & Prompting & 2023 \\
SWE-bench Verified & Claude-3.5 Sonnet with agent scaffold (vendor-reported) & 49.0 & Proprietary & 2024 \\
\midrule
\multirow{2}{*}{WebArena-Lite} & Llama-3.1-8B before / after WebAgent-R1 & 8.5 / 44.8 & End-to-end RL & 2025 \\
 & Qwen-2.5-3B before / after WebAgent-R1 & 6.1 / 33.9 & End-to-end RL & 2025 \\
\midrule
\multirow{3}{*}{AIME 2024} & GPT-4o~\cite{deepseek2025r1} & 9.3 & Prompting & 2024 \\
 & ReTool~\cite{feng2025retool} & 67.0 & SFT+Tools+RL & 2025 \\
 & DeepSeek-R1~\cite{deepseek2025r1} & 79.8 & GRPO & 2025 \\
\midrule
\multirow{3}{*}{HumanEval} & GPT-4 & 67.0 & Prompting & 2023 \\
 & MetaGPT~\cite{hong2023metagpt} & 85.9 & Role-based & 2023 \\
 & DeepSeek-Coder-V2 & 90.2 & SFT+RL & 2024 \\
\midrule
\multirow{3}{*}{OSWorld} & Best model in original paper~\cite{xie2024osworld} & 12.2 & Prompting & 2024 \\
 & Human~\cite{xie2024osworld} & 72.4 & -- & 2024 \\
 & Mobile-Agent-v3 & 37.7 & Multi-agent framework & 2025 \\
\midrule
GAIA & DeepAgent-32B~\cite{li2025deepagent} & 53.3 & ToolPO RL & 2025 \\
\bottomrule
\end{tabular}

\vspace{0.2cm}
\footnotesize
\textbf{Statistical notes:} Values are point estimates from the cited sources, which mostly do not report intervals. For a binomial proportion the 95\% half-width is $1.96\sqrt{\hat{p}(1-\hat{p})/n}$: about $\pm4.4$ points for SWE-bench Verified ($\hat{p}=0.49$, $n=500$) and about $\pm14$ points for AIME 2024 ($\hat{p}=0.80$, $n=30$, single pass). These are back-of-envelope sampling intervals, not published statistics.

\source{Benchmark sources: SWE-bench~\cite{jimenez2024swebench}, WebArena~\cite{zhou2024webarena}, OSWorld~\cite{xie2024osworld}, GAIA~\cite{mialon2023gaia}. WebArena-Lite rows: WebAgent-R1 (arXiv:2505.16421); Mobile-Agent-v3 OSWorld: arXiv:2508.15144 (its 73.3\% figure is AndroidWorld, not OSWorld).}
\end{table}



\textbf{Reading the numbers.} Reported results have improved substantially, although the commonly quoted pairs mix models, benchmark splits, and scaffolds: GPT-4 solved 1.7\% of the original SWE-bench~\cite{jimenez2024swebench}, while an agent scaffold around Claude~3.5 Sonnet is reported at 49\% on the 500-task SWE-bench Verified subset (a different split, so no ratio is computed); end-to-end RL raised Llama-3.1-8B on WebArena-Lite from 8.5\% to 44.8\% (WebAgent-R1); on OSWorld the best model in the original paper reached 12.2\% versus 72.4\% for humans~\cite{xie2024osworld}, and Mobile-Agent-v3 reports 37.7\% in 2025; and on AIME 2024, GPT-4o scores 9.3\% versus 79.8\% for DeepSeek-R1~\cite{deepseek2025r1}. Only some of these pairs isolate RL: WebAgent-R1 (same model before and after end-to-end RL) and DeepSeek-R1 (GRPO) do; the SWE-bench pair reflects a different model, a different split, and an inference-time scaffold.


\subsection{The Safety Problem Gets Worse with Agents}

Here's an uncomfortable truth: the same reasoning that lets an agent fix your code also lets it exploit a vulnerability. Give a model shell access and suddenly ``be helpful'' can conflict with ``don't break things.'' Safety can't be bolted on at the end---it needs to be baked into the reward function from day one.

How do we actually measure whether agents are reliable? Several benchmarks go beyond simple task completion:

\textbf{AgentBench} \cite{liu2023agentbench} evaluates models across eight different environments---terminals, databases, games, web browsers. GPT-4 completes 38.9\% of tasks; open models hover around 15--20\%. Neither is production-ready.

\textbf{GAIA} \cite{mialon2023gaia} curated 466 tasks that look like actual user requests---stuff that requires chaining tools and reasoning across multiple steps. GPT-4 with tools? 30.2\%. Apparently being ``helpful'' is harder than achieving high scores in traditional academic evaluations.

\textbf{TrustLLM} \cite{sun2024trustllm} sliced trustworthiness into measurable components: factuality, harm avoidance, fairness, robustness to adversarial inputs. The awkward finding: pushing safety from 72\% to 89\% sometimes tanks factuality from 91\% to 87\%. There's no free lunch.

\textbf{OS-level risk testing} \cite{yin2024os} found that giving models environmental access amplifies harm potential approximately 3.2$\times$ compared to text-only interaction, though further replication over various model architectures is needed to establish confidence bounds.
 Their 150-task suite, built on OSWorld and spanning deliberate misuse, prompt injection, and model misbehavior, covers everything from ``oops I deleted the wrong file'' to ``oops I leaked credentials.''

\subsubsection{Taxonomy of Agent Safety Risks}

Drawing on the surveyed literature, we identify four categories of safety concerns specific to agentic systems:

\begin{enumerate}
    \item \textbf{Action irreversibility:} Unlike chatbots, agents take real-world actions. Deleting files, sending emails, executing code---these cannot be undone with ``Sorry, I made a mistake.'' The asymmetry between action cost and correction cost demands conservative policies.

        \item \textbf{Goal misspecification:} the potential for reward shortcuts increases as models become more capable. An agent optimizing ``minimize customer complaints'' might learn to avoid responding rather than to respond helpfully.
 Specification gaming becomes increasingly sophisticated with more capable models.

    \item \textbf{Multi-step error propagation:} Errors compound across tool chains. A 95\% reliable tool used 10 times in sequence yields $(0.95)^{10} \approx 60\%$ end-to-end reliability. Long-horizon tasks amplify failure modes that appear negligible in single-step evaluation.

    \item \textbf{Emergent coordination risks:} Multi-agent systems can develop implicit communication protocols that evade human oversight. \citet{kenton2021alignment} term this ``mesa-optimization''---inner objectives emerging from outer training signals.
\end{enumerate}

\subsubsection{International Safety Research Landscape}

Agent safety research spans global institutions with complementary strengths:

\textbf{Chinese contributions:} DeepSeek's work on GRPO \cite{deepseek2025r1} includes safety-aware reward shaping, penalizing harmful action sequences during training. Tsinghua University's SafeAgent benchmark evaluates refusal behavior across 12 harm categories. ByteDance's DAPO \cite{bytedance2025dapo} incorporates dynamic safety thresholds that tighten as training progresses.

\textbf{European contributions:} DeepMind (London) contributes scalable oversight and evaluation research, including reward modeling studies and agent evaluation frameworks. EPFL's work on formal verification of neural policies offers provable safety bounds for constrained action spaces. The EU AI Act's risk-based framework influences benchmark design for high-stakes applications.

\textbf{North American contributions:} Anthropic pioneered Constitutional AI \cite{bai2022constitutional}, enabling self-supervised safety alignment now widely adopted. Their research into scalable oversight mechanisms addresses the ``who watches the watchmen'' problem. OpenAI's red-teaming methodology establishes industry standards for pre-deployment safety evaluation. Stanford's HELM benchmark includes safety sub-metrics alongside capability measures.

This geographic diversity strengthens the field: different regulatory contexts surface different failure modes, and cross-cultural evaluation catches biases that single-region testing misses.

\section{Scaling Laws for Agents}
Recent work shows that ``test-time compute''---giving models extra cycles to spin up more reasoning tokens---often beats just bulking up the model itself. \citet{snell2024scaling} demonstrate that optimally scaling test-time compute can rival much larger models on reasoning tasks.

\subsection{Limitations of Current Tool-Use and Evaluation Research}
This survey focuses on English-language publications from 2023--2025; non-English literature and earlier foundational work may be underrepresented. Benchmark performance metrics reported here reflect best-case results from published papers and may not generalize to deployment conditions. Tool-use evaluations are predominantly conducted on curated benchmarks with known APIs---real-world tool ecosystems involve undocumented APIs, version drift, and authentication complexity not captured in current evaluations. Additionally, safety benchmarks remain limited in scope, testing specific attack vectors rather than comprehensive threat models.

\begin{tcolorbox}[colback=green!5!white, colframe=green!60!black, title=\textbf{Chapter 3 Key Takeaways}]
\begin{enumerate}
    \item \textbf{Tools transform capabilities:} ReAct-style interleaving of reasoning and action beats pure reasoning by 34\% on household tasks.
    \item \textbf{Code agents advanced quickly:} 1.7\% (GPT-4, original SWE-bench, 2023) and 49\% (Claude-3.5 Sonnet scaffold, SWE-bench Verified, 2024) are on different splits, so the gap is not a like-for-like multiplier.
    \item \textbf{GUI automation lags:} on OSWorld, the human reference in Table~\ref{tab:sota_benchmarks} is 72.4\%; agent results remain well below human level on the original OSWorld task set, so pixel-level control remains harder than API calls.
    \item \textbf{Safety scales with capability:} Environmental access amplifies harm potential $\sim$3.2$\times$.
\end{enumerate}
\end{tcolorbox}

\chapter{Beyond the Basics: RAG and Multi-Agent Systems}
\label{ch:advancedarch}

\textit{This chapter explores two architectural patterns that extend agent capabilities beyond single-model reasoning. Section~\ref{sec:rag-matters} explains why retrieval-augmented generation (RAG) solves knowledge staleness and context limitations. Section~4.2 covers multi-agent systems where specialized agents work together on intricate problems. The main insight: these patterns address fundamental limitations of monolithic models.
}

\vspace{0.3cm}
\noindent\textit{See Table~\ref{tab:notation} for notation conventions.}

\section{Why RAG Matters for Agents}
\label{sec:rag-matters}
All knowledge in a neural network is stored in its weights, which are fixed at the end of training. The world keeps changing; the model doesn't. Want to update its knowledge? Retrain it---which burns GPU weeks and calendar days.

RAG sidesteps this by splitting the problem in two. The model handles reasoning. A separate database handles facts. Update the database and the model's answers update instantly---no gradient descent required. This isn't a hack; it's a design principle. Reasoning evolves slowly (training is expensive), but facts evolve constantly (database writes are cheap).

There's another constraint: context windows have hard limits. Llama-3 caps at 128K tokens. GPT-4 Turbo supports 128K tokens. Claude-3.5 reaches 200K tokens. Sounds like a lot until you're debugging a million-line codebase or synthesizing a month's worth of communications. Even with expanded windows, pure parametric models hit a wall because attention scales quadratically and extended context lengths degrade retrieval accuracy.

RAG solves this too: instead of stuffing everything into context, you retrieve what's relevant when you need it. The model sees the question plus the top-K passages most likely to help---nothing more, nothing less.

\begin{figure}[htbp]
\centering
\begin{tikzpicture}[
    proc/.style={rectangle, draw, fill=cyan!20, text width=2.5cm, text centered, minimum height=0.9cm, rounded corners},
    db/.style={cylinder, draw, fill=gray!20, text width=1.8cm, text centered, minimum height=1.2cm, aspect=0.25},
    arrow/.style={->, >=stealth, thick}
]
    \node[proc] (query) {User Query};
    \node[proc, below=1.2cm of query] (retrieve) {Retrieval\\{\scriptsize(BM25 + Dense)}};
    \node[db, right=2cm of retrieve] (kb) {Knowledge\\Base};
    \node[proc, below=1.2cm of retrieve] (context) {Context\\Assembly};
    \node[proc, below=1.2cm of context] (llm) {LLM Generation};
    \node[proc, below=1.2cm of llm] (response) {Final Response};

    \draw[arrow] (query) -- (retrieve);
    \draw[arrow] (retrieve) -- (kb);
    \draw[arrow] (kb) -- (retrieve);
    \draw[arrow] (retrieve) -- (context);
    \draw[arrow] (context) -- (llm);
    \draw[arrow] (llm) -- (response);

    \node[right=1.5cm of context, text width=2.5cm, font=\scriptsize, align=left] {Top-K\\relevant chunks};
\end{tikzpicture}
\caption{How Retrieval-Augmented Generation grounds knowledge. Given a query, the system pulls Top-K relevant chunks from a knowledge base using hybrid search (BM25 + dense embeddings). Retrieved context merges with the query and feeds into the language model. This curbs hallucination and lets the system tap into current facts beyond its training cutoff, boosting factual accuracy 15--30\% on knowledge-heavy tasks.}
\label{fig:rag_architecture}
\source{Adapted from \citet{gao2023retrieval}}
\end{figure}

\textbf{HippoRAG} \cite{gutierrez2024hipporag} draws from hippocampal memory, separating working memory from long-term storage. Splitting memory this way sharpens multi-hop reasoning by fetching far-off facts without jamming up the active context. Combining such memory systems with RL-tuned policies promises "open-book" agents that consult external knowledge appropriately.

\subsection{Memory Architectures and Knowledge Integration}
\label{sec:memory}

A handful of ideas tackle the memory bottleneck:

\textbf{Operating-System-Style Paging:} \citet{packer2023memgpt} implemented virtual memory for language models---swapping context between active attention and persistent storage. This keeps 10,000-turn dialogues humming with just 1.2\% performance drop---proof that OS-style memory tricks port straight to language models.

\textbf{Modular Retrieval Integration:} \citet{shi2023replug} treated the retriever as interchangeable, boosting zero-shot QA from 41.5\% to 56.8\% on Natural Questions. This modular design lets practitioners swap retrievers without having to retrain the entire model from scratch.

\textbf{End-to-End Retrieval Learning:} \citet{izacard2022atlas} jointly optimized retriever and generator, proving that 11B parameters plus learned retrieval equals 175B-parameter performance (64.0\% on Natural Questions) without it---a solid argument for offloading knowledge.

\subsection{Advanced RAG Systems for Agents (2025)}
\label{sec:rag_2025}

Newer papers have tailored RAG for agentic workflows, often baking in RL-based tuning:

\textbf{PoU (Proof-of-Use Verification)}: Addresses a critical RAG failure mode: retrieving relevant passages but failing to actually use them. PoU adds explicit verification that generated answers cite retrieved content, penalizing responses that ignore retrieved context.

\textbf{InfoAgent} \cite{zhang2025infoagent}: A RAG system trained with RL to minimize information retrieval latency while maintaining accuracy. Learns to predict which sources will contain answers, enabling parallel retrieval and 3$\times$ faster response times on information-seeking tasks.

\textbf{RAHL (Retrieval-Augmented Hierarchical Learning)}: Separates retrieval into two levels: (1) coarse-grained topic retrieval selecting document clusters, then (2) fine-grained passage retrieval within selected clusters. The hierarchical approach significantly reduces search space while maintaining high recall on targeted passages.

\begin{definition}[Retrieval-Augmented Policy]
\label{def:rap}
A retrieval-augmented policy $\pi_{\text{RAG}}$ is a composite function:
\begin{equation}
\pi_{\text{RAG}}(a | s, q) = \pi_\theta\left(a \Big| s, \text{Retrieve}_K\left(\mathcal{K}, \phi(q)\right)\right)
\end{equation}
where $\mathcal{K}$ is an external knowledge base, $\phi(\cdot)$ is an embedding function, $\text{Retrieve}_K$ returns the top-$K$ passages ranked by similarity score $\text{sim}(\phi(q), \phi(d))$ for documents $d \in \mathcal{K}$, and $\pi_\theta$ is the base policy conditioned on both state $s$ and retrieved context.
\end{definition}

\begin{toymodel}[When Retrieval Helps Factuality]
\label{prop:hallucination}
Let $c \in [0,1]$ be the probability the knowledge base contains sufficient evidence for a query, $\epsilon_r$ the probability the retriever fails to return that evidence when it exists, and $\epsilon_g$ the probability the generator answers incorrectly \emph{even when given correct evidence} (e.g., misreading or ignoring it). Under a simplified mixture model:
\begin{equation}
H(\pi_{\text{RAG}}) \leq (1 - c) \cdot H(\pi_{\text{base}}) + c \cdot \left(\epsilon_r + (1-\epsilon_r)\epsilon_g\right)
\end{equation}
where $H(\pi_{\text{base}})$ is the base model's hallucination rate without retrieval. The bound assumes (i) non-distraction: when the knowledge base lacks evidence, retrieved context does not raise the hallucination rate above $H(\pi_{\text{base}})$; and (ii) a pessimistic retrieval-failure case, in which every query whose evidence exists but is not retrieved is counted as a hallucination. Under (i) and (ii) it follows from the law of total probability over these three events. This is not a guarantee; it clarifies \textbf{which failure modes}---coverage ($c$), retrieval ($\epsilon_r$), and evidence-use ($\epsilon_g$)---drive outcomes.
\end{toymodel}

\begin{remark}
This simplified model helps interpret empirical RAG results. When retrieval works (high $c$, low $\epsilon_r$), the bottleneck shifts to evidence-use error $\epsilon_g$---the model must actually read and apply retrieved passages. Beefing up knowledge base coverage often yields the largest practical gains because $c$ appears linearly in the bound, but practitioners should also monitor whether models ignore retrieved context (high $\epsilon_g$).
\end{remark}

\section{Multi-Agent Systems: Divide and Conquer}
One model trying to do everything eventually chokes. Context fills up. Attention fragments. Error compounds. The fix? Split the work across specialists.

Think of it like a software team. You don't have one person who writes code, reviews it, tests it, deploys it, and monitors it. You have engineers, reviewers, QA, SREs---each focused on what they do best. Multi-agent systems apply the same logic: assign subtasks to agents that specialize in them. The manager plans. The coder implements. The critic checks. Errors get caught earlier. Parallelism speeds things up. And when one agent screws up, others can compensate.

\begin{figure}[htbp]
\centering
\begin{tikzpicture}[
    agent/.style={circle, draw, fill=purple!20, minimum size=1.3cm, font=\scriptsize},
    comm/.style={<->, >=stealth, dashed, thin},
    arrow/.style={->, >=stealth, thick}
]
    \node[agent] (a1) at (0,1.5) {Agent 1};
    \node[agent] (a2) at (3,1.5) {Agent 2};
    \node[agent] (a3) at (0,0) {Agent 3};
    \node[agent] (a4) at (3,0) {Agent 4};

    \node[rectangle, draw, fill=blue!10, text width=2.5cm, text centered, minimum height=0.8cm, rounded corners] (env) at (1.5,3.5) {Shared Environment};

    \draw[comm] (a1) -- (a2);
    \draw[comm] (a1) -- (a3);
    \draw[comm] (a2) -- (a4);
    \draw[comm] (a3) -- (a4);
    \draw[comm] (a1) -- (a4);
    \draw[comm] (a2) -- (a3);

    \draw[arrow] (a1.north) -- ++(0,0.8) -| (env.south west);
    \draw[arrow] (a2.north) -- ++(0,0.8) -| (env.south east);
    \draw[arrow] (a3.north) -- ++(0,1.5);
    \draw[arrow] (a4.north) -- ++(0,1.5);
\end{tikzpicture}
\caption{Decentralized multi-agent architecture showing independent agents (1-4) coordinating through message passing (dashed lines) while interacting with a shared environment. Each agent operates independently while coordinating through messages. All agents interact with a shared environment, so multiple workers can run side-by-side and specialize (e.g., Agent 1: planning, Agent 2: coding, Agent 3: testing). Spreading out the control logic avoids bottlenecks and lets coordination patterns emerge organically.}
\label{fig:multi_agent}
\source{Adapted from Zhang et al.~\cite{zhang2024agentnet}}
\end{figure}

\begin{itemize}
    \item \textbf{FireAct} \cite{chen2023fireact}: A single-agent method rather than a multi-agent one; it fine-tunes smaller models on ReAct trajectories distilled from a larger model, and such tuned agents can then serve as components of multi-agent systems.
    \item \textbf{AgentNet}: Proposes decentralized evolutionary coordination, achieving better task completion by letting agents evolve interaction patterns \cite{zhang2024agentnet}.
    \item \textbf{MasRouter}: Learns to route tasks dynamically between agents to optimize cost and performance \cite{chen2024masrouter}.
\end{itemize}
Decentralized architectures offer robustness, avoiding single points of failure in centralized controller models.

\subsection{Multi-Agent Collaboration Frameworks}
\label{sec:multi_agent_collab}

No single agent can juggle a mental model of the entire codebase while simultaneously writing tests, catching bugs, and keeping docs in sync. Even if it could, the working memory overhead would guarantee mistakes. The payoff of multi-agent systems is that each agent can specialize---and when they mess up, their teammates can catch it.

Several coordination patterns have emerged:

\textbf{The org-chart approach.} \citet{hong2023metagpt} gave agents job titles: product manager, architect, engineer. Each follows role-specific protocols. The hierarchy scored 85.9\% on HumanEval and cut hallucination by 62\%. Turns out, forcing agents to ``stay in their lane'' improves both accuracy and coherence.

\textbf{End-to-end pipelines.} \citet{qian2023chatdev} built a full software lifecycle---design, code, test, document---run by a chain of agents handing off to each other. They completed 70.4\% of software tasks from spec to working code. That's not a demo; that's a pipeline.

\textbf{Emergent social behavior.} \citet{park2023generative} dropped 25 agents into a virtual town and let them figure things out. No scripting. The agents planned their days, formed relationships, gossiped. Human raters judged 87\% of behaviors as believable. Coordination emerged without anyone designing it.

\textbf{Self-organized teams.} \citet{li2023camel} had agents negotiate who does what through ``inception prompting''---agents defining each other's roles on the fly. Hit 76\% task completion with zero human intervention in task allocation.

\textbf{Plug-and-play topologies.} \citet{wu2023autogen} built a framework where you can wire agents together in different patterns---pipelines, trees, meshes---without rewriting everything. Swap the topology and 90\% of the code still works.

\subsection{RL-Trained Multi-Agent Coordination (2025 Advances)}
\label{sec:multi_agent_rl}

Recent work has moved beyond hand-designed role assignments to learning coordination through reinforcement learning:

\textbf{MATPO (Multi-Agent Trust-Region Policy Optimization)} \cite{mo2025matpo}: Extends PPO's trust region concept to multi-agent settings where agents must coordinate without centralized control. By maintaining separate advantage baselines per agent while sharing a joint reward signal, MATPO achieves \textbf{18.38\% improvement} over single-agent approaches on evaluations focused on collaborative reasoning.
 The key insight: agents learn complementary specializations---one focuses on information retrieval while another handles synthesis.

\textbf{M-GRPO (Multi-Agent GRPO):} Adapts GRPO's group normalization for hierarchical coordination. A manager agent learns to decompose tasks and assign them to worker agents, each with specialized capabilities. Advantages are normalized within each hierarchy level, preventing dominant agents from overshadowing specialists.

\textbf{Emergent Coordination Behaviors:} Training multi-agent systems with shared rewards produces emergent specialization without explicit role assignment. Documented behaviors include:
\begin{itemize}
    \item \textit{Skill differentiation}: Agents spontaneously specialize in retrieval vs. reasoning
    \item \textit{Communication protocols}: Agents develop shorthand for common subtasks
    \item \textit{Error recovery handoffs}: Agents learn to recognize when to delegate failed attempts
\end{itemize}

\begin{heuristic}[Communication Overhead by Topology]
\label{thm:multiagent}
For $n$ agents exchanging messages, communication costs per coordination round (i.e., one synchronization cycle where all agents send and receive messages) scale differently by topology:
\begin{enumerate}
    \item \textbf{Fully connected (decentralized):} $\mathcal{O}(n^2)$ pairwise messages per round---every agent talks to every other agent
    \item \textbf{Tree/hierarchical ($k$ levels):} $\mathcal{O}(n)$ messages per round and $\mathcal{O}(k)$ sequential routing hops, which is $\mathcal{O}(\log n)$ only for a balanced tree with branching factor at least 2---messages aggregate upward, decisions flow downward
    \item \textbf{Centralized hub:} $\mathcal{O}(n)$ messages per round---all agents communicate through a single coordinator, introducing a bottleneck and single-point-of-failure risk
\end{enumerate}
These are asymptotic communication-scaling heuristics, not end-to-end performance guarantees---constants, network latency, and message size significantly affect real-world overhead (e.g., fully connected with $n=3$ agents may outperform hierarchical due to coordination simplicity). The optimal architecture depends on task structure: decentralized topologies offer robustness for long-horizon tasks; hierarchical topologies enable clear decomposition; centralized topologies minimize latency when the coordinator is fast enough.
\end{heuristic}

\begin{remark}
This heuristic is \emph{consistent with} the empirical success of hierarchical multi-agent systems like MetaGPT (achieving 85.9\% on HumanEval \cite{hong2023metagpt}): tree-structured communication lets these systems tackle complex software dev tasks without the $\mathcal{O}(n^2)$ overhead of fully connected coordination. However, the heuristic does not \emph{explain} performance---many other factors (agent capability, task decomposition quality, prompt design) also matter.
\end{remark}

\section{What Actually Goes Inside an Agent?}
\label{sec:architectures}

Setting aside speculation, we ask: what architectural patterns actually work?

\textbf{Three design philosophies.} \citet{sumers2023cognitive} sorted architectures into buckets: symbolic (explicit rules), neural (end-to-end learning), and hybrid (rules + learning). Hybrids win by 15--30\% on long-horizon tasks. The neural part handles messy perception; the symbolic part keeps planning from going off the rails.

\textbf{Planning strategies compared.} \citet{wang2024planning} ran a bake-off: hierarchical decomposition (68\% success), external planner (72\%), iterative self-refinement (76\%). The pattern that kept winning: feed execution results back into planning. Doesn't matter what the agent is doing---closing the loop helps.

\textbf{What everyone actually builds.} \citet{xi2023rise} surveyed 100+ agent implementations. Every single one had an action module. 92\% had a planning module. Only 58\% had perception, and just 43\% had memory. Agents without memory struggle on anything requiring context across steps---which includes most interesting tasks.

\subsection{Limitations of RAG and Multi-Agent Research}
The RAG literature focuses predominantly on text retrieval; multimodal retrieval (images, code, structured data) remains underexplored. Retrieval quality metrics (e.g., recall@k) do not always correlate with downstream task performance, making evaluation methodology an open question. For multi-agent systems, coordination overhead and failure cascades are often underreported---papers typically highlight success cases rather than systematic failure analysis. Additionally, most multi-agent benchmarks test pre-defined role assignments; emergent role discovery and dynamic team composition remain largely unstudied. Communication efficiency claims should be interpreted cautiously, as real-world latency and bandwidth constraints differ substantially from simulated environments.

\begin{tcolorbox}[colback=orange!5!white, colframe=orange!60!black, title=\textbf{Chapter 4 Key Takeaways}]
\begin{enumerate}
    \item \textbf{RAG solves knowledge staleness:} Database updates propagate instantly vs. retraining (GPU-weeks).
    \item \textbf{Multi-agent beats monolithic:} Specialized agents + coordinator outperform single large models on complex tasks.
    \item \textbf{Memory is underutilized:} Only 43\% of implementations include memory---yet it's critical for multi-step tasks.
    \item \textbf{Hybrid architectures win:} Neural + symbolic combinations beat pure approaches by 15--30\%.
\end{enumerate}
\end{tcolorbox}

\chapter{Experimental Methodology and Results}
\label{ch:experiments}

This chapter has two parts. The first describes a proposed agentic evaluation programme (algorithm comparison, tool use, memory, multi-agent coordination, test-time scaling) that was \emph{not} executed. The second (Section~\ref{sec:actual_results}) reports the low-budget GRPO runs that were executed on GSM8K, using only runs with released W\&B records and held-out evaluation files, followed by a planned scaling study (Section~\ref{sec:systematic_study}).

\vspace{0.3cm}

\section{Experimental Design Philosophy}
\label{sec:design_philosophy}

\textbf{Scope note.} Sections~\ref{sec:design_philosophy}--\ref{sec:comprehensive_experiments} describe a \emph{proposed} agentic evaluation programme (WebArena, SWE-bench, OSWorld, PPO/DPO baselines, RAG). None of those experiments was executed for this report. The executed work is limited to the low-budget GRPO runs on GSM8K reported in Section~\ref{sec:actual_results}.

\textbf{Requirement Questioning:} Why run experiments at all? Without controlled comparison, we cannot distinguish algorithmic improvements from noise. Why three benchmarks? One is insufficient, but three orthogonal benchmarks (SWE-bench for code, WebArena for web, OSWorld for OS) would expose different failure modes efficiently. Why 8B models? Larger models are expensive for academic research.

\textbf{Simplified Design (proposed):} Three core experiments (PPO baseline, GRPO test, SFT control) on WebArena with Llama-3-8B would address the key question: does GRPO demonstrate higher sample efficiency than PPO for agentic tasks?

The proposed programme has three lines of investigation: (1) establish baselines to understand what 8B models can do without specialized training; (2) compare algorithms to understand which algorithmic biases matter for which task types; (3) evaluate RAG and tool use as solutions to specific problems---RAG for knowledge limits, tool use for action limits.

\subsection{Research Hypotheses}
From the papers we surveyed, we distill these testable hypotheses for the proposed programme; none is tested in this report. We define \textbf{sample efficiency} operationally as $I_{80}$: the number of environment interactions required to reach 80\% of final performance measured at a predetermined iteration count $N_{\max}$ (we would use $N_{\max} = 10{,}000$ interactions). Hypotheses avoid hard numeric thresholds where the literature does not support them:

\begin{description}
    \item[H1 (Sample-Efficiency Endpoint, GRPO vs.\ PPO):] On interactive agentic benchmarks (WebArena, optionally OSWorld), GRPO reaches $I_{80}$ using \textbf{fewer environment interactions} than PPO under matched compute and reward design. Primary endpoint: $I_{80}(\text{GRPO}) / I_{80}(\text{PPO}) < 1$.

    \item[H2 (Tool-Use Training Generalization):] Models fine-tuned on explicit tool-use trajectories achieve higher WebArena task success than few-shot baselines, and the improvement persists on held-out websites. Primary endpoints: mean TSR improvement and held-out TSR improvement.

    \item[H3 (RAG Conditional Benefit):] On tasks requiring external knowledge, RAG improves \textbf{evidence-supported correctness} when retrieval quality is high, \textbf{conditional on retrieval accuracy/coverage}; failure cases with wrong retrieval would be tested explicitly.

    \item[H4 (Algorithm $\times$ Task Interaction):] The relative advantage of PPO/DPO/GRPO depends on task type, tested as an algorithm $\times$ benchmark interaction with seeds as the replication unit.

    \item[H5 (Scaling Scope, Aligned to Resources):] If multiple model sizes are evaluated, RL fine-tuning gains show diminishing returns with size. If only $\sim$8B models are feasible, H5 is a motivation rather than a tested hypothesis.
\end{description}

\section{Ablation Framework (Proposed)}
Systematic ablation identifies which components matter. The proposed matrix would isolate algorithm choice (PPO vs DPO vs GRPO) from architecture choice (ReAct vs Toolformer) and model choice (Llama vs Qwen). It was not executed.

\begin{table}[h]
\centering
\caption{Proposed Ablation Configuration Matrix (not executed)}
\label{tab:ablation}
\begin{tabular}{lllll}
\toprule
\textbf{ID} & \textbf{Model} & \textbf{Algorithm} & \textbf{Tool Use} & \textbf{Benchmark} \\
\midrule
A1 & Llama-3-8B & PPO & ReAct & SWE-bench \\
A2 & Llama-3-8B & DPO & ReAct & SWE-bench \\
A3 & Llama-3-8B & GRPO & ReAct & SWE-bench \\
A4 & Qwen2.5-7B & DPO & Toolformer & WebArena \\
A5 & Qwen2.5-7B & STEP-DPO & Toolformer & OSWorld \\
\bottomrule
\end{tabular}
\end{table}

\section*{Proposed Experimental Suite (Not Executed)}
\label{sec:comprehensive_experiments}

This section lists five proposed experiments addressing the research directions identified in our survey. They are a design for future work; no results from them are reported.

\subsection*{Experiment 1: Algorithm Comparison (PPO vs DPO vs GRPO)}
\textbf{Core question:} Which RL algorithm achieves the best sample efficiency and final performance on agentic tasks? \textbf{Design:} PPO with a learned reward model; DPO with preference pairs; GRPO with group-relative advantages; SFT-only baseline. \textbf{Benchmarks:} WebArena-Lite, SWE-bench Lite, MATH subset. \textbf{Metrics:} task success rate (TSR), $I_{80}$, across-seed variance, memory footprint. Tests H1.

\subsection*{Experiment 2: Tool Use and Grounding}
\textbf{Core question:} Does RL training reduce tool hallucination compared to SFT, and can agents discover tool use without demonstrations? \textbf{Design:} SFT on tool-use demonstrations; SFT + GRPO with outcome rewards; GRPO from scratch with tools. \textbf{Benchmarks:} ToolBench API tasks, WebArena tool-heavy subset. \textbf{Metrics:} valid API call rate, calls to non-existent APIs, task success. Tests H2.

\subsection*{Experiment 3: Memory Architecture (RAG vs Extended Context)}
\textbf{Core question:} For long-horizon agentic tasks, is retrieval-augmented memory more efficient than extended context windows? \textbf{Design:} 8K context; 32K context; 8K + top-5 retrieval; 8K + adaptive retrieval. \textbf{Metrics:} success on long tasks, compute per step, recall of early-step information, retrieval precision. Tests H3.

\subsection*{Experiment 4: Multi-Agent vs Single-Agent}
\textbf{Core question:} Does decomposition across specialized agents outperform a single agent with an equivalent compute budget? \textbf{Design:} single agent; planner + executor; planner + coder + reviewer; hierarchical manager--worker. \textbf{Metrics:} success per unit compute, inter-agent token overhead, failure modes.

\subsection*{Experiment 5: Scaling Efficiency (Test-Time Compute vs Model Size)}
\textbf{Core question:} Is inference-time sampling and voting more compute-efficient than training larger models~\cite{snell2024scaling}? \textbf{Design:} 8B single generation; 8B best-of-16 with majority vote; 8B best-of-16 with process-reward reranking; 70B single generation. \textbf{Metrics:} accuracy per FLOP, latency.

\section*{Summary of the Proposed Programme}

The five proposed experiments would comprise 19 conditions; with three seeds each, the full programme would require 57 runs. None of these runs, and no Stage~3 RAG integration, was carried out for this report. The proposed pipeline is: Stage~1 SFT on demonstrations; Stage~2 RL fine-tuning (PPO, DPO, or GRPO); Stage~3 retrieval augmentation.

\subsection{Infrastructure and Tools}
The executed runs used the \textbf{Tinker} managed training API with LoRA adapters. The W\&B records of the runs in Table~\ref{tab:measured_runs} (project \texttt{tinker-rl-scaling}) log model \texttt{Qwen/Qwen3-8B} (one run \texttt{Qwen/Qwen3.5-4B}), LoRA rank 32 (rank sweep 8/16/64), learning rate $3\times10^{-5}$ ($5\times10^{-6}$ for the 100-step run), and 50 training steps (100 for one run). Group size and batch size are not recorded in those W\&B configs. Experiment configs are in \texttt{platform\_tinker/atropos/configs/}.

\subsection{Reproducibility Protocol (Proposed)}
For the proposed programme we planned three fixed seeds per condition. The executed Qwen3-8B runs instead used seeds 42, 137, 256, 512, and 999 (Table~\ref{tab:measured_runs}). Hyperparameters for the executed runs are those logged in W\&B (above); the planned PPO clip ratio 0.2, DPO $\beta_{\text{DPO}}=0.1$, and warm-up ratio 0.03 apply only to the unexecuted programme.

\subsection{Statistical Analysis Plan (Proposed)}
For the proposed programme, the unit of replication would be the independent training seed. Algorithm comparisons would use seed-level paired designs with a prespecified multiplicity correction, report effect sizes with intervals, and treat non-significance as inconclusive unless an equivalence test with a stated margin is passed. None of these tests was run. The analysis actually performed is the seed-level summary in Section~\ref{sec:actual_results}.

\subsection{Baseline Comparisons (Proposed)}
Planned baselines: zero-shot, 3-shot, chain-of-thought, SFT-only, and published leaderboard results on SWE-bench, WebArena, and OSWorld.

\section{Executed Experiments: Measured GSM8K GRPO Runs}
\label{sec:actual_results}

This section reports only runs with released records. Earlier drafts of this chapter contained a results summary table, GSM8K and MATH training-curve plots, a HumanEval result for a ``Qwen3-3.8B'' model, and a statistical-analysis section (bootstrap intervals, two-proportion $z$-tests, Spearman trends, Chow breakpoint tests, Mann--Whitney tests, cross-model Bonferroni comparisons). Those items are removed: no released trace or evaluation file reproduces them, the plotted curves were hand-entered coordinates, the step-level tests treated autocorrelated steps and within-prompt rollouts as independent observations, and the measured W\&B runs below contradict their values (e.g., a Qwen3-8B training reward of 100\% versus a measured last-10 mean of 27.5--35.0\%).

\textbf{Reward.} The logged quantity is the online GRPO \emph{training reward} on sampled GSM8K training prompts. It is not held-out accuracy.

\begin{table}[htbp]
\centering
\caption{Completed Tinker GRPO runs with W\&B records (project \texttt{tinker-rl-scaling}; GSM8K training reward). Peak and last-10 values are W\&B run summaries; each row is one run. Source: \texttt{platform\_hybrid/experiments/results/review\_2026-10-18/Sem3\_Group6\_Original\_Report/measured\_runs.json}.}
\label{tab:measured_runs}
\begin{tabular}{llcccc}
\toprule
\textbf{Run} & \textbf{Model} & \textbf{Rank} & \textbf{Steps} & \textbf{Peak} & \textbf{Last-10 mean} \\
\midrule
seed042 & Qwen3-8B & 32 & 50 & 0.625 & 0.275 \\
seed137 & Qwen3-8B & 32 & 50 & 0.625 & 0.275 \\
seed256 & Qwen3-8B & 32 & 50 & 0.625 & 0.325 \\
seed512 & Qwen3-8B & 32 & 50 & 0.875 & 0.300 \\
seed999 & Qwen3-8B & 32 & 50 & 0.875 & 0.350 \\
rank8 (seed 42) & Qwen3-8B & 8 & 50 & 0.625 & 0.2125 \\
rank16 (seed 42) & Qwen3-8B & 16 & 50 & 0.750 & 0.1875 \\
rank64 (seed 42) & Qwen3-8B & 64 & 50 & 0.875 & 0.250 \\
100-step (seed 42, lr $5\times10^{-6}$) & Qwen3-8B & 32 & 100 & 0.750 & 0.275 \\
4B (seed 137) & Qwen3.5-4B & 32 & 50 & 1.000 & 0.825 \\
\bottomrule
\end{tabular}
\end{table}

\textbf{Seed-level summary.} Across the five Qwen3-8B seeds (the independent unit), the last-10 training reward is $0.305$ (SD $0.033$, $n=5$). The rank and step-count variants are single runs and support no comparison.

\textbf{Held-out GSM8K.} The five seed adapters were evaluated on the same 200 GSM8K test problems with greedy decoding (\texttt{platform\_tinker/reports/final/gsm8k\_heldout\_seed*.json}); base Qwen3-8B scored 164/200 (82.0\%). Per-seed accuracies are 83.0, 82.5, 80.5, 84.0, and 86.5\% (mean 83.3\%, SD 2.2 points). A one-sample $t$-test of the five seed accuracies against the single base score gives $t(4)=1.32$, $p=0.26$. The single base evaluation has no variance estimate, and the result neither shows a held-out gain nor establishes equivalence.

\subsection{GRPO Training Pipeline Architecture}

\begin{figure}[htbp]
\centering
\begin{tikzpicture}[
    node distance=1.4cm and 2.2cm,
    box/.style={draw, rounded corners=4pt, minimum width=2.8cm, minimum height=1.1cm,
                text centered, align=center, font=\small\sffamily, fill=blue!8, draw=blue!50!black, thick},
    cloud/.style={draw, ellipse, minimum width=2.6cm, minimum height=1.0cm,
                  text centered, align=center, font=\small\sffamily, fill=orange!10, draw=orange!60!black, thick},
    arrow/.style={-{Stealth[length=7pt]}, thick, draw=gray!70!black},
    darrow/.style={-{Stealth[length=7pt]}, thick, draw=red!60!black, dashed},
]

% Row 1: Dataset / Prompt pool
\node[cloud] (dataset) {Dataset\\(e.g.\ GSM8K)};

% Row 1: Environment
\node[box, right=of dataset] (env) {Environment\\(Atropos)};

% Row 2: Rollout server (below env)
\node[box, below=of env] (rollout) {Rollout Server\\(vLLM / SGLANG)};

% Row 2: Trainer (right of rollout)
\node[box, right=of rollout] (trainer) {Tinker Trainer\\(GRPO + LoRA)};

% Row 1: Checkpoint store (right of env)
\node[box, right=of env] (ckpt) {Checkpoint\\Store};

% Row 2: Updated policy (below ckpt)
\node[cloud, below=of ckpt] (policy) {Updated\\Policy};

% Arrows
\draw[arrow] (dataset) -- node[above, font=\tiny\sffamily]{prompts} (env);
\draw[arrow] (env) -- node[right, font=\tiny\sffamily]{batch} (rollout);
\draw[arrow] (rollout) -- node[above, font=\tiny\sffamily]{trajectories + scores} (trainer);
\draw[arrow] (trainer) -- node[right, font=\tiny\sffamily]{checkpoint} (policy);
\draw[darrow] (policy) -- node[above, font=\tiny\sffamily]{reload weights} (rollout);
\draw[arrow] (env.north east) to[bend left=18] node[above, font=\tiny\sffamily]{reward signal} (trainer.north west);
\draw[darrow] (ckpt.south) -- (policy.north);
\draw[arrow] (trainer.north) to[bend right=15] node[right, font=\tiny\sffamily]{save} (ckpt.east);

\end{tikzpicture}
\caption{GRPO training pipeline. The Atropos Environment samples prompts from the dataset and forwards batches to the Rollout Server (vLLM), which generates $G$ completions per prompt. The Environment scores each completion (binary reward) and passes trajectory--reward pairs to the Tinker Trainer, which performs the GRPO policy gradient update over LoRA parameters. Updated weights are checkpointed and reloaded into the Rollout Server every $N$ steps for on-policy training.}
\label{fig:grpo_pipeline}
\end{figure}

\section{Systematic Scaling Study (Planned)}
\label{sec:systematic_study}

\subsection{Motivation}
The executed runs vary only seed, LoRA rank, and step count on one model. Questions about model family, instruction tuning, and parameter count therefore remain open. This section records a planned three-block study; it has no results in this report. W\&B records for the planned Qwen3-0.6B and Qwen3-14B GSM8K runs (project \texttt{tinker-rl-scaling}) are all marked failed or crashed.

\subsection{Planned Design}

\begin{table}[htbp]
\centering
\caption{Planned three-block design (no results reported)}
\label{tab:systematic_design}
\begin{tabular}{llll}
\toprule
\textbf{Block} & \textbf{Varying} & \textbf{Intended to be fixed} & \textbf{Runs} \\
\midrule
A & Instruction tuning (base vs.\ instruct) & Llama-3.1 family, 8B & 2 \\
B & Model family (Qwen3 vs.\ Llama-3.1) & Base, $\sim$8B & 2 \\
C & Parameter count (0.6B -- 30B) & Qwen3 base & 6 \\
\bottomrule
\end{tabular}
\end{table}

\paragraph{Block A.} \texttt{meta-llama/Llama-3.1-8B} (base) vs.\ \texttt{meta-llama/Llama-3.1-8B-Instruct}, holding family and size fixed.

\paragraph{Block B.} \texttt{Qwen/Qwen3-8B} vs.\ \texttt{meta-llama/Llama-3.1-8B} (both base), holding instruction-tuning status and size fixed.

\paragraph{Block C.} Six Qwen3 models (0.6B, 1.7B, 4B, 8B, 14B, 30B-A3B MoE). \textbf{Confound:} the repository configurations change LoRA rank and learning rate with size (\texttt{gsm8k\_qwen\_0\_6b.yaml}: rank 16, lr $6\times10^{-5}$; \texttt{gsm8k\_qwen\_1\_7b.yaml}: rank 16, lr $5\times10^{-5}$; 8B: rank 32, lr $4\times10^{-5}$). As configured, Block C does not vary parameter count alone; a valid size ladder must fix rank and learning rate or treat them as declared co-varying factors.

\paragraph{Planned analysis.} Seeds are the replication unit. Final-window training reward would be summarized per seed, and capability claims would require held-out evaluation; steps within a run would not be treated as independent observations.

\section{Key Findings and Analysis}
\label{sec:key_findings}

\textbf{Finding 1 --- Training reward is not held-out accuracy.}
For Qwen3-8B, the 50-step GRPO runs reach peak training rewards of 0.625--0.875 but last-10 means of only 0.275--0.350 (seed mean 0.305, $n=5$). Peak values reflect single noisy steps; window means are the more stable summary.

\textbf{Finding 2 --- No demonstrated held-out GSM8K gain.}
The base model already scores 82.0\% on the 200-problem held-out slice; the five trained adapters score 80.5--86.5\% (mean 83.3\%). The one-sample test against the single base score ($p=0.26$) does not show a gain, and the design cannot establish equivalence.

\textbf{Finding 3 --- Single-run variants are descriptive only.}
LoRA rank (8/16/32/64) and the 100-step run are single-seed and differ within the seed-to-seed range of the rank-32 runs (last-10 0.275--0.350); they support no rank or horizon conclusion.

\textbf{Finding 4 --- The agentic programme remains untested.}
The WebArena, SWE-bench, OSWorld, PPO/DPO, RAG, and multi-agent experiments of Section~\ref{sec:comprehensive_experiments} were not run; the survey's conclusions about them rest on the cited literature, not on our experiments.

\begin{tcolorbox}[colback=purple!5!white, colframe=purple!60!black, title=\textbf{Chapter 5 Key Takeaways (Empirical)}]
\begin{enumerate}
    \item \textbf{Measured scope is narrow:} the executed work is low-budget LoRA GRPO on GSM8K (mostly Qwen3-8B, 50 steps).
    \item \textbf{Training reward $\neq$ accuracy:} last-10 training reward averaged 0.305 across five seeds.
    \item \textbf{Held-out effect not demonstrated:} 82.0\% base vs.\ 83.3\% mean across five adapters ($p=0.26$, single base evaluation).
    \item \textbf{Proposals are proposals:} the agentic benchmark suite, PPO/DPO baselines, and RAG experiments were designed but not executed.
\end{enumerate}
\end{tcolorbox}

\chapter{Limitations and Future Directions}
\label{ch:limitations}

\textit{This chapter provides critical reflection on both the field and our survey. Section~\ref{sec:field-limitations} identifies five structural limitations in current agentic AI research (L1--L5). Section~6.2 acknowledges our survey's own constraints. We conclude with concrete future research directions addressing these gaps.}

\vspace{0.3cm}
\section{Limitations of the Field}
\label{sec:field-limitations}
Digging into the literature, we uncovered a handful of deep-seated limitations baked into how agentic AI research is carried out today. These aren't mere implementation quirks---they point to deeper structural problems in how we evaluate, replicate work, and guard against harm:

\textbf{L1: Benchmark-Reality Gap.} Current benchmarks test simplified versions of real problems. They provide clear goals and measurable success criteria---but real deployment involves ambiguous requirements, adversarial inputs, and cascading failures.

\textit{Example:} An agent scoring 49\% on SWE-bench may have memorized common bug patterns rather than learning general debugging principles. The core issue: we have no validated correlation between benchmark scores and real-world utility. We may be overfitting to evaluation suites rather than making genuine progress.

\textbf{L2: Reproducibility Constraints.} State-of-the-art results (DeepSeek-R1's 79.8\% on AIME 2024, Claude-3.5's 49\% on SWE-bench) are achieved with proprietary training data, undisclosed hyperparameters, and compute budgets exceeding academic accessibility. Our experimental framework deliberately targets reproducibility with open models (Llama-3-8B, Qwen2.5-7B) and documented configurations, accepting lower absolute performance for methodological transparency.

\textbf{L3: Safety-Capability Tradeoff.} Every capability improvement expands the potential for harm \cite{yin2024os}. Agents capable of debugging code can introduce vulnerabilities; agents capable of web navigation can exfiltrate data. Today's safety guardrails (Constitutional AI, RLHF-based refusal) offer only partial cover. Formally verifying agent behavior at scale is still an unsolved puzzle.

\textbf{L4: Metric Inadequacy.} Task completion rates---the dominant metric---obscure reasoning quality. An agent may succeed through luck or fail despite correct reasoning for 90\% of the trajectory. Step-wise evaluation methods exist~\cite{lightman2023lets}, though few teams have adopted them in practice.

\textbf{L5: Computational Barriers.} PPO requires 4$\times$ base model memory; even GRPO consumes substantial GPU-hours for convergence. Those bills slow down iteration cycles and lock out researchers who lack big-iron clusters.

\section{Limitations of This Survey}
A few caveats apply to our analysis:

\begin{itemize}
    \item \textbf{Temporal Scope:} The core survey analyzes papers from 2023--2025. We include foundational references from earlier years (e.g., PPO~\cite{schulman2017ppo}, Constitutional AI~\cite{bai2022constitutional}) for context, but systematic coverage focuses on the target period.
    \item \textbf{Language Bias:} We mostly stuck to English-language publications, so relevant work from non-English research groups may have slipped under the radar.
    \item \textbf{Publication Bias:} Our analysis tilts toward wins; dead ends and negative findings are scarce in the published record, so they're scarce here too.
    \item \textbf{Verification Limitations:} We quote numbers straight from the papers, no independent replication. Some of those figures may not hold up in practice.
    \item \textbf{Rapid Obsolescence:} The field evolves rapidly; some findings may be superseded by the time of publication.
\end{itemize}

\subsection{Training Instabilities in RL-Based Agent Training}
\label{sec:training_instabilities}

Compute bills aside, RL training for agentic systems comes with a few landmines that practitioners need to spot and defuse:

\begin{table}[h]
\centering
\caption{Common Training Instabilities in Agentic RL (2024--2025)}
\label{tab:training_instabilities}
\small
\begin{tabular}{@{}p{2.5cm}>{\raggedright\arraybackslash}p{5.5cm}>{\raggedright\arraybackslash}p{4.5cm}@{}}
\toprule
\textbf{Instability} & \textbf{Description} & \textbf{Mitigation Strategy} \\
\midrule
\textbf{Echo Trap} & Agent learns to repeat successful patterns verbatim, losing generalization. Manifests as high training accuracy but poor transfer to novel inputs. & Temperature annealing, diversity regularization, periodic curriculum reset \\
\addlinespace
\textbf{Entropy Collapse} & Policy entropy drops prematurely, causing the agent to commit to suboptimal actions before exploring alternatives. & Entropy bonus scheduling, exploration curricula, DAPO-style dynamic sampling \\
\addlinespace
\textbf{Reward Instability} & Learned reward models drift during training, causing the policy to chase moving targets. Common in RLHF pipelines. & Reward model freezing, ensemble rewards, periodic reward calibration \\
\addlinespace
\textbf{Gradient Spikes} & Large, unexpected gradient magnitudes destabilize training. Especially problematic in long-horizon tasks. & Gradient clipping, LayerNorm, learning rate warmup, PPO's trust region \\
\addlinespace
\textbf{Tool Hallucination} & Agent generates syntactically invalid or semantically inappropriate tool calls. Worsens under distribution shift. & GRPO-TCR regularization, tool validation layers, structured generation \\
\addlinespace
\textbf{Credit Assignment Failure} & Sparse rewards in long trajectories prevent learning which actions contributed to success. & Hindsight Experience Replay, dense intermediate rewards, temporal difference methods \\
\bottomrule
\end{tabular}
\end{table}

\paragraph{Echo Trap Deep Dive} The Echo Trap (also related to model collapse \cite{shumailov2023curse}) is sneaky---it inflates training metrics while gutting real-world usefulness. When an agent discovers that a specific phrasing or action sequence consistently yields rewards, it may ``lock in'' that pattern regardless of context. Detection requires held-out evaluation sets with significant distributional differences from training data.

\paragraph{Tool Hallucination} \textit{Tool hallucination refers to when agents generate syntactically invalid or semantically inappropriate tool calls that don't match available API specifications.} This phenomenon worsens under distribution shift, requiring regularization techniques like GRPO-TCR to mitigate.

\paragraph{Entropy Collapse Prevention} GRPO's group-relative normalization naturally combats entropy collapse by ensuring diverse trajectories within each batch, but explicit entropy bonuses remain necessary for stable long-term training. The recommended approach combines an initial high entropy coefficient (0.1) with gradual decay to 0.01 over training.

\subsection{Observed Emergent Behaviors from RL Training}
\label{sec:emergent}

One of the more eye-opening discoveries from recent RL-trained agents is the pop-up of clever behaviors nobody coded in. This matters because it reveals what RL optimization actually teaches: not just specific actions, but general problem-solving strategies. When test-based rewards drive systematic debugging or compute-penalties drive cost-aware tool selection, the takeaway is that reward signals can mold meta-cognitive strategies. These emergent behaviors provide insights into what's learnable:

\textbf{Code Self-Correction in Testing Loops:} Agents trained on SWE-bench develop iterative debugging strategies: writing initial fixes, running tests, analyzing failures, and refining solutions---without explicit training on this loop. The behavior emerges from outcome supervision alone, a sign that test-based rewards nudge agents toward methodical debugging.

\textbf{Strategic Tool Selection:} Rather than always using the most powerful tool, RL-trained agents learn context-dependent selection. For simple calculations, they use basic operations; for complex queries, they invoke specialized APIs. This cost-aware behavior emerges from compute-penalized rewards without explicit efficiency training.

\textbf{Reasoning Pattern Regularities:} DeepSeek-R1 and similar GRPO-trained models develop consistent reasoning templates: problem decomposition, intermediate verification, and confidence calibration. These patterns pop up across problem types, hinting that RL can uncover broadly useful meta-cognitive tricks.

\textbf{Error Anticipation:} Agents exposed to adversarial tool environments (via AT-GRPO) learn proactive error handling---checking preconditions before actions, validating outputs before proceeding, and maintaining fallback plans. This defensive programming style emerges from training on environments with realistic failure modes.

\textbf{Communication Protocol Development (Multi-Agent):} In multi-agent settings without predefined communication channels, agents develop shorthand protocols for common coordination patterns. These emergent ``languages'' are task-specific but remarkably consistent across training runs, hinting that the push to coordinate keeps landing agents on the same answer.

What does this mean in practice? Carefully crafted rewards can coax out nuanced skills even when we never hard-coded them. The flip side is that unprogrammed behaviors are trickier to foresee and harder to inspect---an alignment headache we cannot ignore.

\subsection{Research Gaps Analysis}

Sifting through the literature, we spot glaring gaps between what today's best systems can do and what production-grade agents actually need:

\begin{table}[h]
\centering
\caption{Critical Research Gaps in Agentic AI (Data as of January 2026)}
\label{tab:research_gaps}
\small
\begin{tabular}{@{}p{2.5cm}>{\raggedright\arraybackslash}p{4cm}>{\raggedright\arraybackslash}p{4cm}>{\raggedright\arraybackslash}p{2.5cm}@{}}
\toprule
\textbf{Gap Category} & \textbf{Verified State (Dec 2025)} & \textbf{Required State} & \textbf{Impact} \\
\midrule
Benchmark Generalization & 39--49\% on SWE-bench, 45\% on WebArena & $>$90\% on novel environments & Critical \\
\addlinespace
Long-Horizon Planning & 20--50 step trajectories (DeepAgent) & 100+ step coherent execution & High \\
\addlinespace
Error Recovery & Emergent self-correction (ReTool) & Robust autonomous recovery & Medium \\
\addlinespace
Safety Guarantees & Post-hoc filtering + AT-GRPO & Formal verification & Critical \\
\addlinespace
Multi-Modal Grounding & 37.7\% OSWorld (Mobile-Agent-v3) & Full sensorimotor integration & High \\
\addlinespace
Sample Efficiency & 1K--10K trajectories (GRPO) & $<$1K trajectories & Medium \\
\addlinespace
Interpretability & CoT reasoning visible & Auditable + verifiable chains & High \\
\bottomrule
\end{tabular}
\vspace{0.5em}
\begin{minipage}{\textwidth}
\footnotesize\textit{Note: ``Verified State'' reflects published benchmark results as of December 2025. Rapid progress continues; check official leaderboards for latest results.}
\end{minipage}
\end{table}

\begin{figure}[htbp]
\centering
\begin{tikzpicture}[scale=0.95, transform shape]
    % === AXES ===
    \draw[->, thick, gray!70] (-5.5, 0) -- (5.5, 0) node[right, font=\small] {Research Maturity};
    \draw[->, thick, gray!70] (0, -3.5) -- (0, 4) node[above, font=\small] {Impact on Deployment};

    % Axis labels
    \node[font=\tiny, gray] at (-4.5, -0.35) {Early Stage};
    \node[font=\tiny, gray] at (4.5, -0.35) {Mature};
    \node[font=\tiny, gray, rotate=90] at (-0.4, -2.5) {Low};
    \node[font=\tiny, gray, rotate=90] at (-0.4, 3) {Critical};

    % === QUADRANT BACKGROUNDS ===
    % Top-left: Critical priority (high impact, low maturity)
    \fill[red!15, rounded corners] (-5.3, 0.15) rectangle (-0.15, 3.7);
    \node[font=\scriptsize\bfseries, red!70!black] at (-2.7, 3.4) {Critical Priority};

    % Top-right: Advancing well (high impact, high maturity)
    \fill[green!15, rounded corners] (0.15, 0.15) rectangle (5.3, 3.7);
    \node[font=\scriptsize\bfseries, green!60!black] at (2.7, 3.4) {Advancing Well};

    % Bottom-left: Long-term research (low impact currently, low maturity)
    \fill[orange!15, rounded corners] (-5.3, -3.3) rectangle (-0.15, -0.15);
    \node[font=\scriptsize\bfseries, orange!70!black] at (-2.7, -3.0) {Long-term Research};

    % Bottom-right: Maintenance mode (low impact, high maturity)
    \fill[blue!10, rounded corners] (0.15, -3.3) rectangle (5.3, -0.15);
    \node[font=\scriptsize\bfseries, blue!60!black] at (2.7, -3.0) {Maintenance Mode};

    % === RESEARCH GAPS AS NODES ===
    % Critical Priority (top-left)
    \node[draw, fill=red!30, rounded corners, font=\tiny, align=center, inner sep=4pt] at (-3.5, 2.5) {Safety\\Guarantees};
    \node[draw, fill=red!30, rounded corners, font=\tiny, align=center, inner sep=4pt] at (-1.5, 2.8) {Benchmark\\Generalization};
    \node[draw, fill=red!30, rounded corners, font=\tiny, align=center, inner sep=4pt] at (-4.2, 1.3) {Robustness-\\Capability};

    % Advancing Well (top-right)
    \node[draw, fill=green!30, rounded corners, font=\tiny, align=center, inner sep=4pt] at (2.0, 2.5) {Tool\\Learning};
    \node[draw, fill=green!30, rounded corners, font=\tiny, align=center, inner sep=4pt] at (4.0, 2.0) {GRPO\\Algorithms};
    \node[draw, fill=green!30, rounded corners, font=\tiny, align=center, inner sep=4pt] at (3.0, 1.0) {Code\\Agents};

    % Long-term Research (bottom-left)
    \node[draw, fill=orange!30, rounded corners, font=\tiny, align=center, inner sep=4pt] at (-3.8, -1.5) {Continual\\Learning};
    \node[draw, fill=orange!30, rounded corners, font=\tiny, align=center, inner sep=4pt] at (-1.8, -2.2) {Emergent\\Coordination};
    \node[draw, fill=orange!30, rounded corners, font=\tiny, align=center, inner sep=4pt] at (-3.0, -0.7) {Latent\\Reasoning};

    % Maintenance Mode (bottom-right)
    \node[draw, fill=blue!25, rounded corners, font=\tiny, align=center, inner sep=4pt] at (2.5, -1.5) {Basic\\Prompting};
    \node[draw, fill=blue!25, rounded corners, font=\tiny, align=center, inner sep=4pt] at (4.2, -0.8) {RLHF\\Foundations};

    % High-impact, medium maturity items (near center-top)
    \node[draw, fill=yellow!40, rounded corners, font=\tiny, align=center, inner sep=4pt] at (0.8, 1.8) {Long-Horizon\\Planning};
    \node[draw, fill=yellow!40, rounded corners, font=\tiny, align=center, inner sep=4pt] at (-0.8, 0.8) {Multi-Modal\\Grounding};

    % === LEGEND ===
    \draw[thick, fill=white, rounded corners] (-5.3, -4.8) rectangle (5.3, -3.7);
    \node[font=\tiny\bfseries] at (0, -3.95) {Priority Guide:};
    \node[draw, fill=red!30, rounded corners, minimum width=0.4cm, minimum height=0.25cm] at (-3.5, -4.4) {};
    \node[right, font=\tiny] at (-3.2, -4.4) {Urgent};
    \node[draw, fill=yellow!40, rounded corners, minimum width=0.4cm, minimum height=0.25cm] at (-1.0, -4.4) {};
    \node[right, font=\tiny] at (-0.7, -4.4) {Active};
    \node[draw, fill=green!30, rounded corners, minimum width=0.4cm, minimum height=0.25cm] at (1.3, -4.4) {};
    \node[right, font=\tiny] at (1.6, -4.4) {Progressing};
    \node[draw, fill=blue!25, rounded corners, minimum width=0.4cm, minimum height=0.25cm] at (3.5, -4.4) {};
    \node[right, font=\tiny] at (3.8, -4.4) {Solved};
\end{tikzpicture}
\caption{Research gap landscape for agentic AI (Jan 2025). Gaps are positioned by research maturity (x-axis) and deployment impact (y-axis). The \textbf{Critical Priority} quadrant requires immediate attention: safety, benchmarks, and robustness.}
\label{fig:research_gaps_quadrant}
\end{figure}

\textbf{Gap 1: Benchmark-Reality Divide.} Current benchmarks (SWE-bench, WebArena) represent sanitized slices of real-world complexity. Agents achieving 35\% on these benchmarks face qualitatively different challenges in production: ambiguous requirements, adversarial users, and cascading failures. Closing this gap requires benchmarks with realistic noise, incomplete specifications, and multi-stakeholder constraints.

\textbf{Gap 2: Compositional Generalization.} Agents trained on task distributions struggle with novel compositions of familiar primitives. An agent that can ``search'' and ``click'' may fail at ``search then click on the third result.'' Such fragility hints that today's RL losses chase surface-level patterns instead of building transferable primitives.

\textbf{Gap 3: Long-Horizon Credit Assignment.} Current methods (PPO, GRPO) falter when assigning credit past the 10--20 step mark. Hierarchical RL and intrinsic motivation offer partial solutions, yet no current solution can dependably handle 100-step tasks demanding sustained, coherent planning.

\textbf{Gap 4: Robustness-Capability Frontier.} Increasing agent capability systematically increases potential for harm. Current approaches treat safety as a post-hoc constraint rather than an architectural principle. We need architectures that are safe by design---ones where boosting capability doesn't widen the attack surface.

\section{Future Research Directions}

\textbf{Latent Reasoning:} Current CoT approaches expose reasoning in output tokens, adding latency and token costs. Research into latent reasoning---where models ``think'' internally without generating visible tokens---might slash latency without sacrificing reasoning accuracy.

\textbf{Emergent Multi-Agent Coordination:} Moving beyond hand-designed role assignments toward agents that self-organize, develop specialized roles, and establish communication protocols through learning remains an open challenge.

\textbf{Continual Learning:} Agents that adapt to changing environments, learn from deployment feedback, and avoid catastrophic forgetting would push practical usefulness up a notch.

\textbf{Verifiable Reasoning:} Connecting RL-trained reasoning with formal verification (producing proofs or certificates of correctness) could unlock trustworthy autonomy for high-stakes deployments.

\textbf{Efficient Exploration:} Developing exploration strategies that balance curiosity-driven learning with safety constraints, especially when agents touch the live world, is a must.

\begin{tcolorbox}[colback=red!5!white, colframe=red!60!black, title=\textbf{Chapter 6 Key Takeaways}]
\begin{enumerate}
    \item \textbf{Benchmarks $\neq$ deployment:} No validated correlation between scores and real-world utility.
    \item \textbf{Reproducibility crisis:} SOTA results require undisclosed data + compute beyond academic reach.
    \item \textbf{Safety-capability tension:} Every capability gain expands harm potential---no free lunch.
    \item \textbf{Open problems:} Verifiable reasoning, continual learning, and emergent coordination remain unsolved.
\end{enumerate}
\end{tcolorbox}

\chapter{Conclusion}
\label{ch:conclusion}

\section*{Summary of Findings}
This systematic review of 60+ papers on reinforcement learning for autonomous AI agents documents a rapidly advancing field. The 2024--2025 period saw researchers pivot sharply from prompt-only techniques toward RL-driven fine-tuning, with large reported gains on major evaluation suites (e.g., WebArena-Lite 8.5\%$\rightarrow$44.8\% for Llama-3.1-8B under end-to-end RL; AIME 2024 9.3\%$\rightarrow$79.8\% from GPT-4o to DeepSeek-R1).

\section*{Key Takeaways}

\noindent\textbf{1. The GRPO Paradigm Shift.} GRPO and its spin-offs are now common in the papers we reviewed. By eliminating the value/critic network (using group normalization) and leveraging verifiable rewards without a learned reward model, GRPO halves memory requirements (policy + reference model only, vs.\ four networks for PPO). This efficiency enabled DeepSeek-R1 to achieve 79.8\% on AIME 2024---putting capable agent training within reach of smaller labs. The proliferation of specialized variants (Flow-GRPO for multi-turn dialogue, M-GRPO for hierarchical coordination, AT-GRPO for adversarial robustness) suggests this lightweight approach will define near-term development.

\noindent\textbf{2. Tool Mastery Is Now Baseline.} Agents without retrieval and tool-calling abilities encounter fundamental capability limitations. The emergence of RL-optimized tool learning frameworks---ReTool achieving 67\% on AIME 2024, DeepAgent executing 100+ step trajectories---demonstrates that tool integration is now essential for production deployment.

\noindent\textbf{3. Multi-Agent Coordination Matures.} When single agents fail, coordinated multi-agent systems often succeed. MATPO's 18.38\% improvement over single-agent baselines and Mobile-Agent-v3's 73.3\% on OSWorld show that RL-trained multi-agent coordination is now production-ready. \textit{Caveat:} These gains come with coordination overhead---increased latency, $\mathcal{O}(n^2)$ communication costs for fully-connected topologies, and failure propagation risks---that deployment teams must factor into system design.

\noindent\textbf{4. Emergent Capabilities and Risks.} RL training produces behaviors not explicitly programmed: code self-correction, strategic tool selection, communication protocol development, and cross-validation strategies. Put differently, a well-shaped reward can unlock nuanced skills---yet these unplanned behaviors are difficult to anticipate or audit, presenting significant challenges for alignment research.

\section*{Implications for Research and Practice}

For \textbf{researchers}, we flag seven gaps that matter because they block real-world deployment---they are not mere academic footnotes. The benchmark-reality divide means improvements on sanitized benchmarks don't translate to messy reality—this undermines the entire evaluation paradigm. Long-horizon planning fails beyond 20 steps because credit assignment breaks down—this is algorithmic, not a data problem. Safety verification remains informal because formal methods face state-space explosion—this is computational complexity, not implementation. Multi-modal grounding fails because vision and language representations don't align naturally—this is architectural. Sample efficiency remains low because every environment interaction is expensive—this is economic. Interpretability fails because emergent behaviors are unpredictable—this is foundational. Each gap reflects a structural constraint, not a missing feature.

For \textbf{practitioners}, we recommend a three-stage training pipeline (SFT$\rightarrow$RL$\rightarrow$RAG) with GRPO as the default algorithm for resource-constrained settings. Our ablation grid enables systematic evaluation of design choices within typical academic compute budgets.

For \textbf{policymakers}, the rapid capability gains documented here---with minimal additional training overhead---make it clear that governance frameworks need to land before mass rollout. The emergent-behavior findings, in particular, deserve a close look.

\section*{Broader Impact}

Autonomous AI agents mark a tectonic shift---from passive tools to goal-chasing systems. Tools respond to commands; agents pursue objectives over extended time horizons without step-by-step instruction. This is not just incremental improvement—it's a different paradigm. The implications are profound because the failure modes are different: tools fail by not responding; agents fail by pursuing the wrong goal, by getting stuck in local optima, by taking unsafe shortcuts. The risks scale with capability because the action space grows exponentially with horizon length. This creates a safety-capability frontier: capability improvements without safety measures are dangerous; safety measures that cripple capability are useless. The tradeoff is structural, not optional. This means:

\begin{itemize}
    \item \textbf{Positive:} Democratized automation of complex knowledge work, improved accessibility, and augmented human capability across software engineering, research, and creative domains.
    \item \textbf{Risks:} Potential for cascading failures in agentic chains, difficulty auditing emergent behaviors, and amplification of existing biases in autonomous decision-making.
\end{itemize}

The current approach of capability-first, safety-later is fundamentally flawed. Once you have a capable system, retrofitting safety is like adding seatbelts to a car after it's built—possible, but much harder than designing them in from the start. Worse, emergent behaviors are unpredictable. You cannot bolt on interpretability after the fact if the system's reasoning is opaque. Robustness cannot be added as a layer if the core architecture is brittle. Safety verification cannot be done post-hoc if the state space is intractable. These must be co-equal objectives from the beginning because they constrain the design space. A system optimized purely for capability will inevitably use shortcuts that violate safety—this is optimization pressure, not malice. The only solution is to bake in constraints from the start.

\section*{Future Directions}

Our planned experiments will rigorously test the hypotheses presented in this survey: comparing Llama-3-8B and Qwen2.5-7B across PPO, DPO, and GRPO on standardized benchmarks, with ablations isolating the contributions of tool integration, augmented retrieval strategies, and multi-agent coordination.
 We hope other researchers will replicate, extend, and critique what we have started here.

The goal isn't abstract or academic. Building agents that remain both capable and safe over long, messy interactions is the defining challenge because shorter interactions are already solved. The fundamental issue is scaling: as interaction length grows linearly, the probability of failure grows exponentially. Each step introduces new error modes: tool failures, misinterpretations, compounding inaccuracies, unexpected states. A 100-step interaction isn't 10x harder than a 10-step interaction—it's exponentially harder because error recovery must work across much longer horizons. "Messy" matters because clean, well-defined tasks are already tractable. The real world is ambiguous, adversarial, and constantly changing. Solving this requires agents that can handle uncertainty, detect their own failures, and recover autonomously. This isn't incremental progress; it's crossing a threshold from narrow AI to systems that can actually operate in the real world.

\chapter*{Acknowledgments}
\addcontentsline{toc}{chapter}{Acknowledgments}

We thank our mentors Narayana Darapaneni (Director, AIML) and Anwesh Reddy Paduri (Senior Data Scientist) for their guidance throughout this project. We acknowledge the open-source community for frameworks including Hugging Face Transformers, TRL, and Unsloth that enabled our experimental framework. We thank the authors of SWE-bench, WebArena, OSWorld, and GAIA for making their benchmarks publicly available, and the DeepSeek team for open-sourcing DeepSeek-R1. Compute resources were provided by Google Colab Pro+. Finally, we thank the anonymous reviewers and our peers in the M.Tech AI/ML program at PES University for their valuable feedback on earlier drafts.

\section*{Author Contributions (CRediT)}
\noindent\textbf{Conceptualization:} All authors. \textbf{Methodology:} Arvind C R, Madhu Kumara L, Sandhya Jeyaraj. \textbf{Literature Review:} All authors contributed equally. \textbf{Formal Analysis:} Mohammad Rafi, Arumugam Chetty K. \textbf{Writing---Original Draft:} All authors. \textbf{Writing---Review \& Editing:} Arvind C R, Dhruva. \textbf{Visualization:} Sandhya Jeyaraj, Madhu Kumara L. \textbf{Supervision:} Narayana Darapaneni, Anwesh Reddy Paduri.

\section*{Declaration of Competing Interests}
The authors declare that they have no known competing financial interests or personal relationships that could have appeared to influence the work reported in this paper.

\section*{Data Availability}
This survey is based on publicly available research papers. The bibliography file and supplementary materials are available upon request from the corresponding author. All referenced benchmarks (SWE-bench, WebArena, OSWorld, GAIA) are publicly accessible through their respective official repositories.

\section*{Code Availability}
The experimental framework described in Chapter~\ref{ch:experiments} is a proposed design for future work. Implementation code and experimental scripts will be made publicly available upon completion of experiments.

% ========== APPENDIX ==========
\appendix

\chapter{Glossary of Acronyms}
\label{app:glossary}

\begin{longtable}{@{}p{2.5cm}p{5cm}p{7cm}@{}}
\toprule
\textbf{Acronym} & \textbf{Full Term} & \textbf{Brief Definition} \\
\midrule
\endfirsthead
\toprule
\textbf{Acronym} & \textbf{Full Term} & \textbf{Brief Definition} \\
\midrule
\endhead
\midrule
\multicolumn{3}{r}{\textit{Continued on next page...}} \\
\endfoot
\bottomrule
\endlastfoot

\multicolumn{3}{l}{\textbf{Core Algorithms}} \\
\midrule
PPO & Proximal Policy Optimization & RL algorithm using clipped objective to ensure stable policy updates \\
DPO & Direct Preference Optimization & Optimizes models using preference pairs directly without needing a separate reward network
 \\
GRPO & Group Relative Policy Optimization & Memory-efficient RL using within-batch ranking instead of critic network \\
RLHF & RL from Human Feedback & Alignment framework using human preferences as a guiding reward signal
 \\
TRPO & Trust Region Policy Optimization & PPO's predecessor with hard KL constraint \\
SFT & Supervised Fine-Tuning & Initial training stage using demonstration data \\
\midrule

\multicolumn{3}{l}{\textbf{Architectures \& Techniques}} \\
\midrule
RAG & Retrieval-Augmented Generation & Combining LLM generation with external knowledge retrieval \\
CoT & Chain-of-Thought & Prompting technique eliciting step-by-step reasoning \\
ReAct & Reasoning + Acting & Framework interleaving reasoning traces with environment actions \\
MDP & Markov Decision Process & Mathematical framework for sequential decision-making \\
POMDP & Partially Observable MDP & MDP where agent cannot observe full state \\
PRM & Step-Level Reward Model & evaluates individual deduction steps rather than just the final result \\
ORM & Outcome Reward Model & Scores only final answers/outcomes \\
\midrule

\multicolumn{3}{l}{\textbf{Benchmarks}} \\
\midrule
SWE-bench & Software Engineering Bench & Evaluates code repair on real GitHub issues \\
WebArena & Web Arena & Tests web navigation and interaction tasks \\
OSWorld & Operating System World & Desktop GUI automation benchmark \\
GAIA & General AI Assistants & Multi-tool reasoning benchmark (466 tasks) \\
AIME & American Invitational Math Exam & Competition math problems used to evaluate reasoning \\
\midrule

\multicolumn{3}{l}{\textbf{Metrics \& Concepts}} \\
\midrule
KL & Kullback-Leibler (divergence) & Measures difference between probability distributions \\
TSR & Task Success Rate & Primary metric: fraction of tasks completed correctly \\
$\epsilon$ & Epsilon (clipping threshold) & PPO hyperparameter bounding policy change (typically 0.1--0.2) \\
$\beta_{\text{KL}}$ & Beta (KL penalty) & PPO/GRPO coefficient constraining drift from reference model \\
$\beta_{\text{DPO}}$ & Beta (preference strength) & DPO hyperparameter controlling preference signal strength \\
$\gamma$ & Gamma (discount factor) & Weights future vs. immediate rewards (0--1) \\
\end{longtable}

\chapter{Hyperparameter Reference}
\label{app:hyperparams}

This appendix provides recommended hyperparameter ranges based on best practices from the surveyed literature. Values represent starting points; task-specific tuning is expected.

\section{PPO Hyperparameters}

\begin{table}[htbp]
\centering
\caption{PPO hyperparameter recommendations for agentic fine-tuning}
\label{tab:ppo_hyperparams}
\begin{tabular}{@{}lccl@{}}
\toprule
\textbf{Parameter} & \textbf{Symbol} & \textbf{Recommended Range} & \textbf{Notes} \\
\midrule
Clipping threshold & $\epsilon$ & 0.1--0.2 & Lower for stability, higher for exploration \\
Discount factor & $\gamma$ & 0.99--0.999 & Higher for long-horizon tasks \\
GAE lambda & $\lambda$ & 0.95--0.99 & Bias-variance trade-off in advantage estimation \\
Value loss coefficient & $c_1$ & 0.5--1.0 & Relative weight of critic vs. policy loss \\
Entropy coefficient & $c_2$ & 0.001--0.01 & Encourages exploration; decay over training \\
Learning rate & $\alpha$ & 1e-6 to 5e-5 & Smaller for larger models \\
Batch size & --- & 32--512 & Memory-constrained; larger improves stability \\
Mini-batch size & --- & 4--64 & Gradient accumulation if memory-limited \\
PPO epochs & --- & 2--10 & Per batch; more risks overfitting \\
KL penalty coefficient & $\beta_{\text{KL}}$ & 0.01--0.1 & Optional; constrains drift from reference \\
\bottomrule
\end{tabular}
\end{table}

\section{DPO Hyperparameters}

\begin{table}[htbp]
\centering
\caption{DPO hyperparameter recommendations}
\label{tab:dpo_hyperparams}
\begin{tabular}{@{}lccl@{}}
\toprule
\textbf{Parameter} & \textbf{Symbol} & \textbf{Recommended Range} & \textbf{Notes} \\
\midrule
Preference strength & $\beta$ & 0.1--0.5 & Lower for noisy labels, higher for clean data \\
Learning rate & $\alpha$ & 1e-7 to 5e-6 & Conservative; DPO is sensitive to LR \\
Batch size & --- & 8--64 & Pair-wise data doubles effective size \\
Warmup ratio & --- & 0.1--0.2 & Fraction of training with LR warmup \\
Label smoothing & --- & 0.0--0.1 & Helps with overconfident preferences \\
Max sequence length & --- & 2048--8192 & Task-dependent; longer for code/reasoning \\
\bottomrule
\end{tabular}
\end{table}

\section{GRPO Hyperparameters}

\begin{table}[htbp]
\centering
\caption{GRPO hyperparameter recommendations}
\label{tab:grpo_hyperparams}
\begin{tabular}{@{}lccl@{}}
\toprule
\textbf{Parameter} & \textbf{Symbol} & \textbf{Recommended Range} & \textbf{Notes} \\
\midrule
Group size & $n$ & 8--64 & More samples $\rightarrow$ better baseline estimate \\
KL penalty & $\beta_{\text{KL}}$ & 0.01--0.05 & Prevents drift from reference model \\
Learning rate & $\alpha$ & 1e-6 to 1e-5 & Similar to PPO range \\
Temperature & $\tau$ & 0.7--1.0 & For response sampling diversity \\
Reward normalization $\epsilon$ & $\epsilon$ & 1e-8 & Numerical stability in std normalization \\
Gradient clipping & --- & 1.0 & Max gradient norm \\
\bottomrule
\end{tabular}
\end{table}

\chapter{Compute Requirements}
\label{app:compute}

This appendix documents computational resources required for training and inference across different scales, supporting reproducibility and resource planning.

\section*{Training Compute Estimates}

\begin{table}[htbp]
\centering
\caption{Estimated compute requirements for RL fine-tuning by model scale}
\label{tab:compute_training}
\begin{tabular}{@{}lccccc@{}}
\toprule
\textbf{Model Scale} & \textbf{Algorithm} & \textbf{GPU Memory} & \textbf{GPU-Hours} & \textbf{Hardware} & \textbf{Cost Est.} \\
\midrule
7B parameters & PPO & 80 GB $\times$ 4 & 100--200 & 4$\times$ A100 & \$400--800 \\
7B parameters & DPO & 40 GB $\times$ 2 & 50--100 & 2$\times$ A100 & \$200--400 \\
7B parameters & GRPO & 40 GB $\times$ 2 & 80--150 & 2$\times$ A100 & \$300--600 \\
\midrule
13B parameters & PPO & 80 GB $\times$ 8 & 200--400 & 8$\times$ A100 & \$800--1600 \\
13B parameters & DPO & 80 GB $\times$ 4 & 100--200 & 4$\times$ A100 & \$400--800 \\
13B parameters & GRPO & 80 GB $\times$ 4 & 150--300 & 4$\times$ A100 & \$600--1200 \\
\midrule
70B parameters & PPO & 80 GB $\times$ 32 & 1000--2000 & 32$\times$ A100 & \$4K--8K \\
70B parameters & DPO & 80 GB $\times$ 16 & 400--800 & 16$\times$ A100 & \$1.6K--3.2K \\
70B parameters & GRPO & 80 GB $\times$ 16 & 600--1200 & 16$\times$ A100 & \$2.4K--4.8K \\
\bottomrule
\end{tabular}

\vspace{0.3cm}
\small\textit{Notes:} GPU-hours assume 10K training steps. Costs based on \$4/GPU-hour (cloud A100). Memory estimates assume full precision; quantization (QLoRA) can reduce by 4$\times$. Actual requirements vary by batch size, sequence length, and optimizer state sharding.
\end{table}

\section*{Inference Compute Estimates}

\begin{table}[htbp]
\centering
\caption{Inference requirements for agentic evaluation}
\label{tab:compute_inference}
\begin{tabular}{@{}lcccc@{}}
\toprule
\textbf{Benchmark} & \textbf{Tasks} & \textbf{Avg. Steps/Task} & \textbf{GPU-Hours (7B)} & \textbf{GPU-Hours (70B)} \\
\midrule
SWE-bench Verified & 500 & 15--30 & 8--16 & 40--80 \\
WebArena & 812 & 10--25 & 10--20 & 50--100 \\
OSWorld & 369 & 20--40 & 12--24 & 60--120 \\
AIME 2024 & 30 & 1--5 & 0.5--1 & 2--5 \\
GAIA & 466 & 5--15 & 4--8 & 20--40 \\
\bottomrule
\end{tabular}
\end{table}

\section*{Academic Resource Recommendations}

For researchers with limited compute budgets:
\begin{itemize}
    \item \textbf{Minimum viable setup:} 2$\times$ RTX 4090 (24 GB each) with QLoRA can fine-tune 7-billion parameter models using GRPO
    \item \textbf{Recommended academic setup:} 4$\times$ A100 (40 GB) enables full fine-tuning of 7-billion parameter models and QLoRA for 13B
    \item \textbf{Cloud alternatives:} Lambda Labs, RunPod, and Vast.ai offer A100s at \$1--2/GPU-hour for burst workloads
    \item \textbf{Free resources:} Google Colab Pro+ provides 40 GB A100 access; Kaggle offers 30 GPU-hours/week
\end{itemize}

\chapter{Reproducibility Checklist}
\label{app:reproducibility}

This checklist separates the executed GSM8K GRPO runs from the \textit{proposed, unexecuted} experiments of Chapter~\ref{ch:experiments} (Llama-3-8B, Qwen2.5-7B on agentic benchmarks). We follow NeurIPS reproducibility guidelines and ML community best practices.

\textit{Note: This checklist applies to our reproducible experimental framework---not to proprietary systems surveyed in the literature (GPT-4, Claude-3.5, etc.) for which open weights, training data, and full hyperparameters are unavailable.}

\section*{Code and Data Availability}

\begin{tabular}{@{}p{0.7\textwidth}cc@{}}
\toprule
\textbf{Item} & \textbf{Status} & \textbf{Location} \\
\midrule
Training code and configuration files (executed GSM8K runs) & \ding{51} & GitHub repository \\
Evaluation scripts for SWE-bench, WebArena, OSWorld & \ding{55} & Not built (programme not executed) \\
GSM8K held-out evaluation records (executed runs) & \ding{51} & \texttt{platform\_tinker/reports/final/} \\
Trained checkpoints for proposed Llama-3-8B / Qwen2.5-7B runs & \ding{55} & Not trained \\
Raw logs for executed runs (W\&B) & partial & W\&B project \texttt{tinker-rl-scaling} \\
Random seeds for executed runs & \ding{51} & W\&B run configs \\
Hardware specifications (managed Tinker API) & partial & Not disclosed by the platform \\
Software dependencies & \ding{51} & GitHub repository \\
\bottomrule
\end{tabular}

\vspace{0.3cm}
\noindent\textbf{Repository:} \texttt{pes-llm-research/tinker-rl-lab} (GitHub).

\noindent\textit{Note: The protocol items below describe the proposed programme of Chapter~\ref{ch:experiments}; unchecked boxes are planned, not completed.}

\section*{Experimental Protocol Checklist}

\begin{enumerate}[leftmargin=*]
    \item \textbf{Randomization}
    \begin{itemize}[nosep]
        \item[$\square$] Fixed random seeds: 42, 123, 456 (three runs per experiment)
        \item[$\square$] Seed controls: PyTorch, NumPy, Python random, CUDA
        \item[$\square$] Deterministic algorithms enabled where possible
    \end{itemize}

    \item \textbf{Hyperparameter Selection}
    \begin{itemize}[nosep]
        \item[$\square$] Grid search over learning rate: \{1e-6, 5e-6, 1e-5\}
        \item[$\square$] Grid search over batch size: \{16, 32, 64\}
        \item[$\square$] Validation set for early stopping (10\% held out)
        \item[$\square$] Final hyperparameters documented in Appendix~\ref{app:hyperparams}
    \end{itemize}

    \item \textbf{Statistical Reporting}
    \begin{itemize}[nosep]
        \item[$\square$] Mean $\pm$ standard deviation across 3 seeds
        \item[$\square$] 95\% confidence intervals via bootstrap (1000 resamples)
        \item[$\square$] Statistical significance tests (McNemar's, $p < 0.01$)
        \item[$\square$] Effect sizes reported (Cohen's $d$ for continuous metrics)
    \end{itemize}

    \item \textbf{Benchmark Evaluation}
    \begin{itemize}[nosep]
        \item[$\square$] Official evaluation harnesses used (no custom metrics)
        \item[$\square$] Test sets never seen during training or validation
        \item[$\square$] Identical prompts across all compared methods
        \item[$\square$] Temperature = 0 for deterministic generation during eval
    \end{itemize}

    \item \textbf{Compute Transparency}
    \begin{itemize}[nosep]
        \item[$\square$] Total GPU-hours reported per experiment
        \item[$\square$] Carbon footprint estimated (ML CO2 Impact calculator)
        \item[$\square$] Wall-clock time logged
        \item[$\square$] Hardware failure/restart documentation
    \end{itemize}
\end{enumerate}

\section*{Known Limitations and Failure Modes}

\begin{itemize}
    \item \textbf{Benchmark variance:} SWE-bench results can vary $\pm$3\% across runs due to LLM non-determinism even with temperature=0 (floating-point ordering effects)
    \item \textbf{Environment drift:} WebArena uses live websites; results may not reproduce exactly if site layouts change
    \item \textbf{API dependencies:} Some experiments require API access (OpenAI, Anthropic) which may have version-dependent behavior
    \item \textbf{Hardware sensitivity:} Results on consumer GPUs (RTX 4090) may differ from A100 due to different numerical precision handling
\end{itemize}

\chapter{Paper-to-Taxonomy Mapping}
\label{app:paper_mapping}

This appendix provides a comprehensive mapping of all surveyed papers to our 19-category taxonomy (Figure~\ref{fig:rl_taxonomy}). The mapping enables systematic analysis of research distribution across the field and supports reproducibility of our literature classification.

\section*{Taxonomy Categories Reference}
\begin{multicols}{2}
\footnotesize
\noindent\textbf{Core Algorithms}
\begin{enumerate}
\item Policy-Based Methods
\item Value-Based Methods
\item Model-Based Methods
\end{enumerate}

\noindent\textbf{Learning Paradigms}
\begin{enumerate}
\setcounter{enumi}{3}
\item RLHF \& Alignment
\item Imitation Learning
\item Offline RL
\end{enumerate}

\noindent\textbf{Advanced Methods}
\begin{enumerate}
\setcounter{enumi}{6}
\item Multi-Agent RL
\item Hierarchical RL
\item Meta-Learning
\item Safe RL
\item Exploration
\end{enumerate}

\columnbreak
\noindent\textbf{Applications}
\begin{enumerate}
\setcounter{enumi}{11}
\item Agentic RL
\item Robotics
\item Reasoning
\item Security \& Safety
\item Benchmarks
\end{enumerate}

\noindent\textbf{Infrastructure}
\begin{enumerate}
\setcounter{enumi}{16}
\item Frameworks
\item Surveys
\item Open-Ended Learning
\end{enumerate}
\end{multicols}

\section*{Complete Paper Mapping}

{\scriptsize
\renewcommand{\arraystretch}{1.15}
\begin{longtable}{|>{\raggedright\arraybackslash\hspace{0pt}}p{2.3cm}|>{\raggedright\arraybackslash\hspace{0pt}}p{7.0cm}|c|>{\raggedright\arraybackslash\hspace{0pt}}p{2.4cm}|}
\caption{Complete paper-to-taxonomy mapping (59 papers).}\label{tab:paper_mapping}\\
\hline
\textbf{Citation Key} & \textbf{Paper Title} & \textbf{\#} & \textbf{Category} \\
\hline
\endfirsthead
\hline
\textbf{Citation Key} & \textbf{Paper Title} & \textbf{\#} & \textbf{Category} \\
\hline
\endhead
\hline
\multicolumn{4}{r}{\emph{Continued on next page...}} \\
\endfoot
\hline
\endlastfoot

% === CORE ALGORITHMS: Policy-Based (Cat 1) ===
\multicolumn{4}{|c|}{\cellcolor{blue!15}\textbf{Core Algorithms: Policy-Based (Category 1)}} \\
\hline
schulman2017ppo & Proximal Policy Optimization Algorithms & 1 & Policy-Based \\
\hline
schulman2015trust & Trust Region Policy Optimization & 1 & Policy-Based \\
\hline
mo2025matpo & MATPO: Multi-aspect Trust Region PO & 1 & Policy-Based \\
\hline
hong2025mgrpo & MGRPO: Multi-granularity RLHF & 1 & Policy-Based \\
\hline
li2025flowgrpo & FlowGRPO: Flow-based Group Relative PO & 1 & Policy-Based \\
\hline
ji2025treegrpo & Tree Search for LLM Agent RL & 1 & Policy-Based \\
\hline
li2024remax & ReMax: Simple RL for Aligning LLMs & 1 & Policy-Based \\
\hline
li2025stppo & ST-PPO: Stabilized Off-policy PPO for Agents & 1 & Policy-Based \\
\hline

% === LEARNING PARADIGMS: RLHF & Alignment (Cat 4) ===
\multicolumn{4}{|c|}{\cellcolor{green!15}\textbf{Learning Paradigms: RLHF \& Alignment (Category 4)}} \\
\hline
rafailov2023direct & DPO: Preference-based direct alignment & 4 & RLHF \& Alignment \\
\hline
deepseek2025r1 & DeepSeek-R1: Incentivizing Reasoning via RL & 4 & RLHF \& Alignment \\
\hline
christiano2017deep & Deep RL from Human Preferences & 4 & RLHF \& Alignment \\
\hline
stiennon2020learning & Learning to Summarize from Human Feedback & 4 & RLHF \& Alignment \\
\hline
bai2022\-constitutional & Constitutional AI: Harmlessness from AI & 4 & RLHF \& Alignment \\
\hline
lee2023rlaif & RLAIF: Scaling RLHF with AI Feedback & 4 & RLHF \& Alignment \\
\hline
ethayarajh2024kto & KTO: Prospect Theoretic Optimization & 4 & RLHF \& Alignment \\
\hline
bytedance2025dapo & DAPO: Decoupled Clip and Dynamic Sampling PO & 4 & RLHF \& Alignment \\
\hline

% === ADVANCED METHODS: Multi-Agent RL (Cat 7) ===
\multicolumn{4}{|c|}{\cellcolor{orange!15}\textbf{Advanced Methods: Multi-Agent RL (Category 7)}} \\
\hline
zhang2024agentnet & AgentNet++: Decentralized Multi-Agent & 7 & Multi-Agent RL \\
\hline
chen2024masrouter & MasRouter: Learning to Route Multi-Agent & 7 & Multi-Agent RL \\
\hline
chen2025agentflow & AgentFlow: Flow-based Multi-Agent & 7 & Multi-Agent RL \\
\hline
hong2023metagpt & MetaGPT: Multi-Agent Collaborative & 7 & Multi-Agent RL \\
\hline
qian2023chatdev & ChatDev: Communicative Agents for SW & 7 & Multi-Agent RL \\
\hline
park2023generative & Generative Agents: Interactive Simulacra & 7 & Multi-Agent RL \\
\hline
li2023camel & CAMEL: Communicative Agents for Mind & 7 & Multi-Agent RL \\
\hline
wu2023autogen & AutoGen: Multi-Agent Conversation & 7 & Multi-Agent RL \\
\hline

% === ADVANCED METHODS: Meta-Learning (Cat 9) ===
\multicolumn{4}{|c|}{\cellcolor{orange!15}\textbf{Advanced Methods: Meta-Learning (Category 9)}} \\
\hline
finn2017\-model\-agnostic & MAML: Model-Agnostic Meta-Learning & 9 & Meta-Learning \\
\hline

% === ADVANCED METHODS: Exploration (Cat 11) ===
\multicolumn{4}{|c|}{\cellcolor{orange!15}\textbf{Advanced Methods: Exploration (Category 11)}} \\
\hline
burda2018\-exploration & RND: Random Network Distillation & 11 & Exploration \\
\hline

% === APPLICATIONS: Agentic RL (Cat 12) ===
\multicolumn{4}{|c|}{\cellcolor{purple!15}\textbf{Applications: Agentic RL (Category 12)}} \\
\hline
yao2023react & ReAct: Synergizing Reasoning and Acting & 12 & Agentic RL \\
\hline
wei2022chain & Chain-of-Thought Prompting & 12 & Agentic RL \\
\hline
schick2023\-toolformer & Toolformer: LLMs Teach Tool Use & 12 & Agentic RL \\
\hline
chen2023fireact & FireAct: Language Agent Fine-tuning & 12 & Agentic RL \\
\hline
feng2025retool & ReTool: RL for Tool Discovery & 12 & Agentic RL \\
\hline
jiang2025verltool & VerlTool: Verifiable RL for Tool Agents & 12 & Agentic RL \\
\hline
le2025toolbrain & ToolBrain: Neural Tool Composition & 12 & Agentic RL \\
\hline
li2025deepagent & DeepAgent: Deep RL for Autonomous Dev & 12 & Agentic RL \\
\hline
zhang2025infoagent & InfoAgent: Information-theoretic Agent & 12 & Agentic RL \\
\hline
qin2023toolllm & ToolLLM: LLMs Master 16000+ APIs & 12 & Agentic RL \\
\hline
patil2023gorilla & Gorilla: LLM Connected with APIs & 12 & Agentic RL \\
\hline
liang2023taskmatrix & TaskMatrix.AI: Completing Tasks with APIs & 12 & Agentic RL \\
\hline
wang2023voyager & Voyager: Open-ended Embodied Agent & 12 & Agentic RL \\
\hline
shen2023hugginggpt & HuggingGPT: Solving AI Tasks & 12 & Agentic RL \\
\hline
wang2025\-mobile\-agent\-v3 & MobileAgent v3: Mobile Automation & 12 & Agentic RL \\
\hline
hong2023cogagent & CogAgent: Visual LM for GUI Agents & 12 & Agentic RL \\
\hline
yang2023appagent & AppAgent: Multimodal Smartphone Agents & 12 & Agentic RL \\
\hline
cheng2024seeclick & SeeClick: GUI Grounding for Visual Agents & 12 & Agentic RL \\
\hline
packer2023memgpt & MemGPT: LLMs as Operating Systems & 12 & Agentic RL \\
\hline

% === APPLICATIONS: Reasoning (Cat 14) ===
\multicolumn{4}{|c|}{\cellcolor{purple!15}\textbf{Applications: Reasoning (Category 14)}} \\
\hline
lightman2023lets & Verifying reasoning step-by-step & 14 & Reasoning \\
\hline
snell2024scaling & Scaling LLM Test-time Compute & 14 & Reasoning \\
\hline
hoffmann2022training & Chinchilla: Compute-optimal LLMs & 14 & Reasoning \\
\hline

% === APPLICATIONS: Security & Safety (Cat 15) ===
\multicolumn{4}{|c|}{\cellcolor{purple!15}\textbf{Applications: Security \& Safety (Category 15)}} \\
\hline
mazeika2024\-harm\-bench & HarmBench: Automated Red Teaming & 15 & Security \\
\hline
yin2024os & OS-Harm: Safety for Computer Use Agents & 15 & Security \\
\hline
sun2024trustllm & TrustLLM: Trustworthiness in LLMs & 15 & Security \\
\hline

% === APPLICATIONS: Benchmarks (Cat 16) ===
\multicolumn{4}{|c|}{\cellcolor{purple!15}\textbf{Applications: Benchmarks (Category 16)}} \\
\hline
jimenez2024\-swe\-bench & SWE-bench: Real-world GitHub Issues & 16 & Benchmarks \\
\hline
zhou2024webarena & WebArena: Realistic Web Environment & 16 & Benchmarks \\
\hline
xie2024osworld & OSWorld: Benchmarking Multimodal Agents & 16 & Benchmarks \\
\hline
liu2023agentbench & AgentBench: Benchmarking LLM Autonomy & 16 & Benchmarks \\
\hline
mialon2023gaia & GAIA: Benchmark for General AI Assistants & 16 & Benchmarks \\
\hline
liu2024\-visual\-agent\-bench & VisualAgentBench: Large Multimodal Models & 16 & Benchmarks \\
\hline

% === INFRASTRUCTURE: Surveys (Cat 18) ===
\multicolumn{4}{|c|}{\cellcolor{cyan!15}\textbf{Infrastructure: Surveys (Category 18)}} \\
\hline
wang2025landscape & Landscape of Agentic RL for LLMs & 18 & Surveys \\
\hline
xi2023rise & Rise of LLM-based Agents: A Survey & 18 & Surveys \\
\hline
sumers2023\-cognitive & Cognitive Architectures for Language Agents & 18 & Surveys \\
\hline
gao2023retrieval & RAG for LLMs: A Survey & 18 & Surveys \\
\hline
wang2024planning & Planning with LLMs: A Survey & 18 & Surveys \\
\hline
gutierrez2024hipporag & HippoRAG: Neurobiologically Inspired Memory & 18 & Surveys \\
\hline
shi2023replug & REPLUG: Retrieval-augmented Black-box LMs & 18 & Surveys \\
\hline
izacard2022atlas & Atlas: Few-shot Learning with Retrieval & 18 & Surveys \\
\hline
wang2023\-self\-knowledge & Self-knowledge Guided Retrieval & 18 & Surveys \\
\hline
sutton2018\-reinforce\-ment & RL: An Introduction (Textbook) & 18 & Surveys \\
\hline

\end{longtable}
}

\section*{Category Distribution Summary}

\begin{table}[htbp]
\centering
\begin{tabular}{|l|c|c|}
\hline
\textbf{Category Cluster} & \textbf{Categories} & \textbf{Paper Count} \\
\hline
Core Algorithms & 1--3 & 8 \\
Learning Paradigms & 4--6 & 0 \\
Advanced Methods & 7--11 & 10 \\
Applications & 12--16 & 31 \\
Infrastructure & 17--19 & 10 \\
\hline
\textbf{Total (Unique)} & \textbf{1--19} & \textbf{59} \\
\hline
\end{tabular}
\caption{Distribution of the 59 papers in Table~\ref{tab:paper_mapping} across taxonomy clusters, recounted from that table (each paper has one primary category). Categories 2--6, 8, 10, 13, 17, and 19 have no mapped paper; those categories are discussed from the cited literature in the main text but are not represented in this mapping.}
\label{tab:category_distribution}
\end{table}

% ========== REFERENCES ==========
% Note: Using \bibname to avoid duplicate headings
\renewcommand{\bibname}{References}
\bibliographystyle{plainnat}
\cleardoublepage\phantomsection
\addcontentsline{toc}{chapter}{References}
\bibliography{refs}

\end{document}
